%% ============================================================================
%%  CAPITOLO 2 — Complementi ai numeri reali
%%  Fonte principale per argomenti e dimostrazioni: Marcellini–Sbordone, cap. 2
%%  (§12 estremo superiore, §13 calcolo combinatorio, §14 binomio di Newton,
%%  §15 numeri complessi). Il capitolo è suddiviso in file nella cartella cap02/.
%% ============================================================================
\aperturacapitolo{Complementi ai numeri reali}{cap:complementi}

\epigraph{I numeri immaginari sono un meraviglioso rifugio dello spirito divino, quasi
un anfibio tra l'essere e il non essere.}{Gottfried Wilhelm Leibniz (1702)}

\lettrine[lines=3, findent=2pt, nindent=0pt]{N}{el capitolo} precedente abbiamo fissato
le regole del gioco: gli assiomi dei numeri reali. Tra essi l'assioma di completezza è
apparso come il più misterioso, e insieme come il più importante: è l'unico che
distingue $\R$ da $\Q$. In questo capitolo impariamo a usarlo davvero. Lo trasformeremo
in uno strumento di lavoro, il concetto di \emph{estremo superiore}, che è la chiave di
quasi tutti i teoremi di esistenza dell'Analisi: l'esistenza del limite delle
successioni monotòne, il teorema di Bolzano--Weierstrass, il teorema di Weierstrass sui
massimi e minimi delle funzioni continue, e così via (si riveda lo schema della
\cref{fig:schema-teoremi}).

Con l'estremo superiore manterremo anche tre promesse fatte nel \cref{cap:numeri}:
dimostreremo la \emph{proprietà di Archimede} (i numeri naturali non sono limitati),
l'esistenza della \emph{radice $n$-esima} di ogni numero reale positivo, e daremo
finalmente senso alla potenza $a^{x}$ con esponente $x$ reale qualsiasi, provando che
la funzione esponenziale è una corrispondenza biunivoca tra $\R$ e $(0,+\infty)$.

Nella seconda parte del capitolo cambiamo registro: il \emph{calcolo combinatorio}
insegna a contare senza elencare, e conduce alla \emph{formula del binomio di Newton}.
In appendice, infine, ampliamo $\R$ al campo $\C$ dei \emph{numeri complessi}, nel quale
ogni equazione di secondo grado ha soluzione e la moltiplicazione assume un significato
geometrico sorprendente: è una rotazione.

\medskip
\noindent\textbf{In questo capitolo.}
\begin{itemize}
  \item Massimo, minimo, maggioranti, estremo superiore e inferiore; il teorema di
  esistenza dell'estremo superiore (\S\ref{sec:sup}).
  \item Le proprietà di $\N$, la proprietà di Archimede, la parte intera, la densità di
  $\Q$ e degli irrazionali (\S\ref{sec:archimede}).
  \item Approfondimento: i numeri reali in forma decimale (\S\ref{subsec:forma-decimale}).
  \item L'esistenza della radice $n$-esima e la funzione esponenziale a esponente reale
  (\S\ref{sec:esponenziale-reale}).
  \item Disposizioni, permutazioni, combinazioni (\S\ref{sec:combinatoria}).
  \item La formula del binomio di Newton e il triangolo di Tartaglia (\S\ref{sec:binomio}).
  \item I numeri complessi: forma algebrica, trigonometrica, esponenziale; radici
  $n$-esime (\S\ref{sec:complessi}).
  \item Punti fondamentali del capitolo (\S\ref{sec:sintesi-cap02}).
\end{itemize}

%% ----------------------------------------------------------------------------
%%  §2.1  Massimo, minimo, estremo superiore, estremo inferiore
%% ----------------------------------------------------------------------------
\section{Massimo, minimo, estremo superiore, estremo inferiore}
\label{sec:sup}

\subsection{Massimo e minimo}

\begin{definizione}[Massimo e minimo]\index{massimo}\index{minimo}
Sia $A$ un insieme di numeri reali. Il \emph{massimo} di $A$, se esiste, è un numero
$M$ dell'insieme $A$ che è maggiore o uguale a ogni altro elemento dell'insieme; il
\emph{minimo} di $A$, se esiste, è un numero $m$ dell'insieme $A$ che è minore o uguale a
ogni altro elemento di $A$. In simboli:
\begin{align}
  M=\max A \quad&\sse\quad \begin{cases} M\ge a \quad \perogni a\in A,\\ M\in A;\end{cases}
  \label{eq:max-def}\\[2pt]
  m=\min A \quad&\sse\quad \begin{cases} m\le a \quad \perogni a\in A,\\ m\in A.\end{cases}
  \label{eq:min-def}
\end{align}
\end{definizione}

Si osservi che la definizione ha \emph{due} richieste, e che la seconda --- $M\in A$ ---
è quella che più spesso si dimentica. Non tutti gli insiemi di numeri reali hanno il
massimo e il minimo. Ad esempio, se $A$ è costituito da tutti i numeri reali positivi,
$A$ non ha né massimo né minimo: non esiste il più piccolo numero reale positivo (se $a>0$
allora anche $a/2$ è positivo, e $a/2<a$), e lo zero non è il minimo, perché non
appartiene ad $A$; analogamente, per ogni $a>0$ il numero $a+1$ è un elemento di $A$
maggiore di $a$.

\begin{proposizione}[Unicità del massimo e del minimo]\label{prop:unicita-max}
Se un insieme $A\subseteq\R$ ha massimo, il massimo è unico; lo stesso vale per il minimo.
\end{proposizione}

\begin{schema}
\Ipotesi{$A\subseteq\R$; \ $M_{1}$ e $M_{2}$ verificano entrambi la \eqref{eq:max-def}.}
\Tesi{$M_{1}=M_{2}$ (e analogamente per il minimo).}
\end{schema}

\begin{strategia}
Lo schema è quello di tutte le dimostrazioni di unicità (\cref{prop:unicita}): si
suppone che esistano due massimi $M_{1}$ e $M_{2}$ e si prova che coincidono. L'idea è
di usare le due richieste della definizione «incrociate»: $M_{1}$ è un elemento di $A$,
quindi $M_{2}$, che è maggiore o uguale di \emph{tutti} gli elementi di $A$, è maggiore
o uguale di $M_{1}$; e viceversa. L'antisimmetria dell'ordine conclude.
\end{strategia}

\begin{dimostrazione}
Se $M_{1}$ e $M_{2}$ sono due massimi di un insieme $A$, allora per definizione
\[
  M_{1}\ge a, \qquad M_{2}\ge a, \qquad \perogni a\in A;
\]
ma dato che $M_{1}$ ed $M_{2}$ sono elementi di $A$, posto $a$ rispettivamente uguale a
$M_{2}$ e ad $M_{1}$ nelle relazioni precedenti, si ottiene $M_{1}\ge M_{2}$ e
$M_{2}\ge M_{1}$, cioè $M_{1}=M_{2}$ per la proprietà antisimmetrica
\ref{ass:antisimmetrica}. Per il minimo si procede allo stesso modo.
\end{dimostrazione}

\subsection{Maggioranti, minoranti, insiemi limitati}

\begin{definizione}[Maggioranti, minoranti, insiemi limitati]\index{maggiorante}\index{minorante}\index{insieme!limitato}
Sia $A\subseteq\R$. Un numero reale $L$ si dice un \emph{maggiorante} di $A$ se $L\ge a$
per ogni $a\in A$; analogamente un numero reale $\ell$ è un \emph{minorante} di $A$ se
$\ell\le a$ per ogni $a\in A$. Diciamo che $A$ è \emph{limitato superiormente} se ammette
almeno un maggiorante, \emph{limitato inferiormente} se ammette almeno un minorante, e
\emph{limitato} se è limitato sia superiormente sia inferiormente:
\begin{equation}\label{eq:limitato}
  A \text{ limitato} \quad\sse\quad \esiste \ell,L\in\R : \ \ell\le a\le L \quad \perogni a\in A.
\end{equation}
\end{definizione}

È bene notare esplicitamente che un insieme non sempre ammette maggioranti o minoranti.
Se $A$ è di nuovo l'insieme dei numeri reali positivi, $A$ non ammette alcun maggiorante,
mentre lo zero (e anche qualsiasi numero reale negativo) è un minorante di $A$. Si noti
anche che i maggioranti, quando esistono, non sono mai uno solo: se $L$ è un maggiorante,
lo è anche ogni numero $L'>L$. Un maggiorante di $A$ che appartiene ad $A$ è, per
definizione, il massimo di $A$.

\begin{idea}[Maggioranti e tolleranze]
Un'immagine concreta. In un banco prova si misura più volte il carico di rottura di un
lotto di rivetti; l'insieme $A$ dei valori misurati è limitato superiormente, e un
qualsiasi maggiorante $L$ è un valore che \emph{nessun} rivetto del lotto ha superato.
Di maggioranti ce ne sono infiniti --- un valore assurdamente alto come $10^{9}$ newton lo
è senz'altro, ma non dice nulla --- e l'informazione utile è il \emph{più piccolo} dei
maggioranti, la stima più stretta possibile. L'estremo superiore che stiamo per definire
è esattamente questo: il miglior maggiorante.
\end{idea}

Tenendo presente la definizione \eqref{eq:modulo} del valore assoluto, si riconosce
facilmente la seguente caratterizzazione degli insiemi limitati, che useremo
moltissime volte.

\begin{proposizione}\label{prop:limitato-modulo}
Un insieme $A\subseteq\R$ è limitato se e soltanto se esiste un numero $M>0$ tale che
\begin{equation}\label{eq:limitato-modulo}
  \abs{a}\le M \qquad \perogni a\in A.
\end{equation}
\end{proposizione}

\begin{schema}
\Ipotesi{$A\subseteq\R$.}
\Tesi{$\bigl(\esiste\ell,L : \ell\le a\le L\ \perogni a\in A\bigr) \ \sse\ \bigl(\esiste M>0 : \abs{a}\le M\ \perogni a\in A\bigr)$.}
\end{schema}

\begin{strategia}
Le due condizioni dicono la stessa cosa in due modi diversi: la prima chiude $A$ in un
intervallo $[\ell,L]$ qualsiasi, la seconda in un intervallo \emph{simmetrico} $[-M,M]$.
Da un intervallo simmetrico si passa a uno qualsiasi gratis ($\ell=-M$, $L=M$); per il
viceversa basta allargare $[\ell,L]$ fino a renderlo simmetrico, scegliendo $M$ uguale al
più grande tra $\abs{\ell}$ e $\abs{L}$ (e almeno uguale a $1$, per avere $M>0$ anche
nel caso degenere $\ell=L=0$). Lo strumento tecnico è la
\cref{prop:modulo-r}, che traduce $\abs{a}\le M$ nella doppia disuguaglianza
$-M\le a\le M$.
\end{strategia}

\begin{dimostrazione}
Se per ipotesi vale la \eqref{eq:limitato-modulo}, allora dalla proprietà
\eqref{eq:mod-le} relativa al valore assoluto si ottiene
\begin{equation}\label{eq:limitato-passo1}
  -M\le a\le M \qquad \perogni a\in A,
\end{equation}
e quindi vale la \eqref{eq:limitato} con $\ell=-M$ e $L=M$. Viceversa, se vale la
\eqref{eq:limitato}, allora vale anche la \eqref{eq:limitato-modulo} (o
equivalentemente la \eqref{eq:limitato-passo1}) con $M=\max\{\abs{\ell},\abs{L},1\}$,
perché
\begin{equation}\label{eq:limitato-passo2}
  -M\le-\abs{\ell}\le\ell\le a\le L\le\abs{L}\le M \qquad \perogni a\in A,
\end{equation}
dove abbiamo usato $-\abs{x}\le x\le\abs{x}$ (\cref{teo:triangolare},
\eqref{eq:triang-passo1}). Il valore $1$ è stato incluso nel massimo solo per garantire
$M>0$ anche nel caso $\ell=L=0$.
\end{dimostrazione}

\subsection{Il teorema di esistenza dell'estremo superiore}

Il risultato che segue, alla base della definizione di estremo superiore, è
conseguenza dell'assioma di completezza \ref{ass:completezza} ed è, di fatto, il modo in
cui quell'assioma verrà usato in tutto il resto del libro.

\begin{teorema}[Esistenza dell'estremo superiore]\label{teo:esistenza-sup}\index{estremo superiore!esistenza}
Sia $A$ un insieme non vuoto di numeri reali, limitato superiormente. Allora l'insieme
dei maggioranti di $A$ ha minimo.
\end{teorema}

\begin{schema}
\Ipotesi{$A\subseteq\R$, \ $A\neq\emptyset$, \ $A$ limitato superiormente.}
\Tesi{Detto $B=\{L\in\R : L\ge a\ \perogni a\in A\}$ l'insieme dei maggioranti di $A$,
esiste $\min B$.}
\end{schema}

\begin{strategia}
L'assioma di completezza parla di \emph{due} insiemi separati, $A$ «a sinistra» e $B$ «a
destra». Abbiamo $A$; ci serve un secondo insieme che stia tutto alla sua destra, e il
candidato naturale è proprio l'insieme $B$ dei maggioranti: per definizione ogni $b\in B$
è $\ge$ di ogni $a\in A$. L'assioma fornisce allora un elemento separatore $M$. Resta da
leggere che cosa dice la doppia disuguaglianza $a\le M\le b$: la metà sinistra dice che
$M$ è un maggiorante, cioè $M\in B$; la metà destra che $M$ è minore o uguale di ogni
elemento di $B$. Le due cose insieme sono la definizione di minimo di $B$.
\end{strategia}

\begin{dimostrazione}
Indichiamo con $B$ l'insieme costituito dai maggioranti di $A$. $B$ è non vuoto, perché
$A$ è limitato superiormente; anche $A$ è non vuoto per ipotesi, e per definizione di
maggiorante $a\le b$ per ogni $a\in A$ e $b\in B$. Applichiamo l'assioma di completezza
\ref{ass:completezza} ai due insiemi $A$, $B$: esiste un numero reale $M$ tale che
\begin{equation}\label{eq:separatore-sup}
  a\le M\le b \qquad \perogni a\in A, \quad \perogni b\in B.
\end{equation}
Dato che $M$ è maggiore o uguale a tutti gli elementi di $A$, $M$ è un maggiorante di
$A$; cioè $M\in B$. Inoltre $M$ è minore o uguale a tutti gli elementi di $B$. Quindi, in
base alla definizione \eqref{eq:min-def}, $M$ è il minimo di $B$.
\end{dimostrazione}

\begin{figure}[htb]
\centering
\begin{tikzpicture}[scale=1.0, every node/.style={font=\footnotesize}]
  \draw[-Latex, grigio] (-0.5,0) -- (11.3,0);
  % l'insieme A
  \foreach \x in {0.4,0.9,1.6,2.2,2.9,3.4,3.9,4.3,4.6,4.8,4.9} \fill[blu] (\x,0) circle (1.7pt);
  \draw[blu, thick, decorate, decoration={brace, amplitude=4pt, mirror}] (0.3,-0.3) -- (5.0,-0.3)
    node[midway, below=5pt] {insieme $A$};
  % i maggioranti
  \fill[rame!25] (5.0,-0.12) rectangle (11.0,0.12);
  \draw[rame, thick] (5.0,-0.12) -- (5.0,0.12);
  \draw[rame, thick, decorate, decoration={brace, amplitude=4pt}] (5.0,0.35) -- (11.0,0.35)
    node[midway, above=5pt] {maggioranti di $A$: l'insieme $B$};
  \fill[verde] (5.0,0) circle (2.3pt);
  \draw[verde] (5.0,-0.12) -- (5.0,-0.55);
  \node[verde, anchor=north west] at (5.05,-0.3) {$M=\sup A=\min B$};
\end{tikzpicture}
\caption{L'estremo superiore è il «primo» dei maggioranti: l'insieme $B$ dei maggioranti
di $A$ è una semiretta chiusa $[M,+\infty)$, e il suo estremo sinistro $M$ è $\sup A$.
Il teorema di esistenza garantisce che la semiretta è \emph{chiusa}, cioè che il primo
maggiorante esiste.}
\label{fig:sup-maggioranti}
\end{figure}

In base al teorema precedente, poniamo la seguente definizione.

\begin{definizione}[Estremo superiore]\index{estremo superiore}
Sia $A$ un insieme di numeri reali non vuoto e limitato superiormente. Diciamo che
$M\in\R$ è l'\emph{estremo superiore} di $A$, e scriviamo $M=\sup A$, se $M$ è il minimo
dei maggioranti di $A$.
\end{definizione}

Per la \cref{prop:unicita-max} applicata all'insieme dei maggioranti, l'estremo
superiore è unico, e quindi è lecito parlare dell'estremo superiore, con l'articolo
determinativo. La
definizione, così com'è, è poco maneggevole: per verificare che un numero è il minimo
dei maggioranti dovremmo confrontarlo con \emph{tutti} i maggioranti. La proposizione
che segue la traduce in una forma operativa, quella che useremo nella pratica.

\begin{proposizione}[Caratterizzazione dell'estremo superiore]\label{prop:caratt-sup}
Sia $A\subseteq\R$ non vuoto e limitato superiormente, e sia $M\in\R$. Allora
\begin{equation}\label{eq:caratt-sup}
  M=\sup A \quad\sse\quad
  \begin{cases}
    \text{(i)}\ \ M\ge a \quad \perogni a\in A;\\[2pt]
    \text{(ii)}\ \perogni\epsilon>0\ \ \esiste a\in A : \ M-\epsilon<a.
  \end{cases}
\end{equation}
\end{proposizione}

\begin{schema}
\Ipotesi{$A\neq\emptyset$ limitato superiormente; $M\in\R$.}
\Tesi{$M$ è il minimo dei maggioranti di $A$ \emph{se e solo se} $M$ è un maggiorante
(i) e nessun numero più piccolo di $M$ lo è (ii).}
\end{schema}

\begin{strategia}
La condizione (i) è la prima metà della definizione di minimo dei maggioranti: dice che
$M$ è un maggiorante. La condizione (ii) è la seconda metà, letta «al contrario»: dire
che $M$ è il \emph{più piccolo} maggiorante equivale a dire che ogni numero più piccolo
di $M$, che possiamo sempre scrivere come $M-\epsilon$ con $\epsilon>0$, \emph{non} è un
maggiorante; e «$M-\epsilon$ non è un maggiorante» significa, negando la definizione,
che c'è un elemento di $A$ che lo supera. Tutta la dimostrazione è un esercizio di
\emph{negazione} di un'affermazione con il «per ogni».
\end{strategia}

\begin{dimostrazione}
($\implica$) Sia $M=\sup A$. Allora $M$ è un maggiorante, cioè vale (i). Sia $\epsilon>0$;
il numero $M-\epsilon$ è minore di $M$, e poiché $M$ è il \emph{minimo} dei maggioranti,
$M-\epsilon$ non è un maggiorante. Negare «$M-\epsilon\ge a$ per ogni $a\in A$» significa
affermare che esiste $a\in A$ con $M-\epsilon<a$: è la (ii).

($\Leftarrow$) Supponiamo che valgano (i) e (ii). Per (i), $M$ è un maggiorante. Sia
$L$ un qualsiasi maggiorante di $A$ e supponiamo per assurdo $L<M$. Posto
$\epsilon=M-L>0$, per (ii) esiste $a\in A$ con $a>M-\epsilon=L$, contro il fatto che $L$ è
un maggiorante. Dunque $L\ge M$ per ogni maggiorante $L$: $M$ è il minimo dei
maggioranti, cioè $M=\sup A$.
\end{dimostrazione}

\begin{figure}[htb]
\centering
\begin{tikzpicture}[scale=1.0, every node/.style={font=\footnotesize}]
  \draw[-Latex, grigio] (-0.5,0) -- (11.3,0);
  \foreach \x in {0.6,1.4,2.3,3.1,3.9,4.6,5.2,5.8,6.3,6.7,7.0,7.25,7.45,7.6,7.7} \fill[blu] (\x,0) circle (1.6pt);
  \fill[verde] (7.9,0) circle (2.2pt) node[above=4pt] {$M$};
  \draw[rosso, thick] (6.5,-0.15) -- (6.5,0.15) node[above=2pt] {$M-\epsilon$};
  \fill[rossochiaro, opacity=0.7] (6.5,-0.1) rectangle (7.9,0.1);
  \foreach \x in {6.7,7.0,7.25,7.45,7.6,7.7} \fill[blu] (\x,0) circle (1.6pt);
  \draw[rosso, decorate, decoration={brace, amplitude=3pt, mirror}] (6.5,-0.3) -- (7.9,-0.3)
    node[midway, below=4pt, text width=4.2cm, align=center] {qui cade sempre almeno\\ un elemento di $A$, per ogni $\epsilon>0$};
  \node[blu, below=4pt] at (2.5,-0.1) {$A$};
\end{tikzpicture}
\caption{La caratterizzazione \eqref{eq:caratt-sup}: per quanto piccolo si scelga
$\epsilon>0$, nell'intervallo $(M-\epsilon,M]$ c'è sempre qualche elemento di $A$. Non
è detto che $M$ stesso appartenga ad $A$.}
\label{fig:caratt-sup}
\end{figure}

Analogamente si verifica che, se $A$ è un insieme non vuoto di numeri reali limitato
inferiormente, allora l'insieme dei minoranti di $A$ ha massimo (basta ripetere la
dimostrazione del \cref{teo:esistenza-sup} scambiando i ruoli dei due insiemi
nell'assioma di completezza).

\begin{definizione}[Estremo inferiore]\index{estremo inferiore}
Sia $A\subseteq\R$ non vuoto e limitato inferiormente. Si dice che un numero $m$ è
l'\emph{estremo inferiore} di $A$, e si scrive $m=\inf A$, se $m$ è il massimo dei
minoranti di $A$. Come per l'estremo superiore, ciò equivale a:
\begin{equation}\label{eq:caratt-inf}
  m=\inf A \quad\sse\quad
  \begin{cases}
    m\le a \quad \perogni a\in A;\\[2pt]
    \perogni\epsilon>0\ \ \esiste a\in A : \ m+\epsilon>a.
  \end{cases}
\end{equation}
\end{definizione}

Quindi: \emph{se un insieme non vuoto è limitato superiormente, esiste l'estremo
superiore ed è un numero reale; se un insieme non vuoto è limitato inferiormente,
esiste l'estremo inferiore ed è un numero reale.}

\begin{osservazione}[Estremo superiore e massimo]\label{oss:sup-max}
Se $A$ ha massimo, allora $\max A=\sup A$: il massimo è un maggiorante che appartiene ad
$A$, e ogni altro maggiorante $L$ deve essere $\ge\max A$ proprio perché $\max A\in A$.
Viceversa, se $\sup A\in A$, allora $\sup A$ è il massimo di $A$. In sintesi:
\emph{l'estremo superiore esiste sempre} (per gli insiemi non vuoti e limitati
superiormente), \emph{il massimo esiste se e solo se l'estremo superiore appartiene
all'insieme}. L'estremo superiore è il «massimo che c'è sempre», il prezzo da pagare è
che non sta necessariamente nell'insieme.
\end{osservazione}

\begin{osservazione}[Completezza ed estremo superiore sono la stessa cosa]\label{oss:completezza-sup}
Abbiamo dedotto il \cref{teo:esistenza-sup} dall'assioma di completezza. Vale anche il
viceversa: se in un sistema di numeri che soddisfa gli assiomi delle operazioni e
dell'ordine ogni insieme non vuoto e limitato superiormente ha estremo superiore, allora
vale l'assioma di completezza. Infatti, siano $A$ e $B$ non vuoti con $a\le b$ per ogni
$a\in A$, $b\in B$. Ogni $b\in B$ è un maggiorante di $A$, quindi $A$ è limitato
superiormente ed esiste $c=\sup A$. Allora $a\le c$ per ogni $a\in A$ (perché $c$ è un
maggiorante) e $c\le b$ per ogni $b\in B$ (perché $c$ è il \emph{minimo} dei maggioranti
e $b$ è un maggiorante). Dunque $c$ separa $A$ e $B$. Per questo motivo molti testi assumono come assioma di
completezza direttamente l'esistenza dell'estremo superiore, mentre altri, come
\cite{tao1}, costruiscono $\R$ a partire da $\Q$ e dimostrano questa proprietà come
teorema: le diverse impostazioni sono equivalenti.
\end{osservazione}

\begin{osservazione}[Nel modello di Dedekind il sup è un'unione]\label{oss:modello-dedekind}
Chi costruisce $\R$ a partire da $\Q$ ottiene il \cref{teo:esistenza-sup} quasi senza
fatica. Nel modello delle sezioni di $\Q$ costruito nel \S\ref{subsec:modello-reali}
(\cref{def:sezione-Q}) una sezione $(A,B)$ è determinata dalla sua classe inferiore $A$, e
l'ordinamento è l'inclusione delle classi inferiori. Se $X$ è un insieme non vuoto di
sezioni, tutte precedute da una stessa sezione, l'unione $\Sigma$ delle loro classi
inferiori è ancora la classe inferiore di una sezione $\sigma$: è non vuota, non è tutto
$\Q$, è «chiusa verso sinistra» e non ha massimo, come nella dimostrazione del
\cref{teo:completezza-modello}. Ogni elemento di $X$ precede $\sigma$, e ogni maggiorante
di $X$, che contiene tutte quelle classi inferiori, contiene anche la loro unione:
$\sigma$ è il minimo dei maggioranti, cioè $\sigma=\sup X$. Il postulato diventa una
semplice proprietà degli insiemi; la fatica si sposta sulle operazioni e sugli assiomi di
campo, dove le sezioni sono molto meno maneggevoli (si veda anche \cite{giusti}).
\end{osservazione}

\begin{attenzione}[In $\Q$ l'estremo superiore può non esistere]
Il \cref{teo:esistenza-sup} è falso nel campo dei razionali. L'insieme
$A=\{a\in\Q : a^{2}<2\}$ è non vuoto e limitato superiormente in $\Q$ (ad esempio $2$ è un
maggiorante razionale), ma \emph{non ha un minimo maggiorante razionale}: il
\cref{teo:Q-non-completo} mostra precisamente che nessun razionale separa $A$ dai
razionali positivi con quadrato maggiore di $2$. In $\R$, invece, $\sup A$ esiste ed è
$\sqrt2$: l'esistenza di $\sqrt2$ segue dal \cref{teo:radice-n}, e il fatto che sia
proprio l'estremo superiore di $A$ dalla densità di $\Q$ (\cref{es:sup-razionali}).
\end{attenzione}

\begin{notastorica}[Bolzano e il «principio dell'estremo superiore»]
Il primo a enunciare e usare consapevolmente l'esistenza dell'estremo superiore fu il
sacerdote e matematico boemo Bernard Bolzano (1781--1848), nel saggio \emph{Rein
analytischer Beweis} (1817), dedicato alla dimostrazione «puramente analitica» del
teorema dei valori intermedi. Bolzano voleva liberare l'Analisi dal ricorso
all'intuizione geometrica («la curva deve pur attraversare l'asse, visto che passa da
sopra a sotto»), e capì che serviva una proprietà dei numeri che garantisse l'esistenza
di un «confine» per ogni insieme limitato. Il suo lavoro, pubblicato a Praga in tedesco,
rimase però quasi sconosciuto per decenni; l'estremo superiore entrò nell'uso comune solo
con le lezioni berlinesi di Karl Weierstrass, negli anni Sessanta dell'Ottocento.
\end{notastorica}

\subsection{I simboli \texorpdfstring{$+\infty$}{+∞} e \texorpdfstring{$-\infty$}{−∞}}

È utile introdurre i simboli $+\infty$ e $-\infty$ per descrivere gli insiemi
\emph{non} limitati. Precisamente, sia $A$ un insieme non vuoto. L'estremo superiore di
$A$ è $+\infty$ se $A$ non è limitato superiormente; l'estremo inferiore di $A$ è
$-\infty$ se $A$ non è limitato inferiormente. In simboli:
\begin{align}
  \sup A=+\infty \quad&\sse\quad \perogni L\in\R\ \ \esiste a\in A : \ a>L; \label{eq:sup-infinito}\\
  \inf A=-\infty \quad&\sse\quad \perogni \ell\in\R\ \ \esiste a\in A : \ a<\ell. \label{eq:inf-infinito}
\end{align}
Nelle relazioni sopra scritte ci si può limitare a considerare $L>0$ e $\ell<0$ (se un
numero positivo non è un maggiorante, non lo è nessun numero più piccolo). Si noti che
la \eqref{eq:sup-infinito} non è altro che la negazione di «$A$ è limitato
superiormente».

Facendo uso dei simboli $+\infty$ e $-\infty$ si può quindi affermare che \emph{ogni
insieme non vuoto di numeri reali ammette sia estremo superiore sia estremo inferiore};
se l'insieme è limitato superiormente l'estremo superiore è finito, e se è limitato
inferiormente l'estremo inferiore è finito.

\begin{attenzione}[$+\infty$ non è un numero reale]
I simboli $\pm\infty$ non sono elementi di $\R$, e con essi non si possono eseguire le
operazioni come con i numeri: scrivere $\sup A=+\infty$ è soltanto un modo abbreviato di
dire che $A$ non ha maggioranti. Nel seguito del corso, studiando i limiti, introdurremo
la \emph{retta reale estesa} $\overline{\R}=\R\cup\{-\infty,+\infty\}$ e daremo regole
precise (e limitate) per operare con questi simboli.
\end{attenzione}

\subsection{Esempi}

\begin{esempio}\label{es:sup-semiretta}
Se $A=\{x\in\R : x>0\}$, allora
\begin{equation}\label{eq:es-semiretta}
  \sup A=+\infty, \qquad \inf A=0,
\end{equation}
e il massimo e il minimo di $A$ non esistono. Infatti, per ogni $L>0$ il numero $L+1$
appartiene ad $A$ ed è maggiore di $L$: vale la \eqref{eq:sup-infinito}. Inoltre $0$ è
un minorante, e per ogni $\epsilon>0$ il numero $a=\epsilon/2\in A$ verifica
$a<0+\epsilon$: per la \eqref{eq:caratt-inf}, $\inf A=0$. Poiché $0\notin A$, il minimo
non esiste (\cref{oss:sup-max}).
\end{esempio}

\begin{esempio}\label{es:sup-B}
Sia $B=\set{\dfrac{n-1}{n} : n\in\N}=\set{0,\ \tfrac12,\ \tfrac23,\ \tfrac34,\ \dots}$
(l'insieme è schematizzato nella \cref{fig:insiemi-B-C}). Risulta
\begin{equation}\label{eq:es-B}
  \sup B=1, \qquad \inf B=\min B=0.
\end{equation}
Infatti $\frac{n-1}{n}=1-\frac1n$ è minore di $1$ per ogni $n$, e maggiore o uguale a
$0$, con uguaglianza per $n=1$: dunque $0=\min B=\inf B$ e $1$ è un maggiorante. Per
provare che $1$ è il \emph{minimo} maggiorante, fissato $\epsilon>0$ cerchiamo
$n\in\N$ con
\[
  1-\frac1n>1-\epsilon \quad\sse\quad \frac1n<\epsilon \quad\sse\quad n>\frac1\epsilon .
\]
Un tale $n$ esiste per la \emph{proprietà di Archimede}, che dimostreremo nel prossimo
paragrafo (\cref{teo:archimede}) senza usare questo esempio: nessun circolo vizioso.
Poiché $1\notin B$, l'insieme $B$ non ha massimo.
\end{esempio}

\begin{esempio}\label{es:sup-C}
Se infine $C=\set{\dfrac{n+1}{n} : n\in\N}=\set{2,\ \tfrac32,\ \tfrac43,\ \tfrac54,\ \dots}$,
si trova
\begin{equation}\label{eq:es-C}
  \sup C=\max C=2, \qquad \inf C=1.
\end{equation}
Il lettore ripeta il ragionamento dell'esempio precedente, osservando che
$\frac{n+1}{n}=1+\frac1n$. Qui la situazione è rovesciata: il massimo esiste, il minimo
no.
\end{esempio}

\begin{figure}[htb]
\centering
\begin{tikzpicture}[x=9cm, every node/.style={font=\footnotesize}]
  % insieme B
  \draw[-Latex, grigio] (-0.06,0) -- (1.12,0);
  \foreach \n in {1,...,80} {\fill[blu] ({(\n-1)/\n},0) circle (1.2pt);}
  \draw[rosso, thick, fill=white] (1,0) circle (2.2pt);
  \foreach \n/\lab in {1/0, 2/\tfrac12, 3/\tfrac23, 4/\tfrac34, 5/\tfrac45} {
    \node[below=2pt] at ({(\n-1)/\n},0) {$\lab$};}
  \node[above=3pt, blu] at (0,0) {$0=\min B$};
  \node[above=3pt, rosso] at (1,0) {$1=\sup B\notin B$};
  \node[left=6pt, blu] at (-0.06,0) {$B$};
  % insieme C
  \begin{scope}[yshift=-1.5cm, x=9cm]
  \draw[-Latex, grigio] (-0.06,0) -- (1.12,0);
  \foreach \n in {1,...,80} {\fill[rame] ({1/\n},0) circle (1.2pt);}
  \draw[rosso, thick, fill=white] (0,0) circle (2.2pt);
  \foreach \n/\lab in {1/2, 2/\tfrac32, 3/\tfrac43, 4/\tfrac54, 5/\tfrac65} {
    \node[below=2pt] at ({1/\n},0) {$\lab$};}
  \node[above=3pt, rosso] at (0,0) {$1=\inf C\notin C$};
  \node[above=3pt, rame] at (1,0) {$2=\max C$};
  \node[left=6pt, rame] at (-0.06,0) {$C$};
  \end{scope}
\end{tikzpicture}
\caption{Gli insiemi $B=\{(n-1)/n\}$ e $C=\{(n+1)/n\}$ degli \cref{es:sup-B,es:sup-C}
(la seconda retta è traslata: l'origine del disegno corrisponde a $1$). Gli elementi si
accumulano verso un estremo che non appartiene all'insieme (cerchietto rosso).}
\label{fig:insiemi-B-C}
\end{figure}

\begin{esempio}[Velocità limite: un estremo superiore che non è un massimo]\label{es:velocita-limite}
Un paracadutista in caduta libera, di peso $W$, è soggetto alla resistenza
aerodinamica $D=\tfrac12\rho V^{2}SC_{D}$. La velocità cresce, tendendo al valore per cui
la resistenza eguaglierebbe il peso; questa \emph{velocità limite} si ottiene da $W=\tfrac12\rho V_{t}^{2}SC_{D}$:
\[
  V_{t}=\sqrt{\frac{2W}{\rho\,S\,C_{D}}}\,.
\]
Con $W=800$~N, $\rho=1{,}225$~kg/m$^{3}$ e $SC_{D}=0{,}5$~m$^{2}$ (posizione «a
pancia in giù») si trova $V_{t}\approx51{,}1$~m/s, cioè circa $184$~km/h. Risolvendo l'equazione del moto (lo faremo nel capitolo sulle equazioni
differenziali) si trova che, partendo da fermo, la velocità all'istante $t$ è
\[
  V(t)=V_{t}\,\frac{1-e^{-2gt/V_{t}}}{1+e^{-2gt/V_{t}}}\,.
\]
L'insieme $E=\{V(t): t\ge0\}$ delle velocità effettivamente raggiunte ha
$\inf E=\min E=0$ (all'istante iniziale) e $\sup E=V_{t}$; ma
$V_{t}\notin E$, perché la frazione è sempre strettamente minore di $1$. La
velocità limite è un estremo superiore che \emph{non} è un massimo: il paracadutista vi
si avvicina quanto si vuole (dopo una quindicina di secondi la differenza è inferiore
all'$1\%$), ma non la raggiunge mai (\cref{fig:velocita-limite}). In termini della
\cref{prop:caratt-sup}: per ogni tolleranza $\epsilon>0$ esiste un istante $t$ in cui
$V(t)>V_{t}-\epsilon$.
\end{esempio}

\begin{figure}[htb]
\centering
\begin{tikzpicture}
\begin{axis}[width=0.92\textwidth, height=5.2cm, xmin=-0.5, xmax=22, ymin=-4, ymax=62,
  xlabel={$t$ [s]}, ylabel={$V$ [m/s]}, xtick={5,10,15,20}, ytick={10,20,30,40,51.1},
  yticklabels={10,20,30,40,$V_{t}$}]
  \addplot[tratteggio, domain=0:22] {51.1};
  \addplot[curva, domain=0:21.5, samples=150] {51.1*tanh(9.81*x/51.1)};
  \addplot[only marks, mark=*, mark size=1.2pt, rame] coordinates {(0,0)};
  \node[font=\scriptsize, rame, anchor=south west] at (axis cs:1.0,0.6) {$\min E=V(0)=0$};
  \node[font=\scriptsize, grigio, anchor=south east] at (axis cs:21.5,51.5) {$\sup E=V_{t}$ (mai raggiunto)};
\end{axis}
\end{tikzpicture}
\caption{La velocità di un paracadutista in caduta libera: cresce sempre, si avvicina
quanto si vuole alla velocità limite $V_{t}$ (tratteggiata), ma non la raggiunge mai.}
\label{fig:velocita-limite}
\end{figure}

\subsection{Estremi di una funzione}

\begin{definizione}[Estremo superiore e inferiore di una funzione]\index{funzione!limitata}
Siano $A$ e $B$ due sottoinsiemi di $\R$ e $f$ una funzione da $A$ verso $B$. Per ogni
sottoinsieme $X$ di $A$, l'estremo inferiore (risp.\ superiore) dell'insieme immagine
$f(X)$ si chiama \emph{estremo inferiore} (risp.\ \emph{superiore}) \emph{di $f$ su
$X$}, e si scrive
\begin{equation}\label{eq:sup-funzione}
  \inf_{x\in X}f(x)=\inf f(X), \qquad \sup_{x\in X}f(x)=\sup f(X).
\end{equation}
Se $f(X)$ è limitato inferiormente (risp.\ superiormente) si dice che la funzione $f$ è
\emph{limitata inferiormente} (risp.\ \emph{superiormente}) su $X$; se $f(X)$ è
limitato, si dice che $f$ è \emph{limitata} su $X$. Se $\sup f(X)$ appartiene a $f(X)$,
cioè se esiste $x_{0}\in X$ con $f(x_{0})=\sup_{X}f$, si dice che $f$ \emph{ha massimo}
su $X$ e che $x_{0}$ è un \emph{punto di massimo}.
\end{definizione}

\begin{esempio}\label{es:sup-funzione}
Sia $f(x)=\dfrac{x}{1+x}$ su $X=[0,+\infty)$. Si ha $f(0)=0\le f(x)$, quindi
$\min_{X}f=0$. Inoltre $f(x)=1-\frac{1}{1+x}<1$ per ogni $x\ge0$, dunque $1$ è un
maggiorante; e per ogni $\epsilon\in(0,1)$, scelto $x>\frac{1-\epsilon}{\epsilon}$, si ha
$1+x>\frac1\epsilon$, cioè $\frac{1}{1+x}<\epsilon$, cioè $f(x)>1-\epsilon$ (per
$\epsilon\ge1$ la richiesta è ancora più debole: basta $x=1$, perché
$f(1)=\frac12>0\ge1-\epsilon$). Per la
\cref{prop:caratt-sup}, $\sup_{X}f=1$, ma $f$ non ha massimo: è la stessa situazione
della velocità limite. La funzione $g(x)=1/x$ su $(0,1]$, invece, non è limitata
superiormente: $\sup_{(0,1]}g=+\infty$ (per ogni $L>1$ si ha $g(1/(2L))=2L>L$).
\end{esempio}

\begin{marginfigure}
\centering
\begin{tikzpicture}
\begin{axis}[width=5.3cm, height=4.2cm, xmin=-0.5, xmax=10, ymin=-0.15, ymax=1.3,
  xtick={5}, ytick={1}]
  \addplot[tratteggio, domain=0:10] {1};
  \addplot[curva, domain=0:9.8, samples=100] {x/(1+x)};
\end{axis}
\end{tikzpicture}
\caption{$f(x)=x/(1+x)$ su $[0,+\infty)$: il minimo vale $0$, l'estremo superiore vale
$1$ e non è un massimo.}
\end{marginfigure}

\begin{python}[Calcolare un estremo superiore per bisezione]
# Se A è chiuso verso sinistra (x in A e y < x implicano y in A),
# come A = {x : x^3 < 5}, l'estremo superiore si approssima per
# bisezione: si tengono un punto a in A e un maggiorante b, e si
# dimezza ogni volta l'intervallo [a, b].
def sup_bisezione(appartiene, a, b, tol=1e-12):
    """a deve stare in A, b deve essere un maggiorante di A."""
    while b - a > tol:
        m = (a + b) / 2
        if appartiene(m):
            a = m          # m sta in A: il sup è almeno m
        else:
            b = m          # m non sta in A: è un maggiorante
    return a, b

# A = {x : x^3 < 5}: 0 sta in A, 5 è un maggiorante, e sup A è
# la radice cubica di 5
a, b = sup_bisezione(lambda x: x**3 < 5, 0.0, 5.0)
print(f"{a:.14f} <= sup A <= {b:.14f}")
print(f"5 ** (1/3) = {5 ** (1/3):.14f}")
\end{python}
\begin{risultato}
1.70997594667654 <= sup A <= 1.70997594667710
5 ** (1/3) = 1.70997594667670
\end{risultato}

Il programma funziona perché l'insieme $A$ ha una struttura particolare (è un intervallo
che si estende verso sinistra): per un insieme qualsiasi, stabilire se un numero è un
maggiorante richiederebbe di confrontarlo con \emph{tutti} gli elementi, cosa che un
calcolatore non può fare in un numero finito di passi. È un primo esempio di come la
nozione di estremo superiore, pur essendo «esistenziale», possa diventare costruttiva in
casi concreti. La dimostrazione del \cref{teo:radice-n}, più avanti, si basa sulla stessa
intuizione --- «stringere» il valore cercato tra gli elementi dell'insieme e i suoi
maggioranti --- ma invece di procedere per bisezione ragiona direttamente sull'estremo
superiore.

%% ----------------------------------------------------------------------------
%%  §2.2  Proprietà di N, proprietà di Archimede, densità di Q
%% ----------------------------------------------------------------------------
\section{I numeri naturali, la proprietà di Archimede, la densità di \texorpdfstring{$\Q$}{Q}}
\label{sec:archimede}

Alla luce delle nozioni introdotte nel paragrafo precedente riportiamo di seguito alcune
proprietà dell'insieme dei numeri naturali e dell'insieme dei numeri razionali. Nell'insieme
$\R$, a partire dall'elemento $1$, si possono determinare gli elementi $2=1+1$, $3=2+1$ e
così via; tali elementi costituiscono l'insieme $\N=\{1,2,3,\dots\}$ dei numeri naturali
(\cref{sec:nzq}), i cui elementi sono ordinati secondo le relazioni $1<2<3<\cdots$. Il
principio di induzione (\cref{teo:induzione}) precisa il significato dei puntini: $\N$ è
il più piccolo sottoinsieme di $\R$ che contiene $1$ e che, insieme a ogni suo elemento
$n$, contiene anche $n+1$.

\subsection{Due proprietà caratteristiche di \texorpdfstring{$\N$}{N}}

Cominciamo con un lemma tecnico che rende rigorosa un'evidenza: tra due naturali
consecutivi non ci sono altri naturali.

\begin{lemma}[Struttura di $\N$]\label{lem:struttura-N}
\leavevmode
\begin{enumerate}
  \item Per ogni $n\in\N$ si ha $n\ge1$.
  \item Se $n\in\N$ e $n>1$, allora $n-1\in\N$.
  \item Se $n,m\in\N$ e $m>n$, allora $m\ge n+1$.
\end{enumerate}
\end{lemma}

\begin{schema}
\Ipotesi{$n,m\in\N$; \ il principio di induzione (\cref{teo:induzione}).}
\Tesi{(i) $n\ge1$; \ (ii) $n>1\implica n-1\in\N$; \ (iii) $m>n\implica m\ge n+1$.}
\end{schema}

\begin{strategia}
Tutte e tre le proprietà riguardano «tutti i numeri naturali», e l'unico strumento che
abbiamo per dimostrare qualcosa per tutti i naturali è il principio di induzione. Per
(i) e (ii) l'induzione è immediata. Per (iii) si fa induzione su $n$, tenendo $m$
generico: il caso $n=1$ si riduce a (i) e (ii); nel passo induttivo, se $m>n+1$ si
«toglie $1$» a $m$ (lecito per (ii)) e ci si riconduce all'ipotesi induttiva.
\end{strategia}

\begin{dimostrazione}
(i) La proprietà vale per $n=1$; se $n\ge1$ allora $n+1\ge1+1>1$ (si ricordi che $0<1$,
\cref{es:zero-minore-uno}). Per induzione vale per ogni $n$.

(ii) Sia $\mathcal{P}(n)$ la proposizione «$n=1$ oppure $n-1\in\N$». $\mathcal{P}(1)$ è
vera. Se vale $\mathcal{P}(n)$, allora $\mathcal{P}(n+1)$ è vera perché
$(n+1)-1=n\in\N$. Dunque $\mathcal{P}(n)$ vale per ogni $n$; se $n>1$, allora $n\neq1$, e
quindi $n-1\in\N$.

(iii) Sia $\mathcal{P}(n)$ la proposizione «per ogni $m\in\N$, se $m>n$ allora
$m\ge n+1$». Per $n=1$: se $m>1$, per (ii) $m-1\in\N$ e per (i) $m-1\ge1$, cioè $m\ge2$.
Supponiamo vera $\mathcal{P}(n)$ e sia $m\in\N$ con $m>n+1$. Allora $m>1$, quindi
$m-1\in\N$ e $m-1>n$; per l'ipotesi induttiva $m-1\ge n+1$, cioè $m\ge n+2$. Questo è
$\mathcal{P}(n+1)$.
\end{dimostrazione}

Con ragionamenti analoghi (per induzione) si prova che la somma e il prodotto di numeri
naturali sono numeri naturali; lo diamo per acquisito. L'insieme $\N$ gode di due
proprietà caratteristiche, che enunciamo e dimostriamo.

\begin{teorema}[Principio del buon ordinamento]\label{teo:buon-ordinamento}\index{principio del buon ordinamento}
Ogni sottoinsieme non vuoto di $\N$ ha minimo.
\end{teorema}

\begin{schema}
\Ipotesi{$A\subseteq\N$, \ $A\neq\emptyset$.}
\Tesi{Esiste $\min A$.}
\end{schema}

\begin{strategia}
Ragioniamo per assurdo: supponiamo che $A$ non abbia minimo e dimostriamo che allora $A$ è
vuoto. L'idea è di «risalire» i naturali uno alla volta con l'induzione, mostrando che
nessuno di essi può stare in $A$: se $1,2,\dots,n$ non stanno in $A$ e ci stesse $n+1$,
allora $n+1$ sarebbe il più piccolo elemento di $A$ (non ci sono naturali tra $n$ e
$n+1$, \cref{lem:struttura-N}(iii)) --- ma avevamo supposto che il minimo non ci fosse.
\end{strategia}

\begin{dimostrazione}
Supponiamo per assurdo che $A$ non abbia minimo, e sia $\mathcal{P}(n)$ la proposizione
«nessun naturale $k\le n$ appartiene ad $A$».

\emph{Base.} Se fosse $1\in A$, poiché $1\le a$ per ogni $a\in A$ (\cref{lem:struttura-N}(i)),
$1$ sarebbe il minimo di $A$: escluso. Dunque $1\notin A$; e poiché l'unico naturale
$k\le1$ è $k=1$ (per il \cref{lem:struttura-N}(i) e la proprietà antisimmetrica),
$\mathcal{P}(1)$ è vera.

\emph{Passo induttivo.} Supponiamo vera $\mathcal{P}(n)$. Se fosse $n+1\in A$, allora
ogni $a\in A$ verificherebbe $a>n$ (per $\mathcal{P}(n)$), quindi $a\ge n+1$ per il
\cref{lem:struttura-N}(iii): $n+1$ sarebbe il minimo di $A$, escluso. Quindi
$n+1\notin A$; e poiché i naturali $k\le n+1$ sono quelli $\le n$ e $n+1$ stesso
(ancora per il lemma), vale $\mathcal{P}(n+1)$.

Per il principio di induzione $\mathcal{P}(n)$ vale per ogni $n$; in particolare ogni
$n\in\N$ non appartiene ad $A$ (si applichi $\mathcal{P}(n)$ con $k=n$). Dunque
$A=\emptyset$, contro l'ipotesi.
\end{dimostrazione}

\begin{teorema}\label{teo:max-N}
Ogni sottoinsieme non vuoto di $\N$, limitato superiormente, ha massimo.
\end{teorema}

\begin{schema}
\Ipotesi{$A\subseteq\N$, \ $A\neq\emptyset$, \ $A$ limitato superiormente.}
\Tesi{Esiste $\max A$.}
\end{schema}

\begin{strategia}
Per il \cref{teo:esistenza-sup}, $A$ ha un estremo superiore $M$; dobbiamo provare che
$M\in A$. Per la caratterizzazione \eqref{eq:caratt-sup} con $\epsilon=1$, c'è un
elemento $n\in A$ «vicino» a $M$, precisamente $n>M-1$. Ora sfruttiamo il fatto che gli
elementi di $A$ sono naturali, dunque «distano almeno $1$» l'uno dall'altro: un elemento
di $A$ più grande di $n$ dovrebbe essere $\ge n+1>M$, impossibile. Quindi $n$ è il
massimo.
\end{strategia}

\begin{dimostrazione}
Sia $M=\sup A$, che esiste per il \cref{teo:esistenza-sup}. Per la
\cref{prop:caratt-sup} con $\epsilon=1$, esiste $n\in A$ con $n>M-1$. Se esistesse
$m\in A$ con $m>n$, per il \cref{lem:struttura-N}(iii) sarebbe $m\ge n+1>M$, contro il
fatto che $M$ è un maggiorante di $A$. Dunque $m\le n$ per ogni $m\in A$, e $n\in A$:
$n=\max A$.
\end{dimostrazione}

\begin{osservazione}
I due teoremi si estendono ai sottoinsiemi di $\Z$: \emph{ogni sottoinsieme non vuoto di
$\Z$ limitato inferiormente ha minimo, e ogni sottoinsieme non vuoto di $\Z$ limitato
superiormente ha massimo}. Basta «traslare» l'insieme dentro $\N$: se $A\subseteq\Z$ ha un
minorante $\ell$ e $k\in\N$ è tale che $k>-\ell$, l'insieme $\{a+k : a\in A\}$ è contenuto
in $\N$, e il suo minimo, diminuito di $k$, è il minimo di $A$ (l'esistenza di un tale $k$
è garantita dalla proprietà di Archimede, che dimostriamo subito e la cui dimostrazione
non usa questa osservazione, quindi non c'è circolarità). Per il massimo si applica il
risultato sul minimo all'insieme $-A=\{-a : a\in A\}$, che è limitato inferiormente se
$A$ è limitato superiormente.
\end{osservazione}

\subsection{La proprietà di Archimede}

Tenendo conto delle proprietà precedenti e del teorema di esistenza dell'estremo
superiore si può dimostrare il risultato seguente, che nel \cref{cap:numeri} abbiamo già
usato una volta (nella dimostrazione del \cref{teo:Q-non-completo}) aggirandolo.

\begin{teorema}[Proprietà di Archimede]\label{teo:archimede}\index{proprietà di Archimede}
Per ogni $x\in\R$ esiste $n\in\N$ tale che $n>x$.
\end{teorema}

\begin{schema}
\Ipotesi{$x\in\R$.}
\Tesi{$\esiste n\in\N : n>x$. \ Equivalentemente: $\N$ non è limitato superiormente,
$\sup\N=+\infty$.}
\end{schema}

\begin{strategia}
La tesi dice che $\N$ non ha maggioranti; ragioniamo per assurdo supponendo che ne abbia.
Allora, per il \cref{teo:esistenza-sup}, $\N$ avrebbe un \emph{minimo} maggiorante $M$.
Ma $\N$ è «invariante per traslazione di $1$»: se $n\in\N$ anche $n+1\in\N$. Quindi
$n+1\le M$ per ogni $n$, cioè $n\le M-1$ per ogni $n$: anche $M-1$ sarebbe un
maggiorante, e più piccolo di $M$. Assurdo. Si noti che l'ingrediente decisivo è la
completezza: senza l'esistenza del sup il ragionamento non parte.
\end{strategia}

\begin{dimostrazione}
Se la proprietà di Archimede fosse falsa, esisterebbe $x\in\R$ con $n\le x$ per ogni
$n\in\N$: l'insieme $\N$ sarebbe limitato superiormente e quindi, per il
\cref{teo:esistenza-sup}, dotato di estremo superiore $M\in\R$. In particolare, per ogni
numero naturale $n$ sarebbe $n\le M$. Poiché anche $n+1$ è un numero naturale,
risulterebbe $n+1\le M$, cioè $n\le M-1$, per ogni $n\in\N$; il che è assurdo perché
$M-1$ sarebbe un maggiorante di $\N$, contrariamente al fatto che $M$ è il più piccolo dei
maggioranti.
\end{dimostrazione}

In altre parole, la proprietà di Archimede afferma che l'insieme $\N$ dei numeri naturali
non è limitato superiormente. Ne seguono alcune conseguenze di uso continuo.

\begin{corollario}\label{cor:archimede}
\leavevmode
\begin{enumerate}
  \item Per ogni $\epsilon>0$ esiste $n\in\N$ tale che $\dfrac1n<\epsilon$; quindi
  $\inf\set{\dfrac1n : n\in\N}=0$.
  \item Se $a>0$ e $b\in\R$, esiste $n\in\N$ tale che $na>b$.
\end{enumerate}
\end{corollario}

\begin{schema}
\Ipotesi{(i) $\epsilon>0$; \ (ii) $a>0$, \ $b\in\R$.}
\Tesi{(i) $\esiste n\in\N : \frac1n<\epsilon$, \ e quindi $\inf\{1/n : n\in\N\}=0$; \
(ii) $\esiste n\in\N : na>b$.}
\end{schema}

\begin{strategia}
Entrambe le affermazioni sono la proprietà di Archimede applicata a un numero reale
scelto in modo opportuno: $1/\epsilon$ per la (i), $b/a$ per la (ii). Le disuguaglianze
richieste si ottengono poi moltiplicando o dividendo per numeri positivi, che non
cambiano il verso. Per l'estremo inferiore si usa la caratterizzazione
\eqref{eq:caratt-inf}.
\end{strategia}

\begin{dimostrazione}
(i) Per il \cref{teo:archimede} esiste $n>1/\epsilon$, cioè $1/n<\epsilon$. Poiché
$1/n>0$ per ogni $n$, lo zero è un minorante di $\{1/n\}$, e per quanto appena visto è
il massimo dei minoranti (\eqref{eq:caratt-inf}). (ii) Per il \cref{teo:archimede}
esiste $n>b/a$; moltiplicando per $a>0$ si ottiene $na>b$.
\end{dimostrazione}

\begin{idea}[Il significato fisico della proprietà di Archimede]
La forma (ii) del corollario è quella originale, geometrica: \emph{dati due segmenti, un
opportuno multiplo del più corto supera il più lungo}. Con un righello da $30$ cm si può
misurare (riportandolo un numero sufficiente di volte) la lunghezza di qualsiasi
corridoio; con un passo di integrazione $\Delta t$, per piccolo che sia, una simulazione
numerica raggiunge in un numero finito di passi qualsiasi istante finale. In una
«grandezza non archimedea» esisterebbero invece segmenti \emph{infinitesimi}, che
riportati un numero qualsiasi di volte non raggiungono mai un segmento dato: i numeri
reali non contengono quantità di questo tipo.
\end{idea}

\begin{notastorica}[Archimede, Eudosso e gli infinitesimi]
Il principio porta il nome di Archimede di Siracusa (287--212 a.C.), che lo enuncia come
postulato nel trattato \emph{Sulla sfera e il cilindro} e lo usa sistematicamente nel
\emph{metodo di esaustione}, con cui calcolò l'area del segmento di parabola e il volume
della sfera; nella prefazione alla \emph{Quadratura della parabola} lo stesso Archimede
osserva però che il lemma era già stato usato dai geometri precedenti, e la tradizione
lo attribuisce a Eudosso di Cnido (IV secolo a.C.); una sua forma compare nella definizione~4 del libro~V degli \emph{Elementi} di Euclide.
Per tutto il Seicento e il Settecento, Leibniz e i suoi successori calcolarono con gli
«infinitesimi», quantità più piccole di ogni numero positivo eppure non nulle: grandezze,
appunto, non archimedee. L'Analisi dell'Ottocento li bandì, sostituendoli con la nozione
di limite. Furono riabilitati in modo rigoroso solo nel 1961 da Abraham Robinson con
l'\emph{analisi non standard}, che lavora in un campo ordinato più grande di $\R$ nel quale
la proprietà di Archimede, per l'appunto, non vale.
\end{notastorica}

\subsection{La parte intera}

\begin{proposizione}[Parte intera]\label{prop:parte-intera}\index{parte intera}
Per ogni $x\in\R$ esiste uno e un solo $m\in\Z$ tale che
\begin{equation}\label{eq:parte-intera}
  m\le x<m+1.
\end{equation}
Il numero $m$ si chiama \emph{parte intera} di $x$ e si indica con $\lfloor x\rfloor$
(oppure con $[x]$).
\end{proposizione}

\begin{schema}
\Ipotesi{$x\in\R$.}
\Tesi{$\esiste!\,m\in\Z : m\le x<m+1$.}
\end{schema}

\begin{strategia}
Per l'esistenza, $m$ dev'essere «il più grande intero $\le x$»: lo otterremo come
massimo dell'insieme degli interi minori o uguali a $x$. Per non uscire da $\N$, dove
abbiamo i teoremi già pronti, trasliamo tutto di un naturale $k$ abbastanza grande
(fornito da Archimede) e usiamo invece il \emph{minimo} dei naturali che superano $x+k$.
Per l'unicità si osserva che due interi che soddisfano la \eqref{eq:parte-intera}
distano meno di $1$, e due interi distinti distano almeno $1$.
\end{strategia}

\begin{dimostrazione}
\emph{Esistenza.} Per il \cref{teo:archimede} esiste $k\in\N$ con $k>-x$, cioè $x+k>0$,
e (di nuovo per Archimede) l'insieme $T=\{j\in\N : j>x+k\}$ è non vuoto. Per il
\cref{teo:buon-ordinamento} $T$ ha minimo $j_{0}$. Allora $j_{0}>x+k$; inoltre
$j_{0}-1\le x+k$: se $j_{0}=1$ ciò segue da $x+k>0$, altrimenti $j_{0}-1\in\N$
(\cref{lem:struttura-N}(ii)) e $j_{0}-1\notin T$ perché è minore del minimo. Posto
$m=j_{0}-1-k\in\Z$, le due disuguaglianze $j_{0}-1\le x+k<j_{0}$ diventano $m\le x<m+1$.

\emph{Unicità.} Se anche $m'\in\Z$ verifica $m'\le x<m'+1$, sommando $m\le x<m+1$ con
$-m'-1<-x\le-m'$ si ottiene $-1<m-m'<1$. Ma un intero non nullo $h$ ha $\abs{h}\in\N$,
quindi $\abs{h}\ge1$ (\cref{lem:struttura-N}(i)); dunque $m-m'=0$.
\end{dimostrazione}

\begin{marginfigure}
\centering
\begin{tikzpicture}
\begin{axis}[width=5.4cm, height=4.6cm, xmin=-2.4, xmax=3.4, ymin=-2.6, ymax=3.4,
  axis lines=left, xtick={-2,-1,0,1,2,3}, ytick={-2,-1,0,1,2,3}]
  \foreach \k in {-2,...,2} {
    \addplot[curva, domain=\k:\k+0.97] {\k};
    \addplot[only marks, mark=*, mark size=1.2pt, blu] coordinates {(\k,\k)};
    \addplot[only marks, mark=o, mark size=1.2pt, blu, fill=white] coordinates {(\k+1,\k)};}
\end{axis}
\end{tikzpicture}
\caption{Il grafico della funzione parte intera $y=\lfloor x\rfloor$: una «scala» con
gradini chiusi a sinistra e aperti a destra.}
\label{fig:parte-intera}
\end{marginfigure}

Ad esempio $\lfloor 2{,}7\rfloor=2$, $\lfloor 3\rfloor=3$, $\lfloor-2{,}7\rfloor=-3$
(\emph{non} $-2$: $-3\le-2{,}7<-2$). In Python la parte intera è
\texttt{math.floor}, da non confondere con \texttt{int}, che tronca verso lo zero:
\texttt{int(-2.7)} restituisce \texttt{-2}.

\subsection{La densità dei razionali}

Dalla proprietà di Archimede si ricava che l'insieme $\Q$ dei numeri razionali è
\emph{denso} in $\R$.

\begin{definizione}[Insieme denso]\index{insieme!denso}
Un insieme $D\subseteq\R$ si dice \emph{denso} in $\R$ se per ogni coppia $a,b$ di numeri
reali con $a<b$ esiste un elemento di $D$ compreso fra $a$ e $b$.
\end{definizione}

\begin{teorema}[Densità dei numeri razionali]\label{teo:densita-Q}\index{densità!dei razionali}
L'insieme $\Q$ dei numeri razionali è denso in $\R$: per ogni $a,b\in\R$ con $a<b$
esiste $q\in\Q$ tale che $a<q<b$.
\end{teorema}

\begin{schema}
\Ipotesi{$a,b\in\R$, \ $a<b$.}
\Tesi{$\esiste m\in\Z,\ n\in\N : \ a<\dfrac{m}{n}<b$.}
\end{schema}

\begin{strategia}
Cerchiamo una frazione $m/n$ nell'intervallo $(a,b)$. L'idea è di scegliere prima il
\emph{denominatore} $n$ in modo che la «griglia» delle frazioni di denominatore $n$,
cioè i punti $\dots,-\frac1n,0,\frac1n,\frac2n,\dots$, abbia un passo $\frac1n$ più
piccolo della lunghezza $b-a$ dell'intervallo: allora almeno un punto della griglia deve
cadere dentro (\cref{fig:densita-Q}). Il denominatore lo fornisce Archimede ($n>1/(b-a)$).
Come numeratore prendiamo il primo intero che supera $na$, cioè $m=\lfloor na\rfloor+1$:
è il primo punto della griglia a destra di $a$, e non può «scavalcare» $b$ perché il passo
è troppo corto.
\end{strategia}

\begin{dimostrazione}
Per il \cref{teo:archimede} esiste $n\in\N$ con $n>\frac{1}{b-a}$; per cui risulterà
$nb-na>1$. Sia $m=\lfloor na\rfloor+1$ (\cref{prop:parte-intera}); allora
$m-1\le na<m$, e quindi
\[
  na<m=(m-1)+1\le na+1<na+(nb-na)=nb .
\]
Dalle disuguaglianze $na<m<nb$, dividendo per $n>0$, segue $a<\frac mn<b$.
\end{dimostrazione}

\begin{figure}[htb]
\centering
\begin{tikzpicture}[x=1.05cm, every node/.style={font=\footnotesize}]
  \draw[-Latex, grigio] (-0.3,0) -- (10.5,0);
  \foreach \k in {0,...,10} {\draw[grigio] (\k,0.09) -- (\k,-0.09);}
  \node[below=3pt, grigio] at (4,0) {$\frac{m-1}{n}$};
  \node[below=3pt, verde] at (5,0) {$\frac{m}{n}$};
  \node[below=3pt, grigio] at (6,0) {$\frac{m+1}{n}$};
  \fill[blu] (4.35,0) circle (2pt) node[above=3pt, blu] {$a$};
  \fill[blu] (5.7,0) circle (2pt) node[above=3pt, blu] {$b$};
  \fill[verde] (5,0) circle (2.3pt);
  \draw[grigio, decorate, decoration={brace, amplitude=3pt}] (7,0.25) -- (8,0.25) node[midway, above=3pt] {passo $\frac1n<b-a$};
  \draw[blu, thick] (4.35,-0.55) -- (5.7,-0.55);
  \node[blu, below] at (5.03,-0.6) {$b-a$};
\end{tikzpicture}
\caption{L'idea della dimostrazione del \cref{teo:densita-Q}: se il passo $1/n$ della
griglia delle frazioni di denominatore $n$ è più corto dell'intervallo $(a,b)$, almeno
un punto $m/n$ della griglia cade all'interno.}
\label{fig:densita-Q}
\end{figure}

\begin{corollario}[Densità degli irrazionali]\label{cor:densita-irrazionali}
Anche l'insieme $\R\setminus\Q$ dei numeri irrazionali è denso in $\R$.
\end{corollario}

\begin{schema}
\Ipotesi{$a,b\in\R$, \ $a<b$.}
\Tesi{$\esiste t\in\R\setminus\Q : a<t<b$.}
\end{schema}

\begin{strategia}
Non occorre ripetere la costruzione: basta «spostare» il problema. Traslato
l'intervallo $(a,b)$ di $-\sqrt2$, il teorema di densità vi trova un razionale $q$;
riportato indietro, $q+\sqrt2$ cade in $(a,b)$ ed è irrazionale. Al posto di $\sqrt2$
andrebbe bene qualsiasi irrazionale fissato.
\end{strategia}

\begin{dimostrazione}
Siano $a<b$. Applicando il \cref{teo:densita-Q} alla coppia $a-\sqrt2<b-\sqrt2$ si trova
$q\in\Q$ con $a-\sqrt2<q<b-\sqrt2$, cioè $a<q+\sqrt2<b$ (l'esistenza di $\sqrt2$ in
$\R$ è garantita dal \cref{teo:radice-n}); e $q+\sqrt2$ è irrazionale,
perché somma di un razionale e di un irrazionale (esercizio del \cref{sec:nzq}).
\end{dimostrazione}

\begin{osservazione}[Denso non vuol dire «pieno»]
Il \cref{teo:densita-Q} e il \cref{cor:densita-irrazionali} dicono che razionali e
irrazionali sono «mescolati» in modo inestricabile: in ogni intervallo, per piccolo che
sia, ci sono infiniti razionali e infiniti irrazionali. Eppure $\Q$ ha i «buchi» del
\cref{teo:Q-non-completo}: essere denso non significa essere completo. È una distinzione
sottile che si rivela profonda: Cantor dimostrò nel 1874 che i razionali si possono
mettere in fila e numerare (sono un'infinità \emph{numerabile}), mentre i reali no. In
un senso preciso, «quasi tutti» i numeri reali sono irrazionali.
\end{osservazione}

\begin{python}[La dimostrazione della densità, trasformata in un algoritmo]
# Il teorema di densità è costruttivo: la dimostrazione dice
# esattamente quali n e m scegliere. Traduciamola in codice, con
# l'aritmetica esatta delle frazioni.
from fractions import Fraction as F
import math

def razionale_tra(a, b):
    """Restituisce m/n con a < m/n < b, seguendo la dimostrazione."""
    assert a < b
    n = math.floor(1 / (b - a)) + 1     # Archimede: n > 1/(b - a)
    m = math.floor(n * a) + 1           # primo intero che supera n*a
    q = F(m, n)
    assert a < q < b                    # la tesi del teorema
    return q

print(razionale_tra(F(1, 3), F(1, 2)))         # tra 1/3 e 1/2
print(razionale_tra(math.sqrt(2), 1.4143))     # tra √2 e 1.4143
print(razionale_tra(math.pi, math.pi + 1e-6))  # intervallo minuscolo
\end{python}
\begin{risultato}
3/7
16363/11570
3141593/1000000
\end{risultato}

\begin{esercizio}
Dimostrare che tra due numeri reali distinti esistono \emph{infiniti} numeri razionali.
\begin{soluzione}
Siano $a<b$. Per il \cref{teo:densita-Q} esiste $q_{1}\in(a,b)$ razionale; applicandolo
di nuovo a $(a,q_{1})$ si trova $q_{2}\in(a,q_{1})$, poi $q_{3}\in(a,q_{2})$ e così via:
per induzione si costruisce una successione di razionali $q_{1}>q_{2}>q_{3}>\cdots$
tutti distinti e tutti in $(a,b)$. \emph{Scelta logica:} «infiniti» si dimostra
esibendo un procedimento che non si arresta mai; il teorema di densità è il «motore»
che a ogni passo produce un elemento nuovo.
\end{soluzione}
\end{esercizio}

\subsection{Approfondimento: i numeri reali in forma decimale}
\label{subsec:forma-decimale}

Fin dalla scuola siamo abituati a pensare un numero reale come un \emph{allineamento
decimale}, eventualmente illimitato: $\frac14=0{,}25$, $\frac13=0{,}333\ldots$,
$\sqrt2=1{,}41421356\ldots$. Negli assiomi del \cref{sec:assiomi}, però, di cifre non si
parla affatto, e i puntini non hanno per ora alcun significato preciso. Seguendo
\cite{giusti}, mostriamo che la rappresentazione decimale è una \emph{conseguenza} degli
assiomi: le cifre di un numero si ricavano dalla parte intera, e il numero si ricostruisce
dalle sue cifre con l'estremo superiore. Le serie numeriche, con le quali un allineamento
decimale si legge come una somma infinita, arriveranno solo nel \cref{cap:serie}: qui non
ne avremo bisogno. Ci limitiamo ai numeri $x\ge0$; per $x<0$ si rappresenta $-x$ e si
antepone il segno meno.

\begin{definizione}[Troncamenti e cifre decimali]\label{def:allineamento-decimale}\index{rappresentazione decimale}\index{cifre decimali}
Sia $x\ge0$. Per ogni intero $n\ge0$ il numero razionale
\begin{equation}\label{eq:troncamento}
  x_{n}=\frac{\lfloor10^{n}x\rfloor}{10^{n}}
\end{equation}
si chiama \emph{troncamento} di $x$ alla $n$-esima cifra decimale. Si pone poi
$c_{0}=\lfloor x\rfloor$ e, per ogni $k\in\N$,
\begin{equation}\label{eq:cifra-decimale}
  c_{k}=\lfloor10^{k}x\rfloor-10\,\lfloor10^{k-1}x\rfloor ;
\end{equation}
$c_{k}$ è la \emph{$k$-esima cifra decimale} di $x$, e si scrive
$x=c_{0}{,}c_{1}c_{2}c_{3}\ldots$ (\emph{allineamento decimale} di $x$), dove l'intero
$c_{0}$ si scrive nel modo usuale.
\end{definizione}

Ad esempio, per $x=\frac{27}{8}$ si trova $\lfloor x\rfloor=3$, $\lfloor10x\rfloor=33$,
$\lfloor100x\rfloor=337$, $\lfloor1000x\rfloor=3375$ e poi
$\lfloor10^{k}x\rfloor=3375\cdot10^{k-3}$; quindi $c_{1}=33-30=3$, $c_{2}=337-330=7$,
$c_{3}=3375-3370=5$ e $c_{k}=0$ per $k\ge4$: $\frac{27}{8}=3{,}375$. La formula
\eqref{eq:cifra-decimale} non fa che leggere la cifra delle unità dell'intero
$\lfloor10^{k}x\rfloor$: moltiplicare per $10^{k}$ sposta la virgola di $k$ posti verso
destra, e la parte intera butta via ciò che resta dopo la virgola.

\begin{teorema}[Esistenza dell'allineamento decimale]\label{teo:cifre-decimali}
Sia $x\ge0$, con troncamenti $x_{n}$ e cifre $c_{k}$ come nella
\cref{def:allineamento-decimale}. Allora:
\begin{enumerate}
  \item $c_{k}\in\{0,1,\dots,9\}$ per ogni $k\in\N$;
  \item per ogni $n\ge0$ (la somma è vuota se $n=0$)
  \begin{equation}\label{eq:troncamento-stima}
    x_{n}=c_{0}+\sum_{k=1}^{n}\frac{c_{k}}{10^{k}},
    \qquad x_{n}\le x<x_{n}+\frac{1}{10^{n}} ;
  \end{equation}
  \item $x=\sup\{x_{n} : n\ge0\}$;
  \item le cifre non sono \emph{definitivamente uguali a $9$}: non esiste $m\in\N$ tale
  che $c_{k}=9$ per ogni $k>m$.
\end{enumerate}
\end{teorema}

\begin{schema}
\Ipotesi{$x\ge0$; \ $x_{n}=\lfloor10^{n}x\rfloor/10^{n}$, \ $c_{0}=\lfloor x\rfloor$, \
$c_{k}=\lfloor10^{k}x\rfloor-10\lfloor10^{k-1}x\rfloor$.}
\Tesi{(i) $c_{k}\in\{0,\dots,9\}$; \ (ii) $x_{n}=c_{0}+\sum_{k=1}^{n}c_{k}10^{-k}$ e
$x_{n}\le x<x_{n}+10^{-n}$; \ (iii) $x=\sup_{n}x_{n}$; \ (iv) per nessun $m$ si ha
$c_{k}=9$ per ogni $k>m$.}
\end{schema}

\begin{strategia}
Tutto discende dalla doppia disuguaglianza che definisce la parte intera
(\cref{prop:parte-intera}), scritta per l'intero $m_{n}=\lfloor10^{n}x\rfloor$:
$m_{n}\le10^{n}x<m_{n}+1$. Dividendo per $10^{n}$ si ottiene la stima (ii):
\emph{il troncamento approssima $x$ per difetto con un errore minore di $10^{-n}$}.
Moltiplicando invece per $10$ la disuguaglianza scritta per $m_{k-1}$ si chiude $10^{k}x$
tra due interi che distano $10$, e questo costringe la cifra $c_{k}$ tra $0$ e $9$; la
formula di $x_{n}$ come somma delle cifre è una somma telescopica. Il punto (iii) è la
caratterizzazione dell'estremo superiore, con la proprietà di Archimede a garantire che
l'errore $10^{-n}$ diventi piccolo quanto si vuole. Il punto (iv) si prova per assurdo:
una coda di soli $9$ spingerebbe i troncamenti tanto vicino al numero $x_{m}+10^{-m}$ da
costringere $x$ a coincidere con esso, mentre la stima (ii) lo vieta.
\end{strategia}

\begin{dimostrazione}
Poniamo $m_{n}=\lfloor10^{n}x\rfloor\in\Z$, cosicché $x_{n}=m_{n}/10^{n}$ e, per la
\eqref{eq:parte-intera},
\begin{equation}\label{eq:mn-parte-intera}
  m_{n}\le10^{n}x<m_{n}+1 \qquad\text{per ogni } n\ge0 .
\end{equation}
Useremo un fatto elementare: \emph{se $j\in\Z$ e $j\le t$, allora $j\le\lfloor t\rfloor$}.
Infatti, se fosse $j>\lfloor t\rfloor$, l'intero $j-\lfloor t\rfloor$ sarebbe positivo,
quindi $\ge1$ (\cref{lem:struttura-N}(i)), e risulterebbe $j\ge\lfloor t\rfloor+1>t$.

(i) Moltiplicando per $10$ la \eqref{eq:mn-parte-intera} scritta con $n=k-1$ si ottiene
$10\,m_{k-1}\le10^{k}x<10\,m_{k-1}+10$. L'intero $10\,m_{k-1}$ non supera $10^{k}x$, e
quindi $10\,m_{k-1}\le m_{k}$; d'altra parte $m_{k}\le10^{k}x<10\,m_{k-1}+10$. Dunque
$c_{k}=m_{k}-10\,m_{k-1}$ è un intero con $0\le c_{k}<10$, cioè $c_{k}\in\{0,1,\dots,9\}$.

(ii) Dividendo la \eqref{eq:mn-parte-intera} per $10^{n}>0$ si ha
$x_{n}\le x<x_{n}+10^{-n}$. Inoltre $c_{k}10^{-k}=m_{k}10^{-k}-m_{k-1}10^{-(k-1)}
=x_{k}-x_{k-1}$, e la somma è telescopica:
\[
  c_{0}+\sum_{k=1}^{n}\frac{c_{k}}{10^{k}}=x_{0}+\sum_{k=1}^{n}(x_{k}-x_{k-1})=x_{n}.
\]

(iii) Per la (ii), $x$ è un maggiorante dell'insieme dei troncamenti. Osserviamo poi che
\begin{equation}\label{eq:dieci-meno-n}
  \perogni\epsilon>0\ \ \esiste n\in\N : \ \frac{1}{10^{n}}<\epsilon ,
\end{equation}
e che la stessa disuguaglianza vale allora per ogni indice più grande, perché le potenze
di $10$ crescono. Infatti per il \cref{cor:archimede}(i) esiste $n$ con $\frac1n<\epsilon$,
e per la disuguaglianza di Bernoulli (\cref{teo:bernoulli}) $10^{n}=(1+9)^{n}\ge1+9n>n$,
quindi $\frac{1}{10^{n}}<\frac1n<\epsilon$. Fissato $\epsilon>0$ e scelto $n$ come nella
\eqref{eq:dieci-meno-n}, la (ii) dà $x_{n}>x-10^{-n}>x-\epsilon$. Per la
\cref{prop:caratt-sup}, $x=\sup_{n}x_{n}$.

(iv) Supponiamo per assurdo che sia $c_{k}=9$ per ogni $k>m$, e poniamo $y=x_{m}+10^{-m}$.
Per $n>m$, usando la (ii) e l'identità $9\cdot10^{-k}=10^{-(k-1)}-10^{-k}$, che rende
telescopica anche questa somma,
\[
  x_{n}=x_{m}+\sum_{k=m+1}^{n}\frac{9}{10^{k}}=x_{m}+\frac{1}{10^{m}}-\frac{1}{10^{n}}
  =y-\frac{1}{10^{n}} .
\]
D'altra parte la (ii) con $n=m$ dà $x<y$, e con $n>m$ dà $x\ge x_{n}=y-10^{-n}$. Dunque
$0<y-x\le10^{-n}$ per ogni $n>m$, contro la \eqref{eq:dieci-meno-n} con $\epsilon=y-x$.
\end{dimostrazione}

\begin{figure}[htb]
\centering
\begin{tikzpicture}[x=0.8cm, every node/.style={font=\scriptsize}]
  % prima riga: [1, 2] diviso in decimi
  \draw[grigio] (0,0) -- (10,0);
  \foreach \k in {0,...,10} {\draw[grigio] (\k,0.08) -- (\k,-0.08);}
  \fill[rame!30] (4,-0.09) rectangle (5,0.09);
  \draw[rame, thick] (4,0) -- (5,0);
  \foreach \p/\lab in {0/1, 4/{1{,}4}, 5/{1{,}5}, 10/2} {\node[below=2pt] at (\p,0) {$\lab$};}
  \fill[blu] (4.142,0) circle (1.5pt) node[above=2pt] {$\sqrt2$};
  \node[anchor=west] at (10.3,0) {$x_{1}=1{,}4$, \ $c_{1}=4$};
  % seconda riga: [1,4; 1,5] diviso in centesimi
  \begin{scope}[yshift=-1.7cm]
  \draw[grigio] (0,0) -- (10,0);
  \foreach \k in {0,...,10} {\draw[grigio] (\k,0.08) -- (\k,-0.08);}
  \fill[rame!30] (1,-0.09) rectangle (2,0.09);
  \draw[rame, thick] (1,0) -- (2,0);
  \foreach \p/\lab in {0/{1{,}4}, 1/{1{,}41}, 2/{1{,}42}, 10/{1{,}5}} {\node[below=2pt] at (\p,0) {$\lab$};}
  \fill[blu] (1.421,0) circle (1.5pt) node[above=2pt] {$\sqrt2$};
  \node[anchor=west] at (10.3,0) {$x_{2}=1{,}41$, \ $c_{2}=1$};
  \end{scope}
  % terza riga: [1,41; 1,42] diviso in millesimi
  \begin{scope}[yshift=-3.4cm]
  \draw[grigio] (0,0) -- (10,0);
  \foreach \k in {0,...,10} {\draw[grigio] (\k,0.08) -- (\k,-0.08);}
  \fill[rame!30] (4,-0.09) rectangle (5,0.09);
  \draw[rame, thick] (4,0) -- (5,0);
  \foreach \p/\lab in {0/{1{,}41}, 10/{1{,}42}} {\node[below=2pt] at (\p,0) {$\lab$};}
  \node[below=2pt, anchor=north east] at (4.1,0) {$1{,}414$};
  \node[below=2pt, anchor=north west] at (4.9,0) {$1{,}415$};
  \fill[blu] (4.2136,0) circle (1.5pt) node[above=2pt] {$\sqrt2$};
  \node[anchor=west] at (10.3,0) {$x_{3}=1{,}414$, \ $c_{3}=4$};
  \end{scope}
  % la lente: l'intervallo evidenziato diventa la riga successiva
  \draw[grigio, dashed] (4,-0.5) -- (0,-1.45);
  \draw[grigio, dashed] (5,-0.5) -- (10,-1.45);
  \draw[grigio, dashed] (1,-2.2) -- (0,-3.15);
  \draw[grigio, dashed] (2,-2.2) -- (10,-3.15);
\end{tikzpicture}
\caption{Le cifre decimali come una lente d'ingrandimento. A ogni passo l'intervallo
$[x_{n},x_{n}+10^{-n})$ che contiene $x=\sqrt2$ (in color rame) viene diviso in dieci
parti uguali, e la cifra $c_{n+1}$ dice in quale delle dieci cade $x$. I troncamenti
$x_{n}$, estremi sinistri degli intervalli, crescono e hanno $x$ come estremo superiore
(\cref{teo:cifre-decimali}).}
\label{fig:decimali-lente}
\end{figure}

Il teorema precedente va dal numero alle cifre. Il cammino inverso, dalle cifre al numero,
è quello in cui interviene in modo essenziale la completezza di $\R$: assegnate
arbitrariamente le cifre, occorre \emph{trovare} un numero che le abbia.

\begin{teorema}[Ogni allineamento rappresenta un numero reale]\label{teo:allineamento-numero}
Siano $c_{0}\ge0$ un intero e $c_{1},c_{2},c_{3},\dots$ cifre di $\{0,1,\dots,9\}$, e
poniamo
\[
  s_{n}=c_{0}+\sum_{k=1}^{n}\frac{c_{k}}{10^{k}} \qquad (n\ge0).
\]
Allora l'insieme $\{s_{n} : n\ge0\}$ è limitato superiormente, e il suo estremo superiore
$s$ verifica $s_{n}\le s\le s_{n}+10^{-n}$ per ogni $n$. Inoltre:
\begin{enumerate}
  \item se le cifre $c_{k}$ non sono definitivamente uguali a $9$, l'allineamento decimale
  di $s$ è proprio $c_{0}{,}c_{1}c_{2}c_{3}\ldots$, ed $s$ è l'unico numero reale con
  questa proprietà;
  \item se invece $c_{k}=9$ per ogni $k>m$, allora $s=s_{m}+10^{-m}$, e le cifre di $s$
  sono nulle dalla $(m+1)$-esima in poi.
\end{enumerate}
\end{teorema}

\begin{schema}
\Ipotesi{$c_{0}\in\Z$, $c_{0}\ge0$; \ $c_{k}\in\{0,\dots,9\}$ per $k\in\N$; \
$s_{n}=c_{0}+\sum_{k=1}^{n}c_{k}10^{-k}$.}
\Tesi{Esiste $s=\sup_{n}s_{n}$, e $s_{n}\le s\le s_{n}+10^{-n}$; \ (i) se le cifre non
sono definitivamente $9$, le cifre di $s$ sono le $c_{k}$, ed $s$ è l'unico reale con
queste cifre; \ (ii) se $c_{k}=9$ per $k>m$, allora $s=s_{m}+10^{-m}$, con cifre nulle
dopo la $m$-esima.}
\end{schema}

\begin{strategia}
Le somme $s_{n}$ crescono, perché si aggiungono termini $\ge0$, ma non possono crescere
troppo: ciò che si aggiunge dopo l'indice $n$ è al più $\sum_{k>n}9\cdot10^{-k}$, una
somma telescopica minore di $10^{-n}$. Quindi $s_{n}+10^{-n}$ è un maggiorante, e il
candidato naturale per il numero cercato è $s=\sup_{n}s_{n}$: è la completezza di $\R$
(\cref{teo:esistenza-sup}) a garantire che \emph{esista}. Per il punto (i) si migliora la
stima in una disuguaglianza stretta, $s<s_{n}+10^{-n}$, sfruttando una cifra $c_{j}\le8$
con $j>n$; a quel punto $10^{n}s$ è stretto tra l'intero $10^{n}s_{n}$ e il suo
successivo, e l'unicità della parte intera dice che $s_{n}$ è il troncamento di $s$. Il
punto (ii) è un calcolo diretto, chiuso dalla caratterizzazione dell'estremo superiore.
\end{strategia}

\begin{dimostrazione}
Le $s_{n}$ crescono, perché $s_{n}-s_{n-1}=c_{n}10^{-n}\ge0$. Se $m>n$, poiché
$c_{k}\le9$ e $9\cdot10^{-k}=10^{-(k-1)}-10^{-k}$,
\begin{equation}\label{eq:coda-decimale}
  s_{m}-s_{n}=\sum_{k=n+1}^{m}\frac{c_{k}}{10^{k}}\le\sum_{k=n+1}^{m}\frac{9}{10^{k}}
  =\frac{1}{10^{n}}-\frac{1}{10^{m}}<\frac{1}{10^{n}} .
\end{equation}
Dunque $s_{m}<s_{n}+10^{-n}$ se $m>n$, e $s_{m}\le s_{n}$ se $m\le n$: il numero
$s_{n}+10^{-n}$ è un maggiorante di tutte le $s_{m}$. In particolare (con $n=0$)
l'insieme delle $s_{m}$ è limitato superiormente da $c_{0}+1$, e per il
\cref{teo:esistenza-sup} esiste $s=\sup_{m}s_{m}$; si ha $s\ge s_{n}$ perché $s$ è un
maggiorante, e $s\le s_{n}+10^{-n}$ perché $s$ è il minimo dei maggioranti.

(i) Fissato $n\ge0$, per ipotesi esiste $j>n$ con $c_{j}\le8$. Per $m\ge j$, ripetendo la
stima \eqref{eq:coda-decimale} ma usando $c_{j}\le9-1$ nel termine di indice $j$,
\[
  s_{m}-s_{n}\le\frac{1}{10^{n}}-\frac{1}{10^{m}}-\frac{1}{10^{j}}
  <\frac{1}{10^{n}}-\frac{1}{10^{j}} ;
\]
per $m<j$ si ha $s_{m}\le s_{j}$, e la stessa disuguaglianza vale ancora. Dunque
$s_{n}+10^{-n}-10^{-j}$ è un maggiorante delle $s_{m}$, e
\[
  s_{n}\le s\le s_{n}+\frac{1}{10^{n}}-\frac{1}{10^{j}}<s_{n}+\frac{1}{10^{n}} .
\]
Moltiplicando per $10^{n}$ si ottiene $10^{n}s_{n}\le10^{n}s<10^{n}s_{n}+1$, e
$10^{n}s_{n}=10^{n}c_{0}+\sum_{k=1}^{n}c_{k}10^{n-k}$ è un intero. Per l'unicità nella
\cref{prop:parte-intera}, $\lfloor10^{n}s\rfloor=10^{n}s_{n}$: il troncamento $n$-esimo
di $s$ è $s_{n}$. Allora $\lfloor s\rfloor=s_{0}=c_{0}$ e, per $k\in\N$, la $k$-esima
cifra di $s$ è $10^{k}s_{k}-10\cdot10^{k-1}s_{k-1}=10^{k}(s_{k}-s_{k-1})=c_{k}$. Se infine
$y\ge0$ ha le stesse cifre di $s$, per la \eqref{eq:troncamento-stima} i suoi troncamenti
sono $y_{n}=s_{n}$, e per il \cref{teo:cifre-decimali}(iii)
$y=\sup_{n}y_{n}=\sup_{n}s_{n}=s$.

(ii) Se $c_{k}=9$ per ogni $k>m$, la stima \eqref{eq:coda-decimale} è un'uguaglianza:
$s_{n}=s_{m}+10^{-m}-10^{-n}$ per $n>m$. Il numero $y=s_{m}+10^{-m}$ è quindi un
maggiorante delle $s_{n}$ (per $n\le m$ si ha $s_{n}\le s_{m}<y$) e, dato $\epsilon>0$,
scelto $n>m$ con $10^{-n}<\epsilon$ come nella \eqref{eq:dieci-meno-n}, si ha
$s_{n}>y-\epsilon$:
per la \cref{prop:caratt-sup}, $s=y$. Infine $10^{n}s=10^{n-m}(10^{m}s_{m}+1)$ è intero
per $n\ge m$, quindi $\lfloor10^{n}s\rfloor=10^{n}s$, e per $k>m$ la $k$-esima cifra di
$s$ è $10^{k}s-10\cdot10^{k-1}s=0$.
\end{dimostrazione}

\begin{osservazione}[Le code di $9$ e il legame con le serie]\label{oss:allineamenti-nove}
I due teoremi stabiliscono una corrispondenza biunivoca tra i numeri reali $x\ge0$ e gli
allineamenti decimali $c_{0}{,}c_{1}c_{2}\ldots$ le cui cifre non sono definitivamente
uguali a $9$. Gli allineamenti esclusi non sono «sbagliati»: per il
\cref{teo:allineamento-numero}(ii) rappresentano, tramite l'estremo superiore, numeri che
hanno anche un allineamento finito, ad esempio
\[
  0{,}4999\ldots=\sup\{0{,}4;\ 0{,}49;\ 0{,}499;\ \dots\}=0{,}5 ,
\]
e allo stesso modo $0{,}999\ldots=1$. Nel \cref{cap:serie} un allineamento sarà letto
come somma della serie $\sum_{k}c_{k}10^{-k}$: poiché le somme parziali crescono, la somma
è il loro estremo superiore (\cref{teo:monotone}), e i due punti di vista coincidono. Lì
vedremo anche che un allineamento periodico rappresenta un razionale
(\cref{prop:periodici}); il viceversa è l'\cref{es:decimali-periodici}. Infine, il numero
$10$ non ha nulla di speciale: in base $b\ge2$ le cifre stanno in $\{0,\dots,b-1\}$ e si
escludono le code di cifre $b-1$. I calcolatori usano la base $2$, in cui $\frac1{10}$ è
periodico: per questo \texttt{0.1} non è esattamente un decimo e
\texttt{0.1 + 0.2 == 0.3} è falso (\emph{Attenzione} alla fine del \cref{sec:assiomi}).
\end{osservazione}

\begin{python}[Le cifre di un numero, dalla definizione]
# Le cifre di x dalla definizione c_k = [b^k x] - b [b^(k-1) x],
# in aritmetica intera esatta: m(k) deve restituire [b^k x].
import math
from fractions import Fraction as F

def cifre(m, n, b=10):
    """Restituisce la stringa 'c_0,c_1c_2...c_n' nella base b."""
    c = [m(k) - b * m(k - 1) for k in range(1, n + 1)]
    return f"{m(0)}," + "".join(str(ck) for ck in c)

# [10^k sqrt(2)] è il più grande intero j con j^2 <= 2 * 10^(2k)
print("sqrt(2) =", cifre(lambda k: math.isqrt(2 * 100**k), 40))
# 1/7: le cifre si ripetono con periodo 6
print("1/7     =", cifre(lambda k: math.floor(F(10**k, 7)), 36))
# 1/10 in base 2: periodico, quindi 0.1 non è esatto in un float
print("1/10    =", cifre(lambda k: math.floor(F(2**k, 10)), 36, b=2))
print("valore esatto di 0.1:", F(0.1))
\end{python}
\begin{risultato}
sqrt(2) = 1,4142135623730950488016887242096980785696
1/7     = 0,142857142857142857142857142857142857
1/10    = 0,000110011001100110011001100110011001
valore esatto di 0.1: 3602879701896397/36028797018963968
\end{risultato}

Le cifre di $\sqrt2$ sono ottenute senza numeri in virgola mobile: per un intero
$j\ge0$, $j\le10^{k}\sqrt2$ equivale a $j^{2}\le2\cdot10^{2k}$, e \texttt{math.isqrt}
trova esattamente il più grande di tali $j$. È l'antica estrazione della radice «cifra
per cifra». L'ultima riga mostra il razionale (di denominatore $2^{55}$) che il
calcolatore memorizza al posto di $\frac1{10}$.

\begin{notastorica}[Stevin e l'«arte dei decimi»]
Le frazioni con denominatore $10$, $100$, $1000,\dots$ erano già note a matematici
cinesi e arabi, ma il loro uso sistematico in Europa si deve all'ingegnere fiammingo Simon
Stevin (1548--1620), che nel 1585 pubblicò a Leida un opuscolo di una trentina di pagine,
\emph{De Thiende} («Il decimo»), tradotto in francese lo stesso anno come \emph{La Disme}.
Lo dedicava ad astronomi, agrimensori, stazzatori di botti, zecchieri e mercanti: a
chiunque dovesse fare conti precisi con grandezze continue. Scriveva i decimali indicando
in un cerchietto l'esponente della potenza di un decimo, ad esempio $7{,}3486$ come
7\,\textcircled{\tiny 0}\,3\,\textcircled{\tiny 1}\,4\,\textcircled{\tiny 2}\,8\,%
\textcircled{\tiny 3}\,6\,\textcircled{\tiny 4}; la notazione, ripresa e semplificata da
Clavio e da Nepero, divenne quella attuale. Stevin prevedeva anche che monete, pesi e
misure sarebbero diventati decimali: accadde due secoli dopo, con il sistema metrico. Nella
sua \emph{Arithmétique}, dello stesso 1585, sosteneva infine che le radici e gli
irrazionali sono numeri a pieno titolo, e non quantità «assurde» o «sorde». Era, osserva
\cite{giusti}, la risposta che l'età dei mercanti dava all'incommensurabile scoperto dai
pitagorici: di $\sqrt2$ non si chiedeva che cosa fosse, ma come calcolarlo con tutte le
cifre necessarie. Ne resta traccia nella definizione scolastica di numero irrazionale come
«decimale illimitato non periodico», che prende per definizione ciò che abbiamo appena
dimostrato come teorema: le cifre esistono perché $\R$ è completo, non viceversa.
\end{notastorica}

\begin{esercizio}[I razionali hanno allineamento periodico]\label{es:decimali-periodici}
Sia $x=\frac pq$, con $p\ge0$ intero e $q\in\N$. Dimostrare che le cifre decimali di $x$
sono definitivamente periodiche: esistono $i\ge0$ e $\ell\in\N$ tali che
$c_{k+\ell}=c_{k}$ per ogni $k>i$.
\begin{soluzione}
Siano $m_{k}=\lfloor10^{k}x\rfloor$ e $r_{k}=10^{k}p-q\,m_{k}$: moltiplicando per $q$ la
\eqref{eq:mn-parte-intera} si trova $0\le r_{k}<q$, cioè $r_{k}$ è il resto della
divisione di $10^{k}p$ per $q$. Allora
\[
  10^{k+1}x=\frac{10\,(q\,m_{k}+r_{k})}{q}=10\,m_{k}+\frac{10\,r_{k}}{q},
\]
e poiché $\lfloor h+t\rfloor=h+\lfloor t\rfloor$ per ogni intero $h$ (per l'unicità della
parte intera), $m_{k+1}=10\,m_{k}+\lfloor10\,r_{k}/q\rfloor$. Dunque
\[
  c_{k+1}=\Bigl\lfloor\frac{10\,r_{k}}{q}\Bigr\rfloor, \qquad
  r_{k+1}=10\,r_{k}-q\,c_{k+1} :
\]
cifra e resto successivi dipendono \emph{soltanto} dal resto $r_{k}$: è la divisione in
colonna. I $q+1$ resti $r_{0},\dots,r_{q}$ stanno in $\{0,\dots,q-1\}$, che ha $q$
elementi, quindi due coincidono: $r_{i}=r_{i+\ell}$ con $\ell\ge1$. Da lì la divisione
ricomincia identica: per induzione $r_{k+\ell}=r_{k}$ per $k\ge i$, e quindi
$c_{k+\ell}=c_{k}$ per $k>i$. Per $\frac17$ i resti sono $1,3,2,6,4,5,1,\dots$ e il periodo
è $142857$. \emph{Scelta logica:} invece di una formula per le cifre, si cerca una
grandezza che ne determina il seguito (il resto) e che assume un numero finito di valori:
appena un valore si ripete, si ripete tutto («principio dei cassetti»).
\end{soluzione}
\end{esercizio}

%% ----------------------------------------------------------------------------
%%  §2.3  Radici n-esime e funzione esponenziale a esponente reale
%% ----------------------------------------------------------------------------
\section{Radici \texorpdfstring{$n$}{n}-esime e potenze a esponente reale}
\label{sec:esponenziale-reale}

Nel \cref{sec:potenze} abbiamo definito la radice $n$-esima come funzione inversa della
potenza $x\mapsto x^{n}$, rimandando a più avanti la prova che ogni numero $y\ge0$
\emph{ha} una radice $n$-esima; e abbiamo usato potenze $a^{x}$ con esponente reale senza
averle definite. Con l'estremo superiore possiamo ora colmare entrambe le lacune.

\subsection{L'esistenza della radice \texorpdfstring{$n$}{n}-esima}

Premettiamo una stima elementare, che controlla di quanto cambia una potenza quando
cambia la base.

\begin{lemma}[Stima della differenza di due potenze]\label{lem:stima-potenze}
Siano $0\le v<u$ e $n\in\N$. Allora
\begin{equation}\label{eq:stima-potenze}
  n\,v^{n-1}(u-v)\ \le\ u^{n}-v^{n}\ \le\ n\,u^{n-1}(u-v).
\end{equation}
\end{lemma}

\begin{schema}
\Ipotesi{$0\le v<u$; \ $n\in\N$ (con la convenzione $v^{0}=1$ anche per $v=0$).}
\Tesi{$n\,v^{n-1}(u-v)\le u^{n}-v^{n}\le n\,u^{n-1}(u-v)$.}
\end{schema}

\begin{strategia}
La chiave è la scomposizione
$u^{n}-v^{n}=(u-v)\bigl(u^{n-1}+u^{n-2}v+\dots+uv^{n-2}+v^{n-1}\bigr)$, che generalizza
$u^{2}-v^{2}=(u-v)(u+v)$. Nella parentesi ci sono $n$ addendi, ciascuno dei quali è
compreso tra il più piccolo, $v^{n-1}$, e il più grande, $u^{n-1}$. Stimando ogni
addendo dal basso e dall'alto si ottengono le due disuguaglianze.
\end{strategia}

\begin{dimostrazione}
Per $n=1$ la \eqref{eq:stima-potenze} si riduce a $u-v\le u-v\le u-v$. Sia $n\ge2$.
Poiché $u>0$, possiamo applicare la formula della progressione geometrica
(\cref{prop:geometrica}) con ragione $x=v/u\neq1$ e $n-1$ al posto di $n$:
\[
  1+\frac vu+\Bigl(\frac vu\Bigr)^{2}+\dots+\Bigl(\frac vu\Bigr)^{n-1}
  =\frac{1-(v/u)^{n}}{1-v/u}.
\]
Moltiplicando entrambi i membri per $u^{n-1}(u-v)=u^{n}(1-v/u)$ si trova
\begin{equation}\label{eq:scomposizione-potenze}
  u^{n}-v^{n}=(u-v)\bigl(u^{n-1}+u^{n-2}v+\dots+u\,v^{n-2}+v^{n-1}\bigr).
\end{equation}
Il generico addendo della parentesi è $u^{n-1-k}v^{k}$, con $0\le k\le n-1$; poiché
$0\le v<u$, per la monotonia delle potenze (\cref{prop:potenza-crescente}) risulta
$v^{n-1}\le u^{n-1-k}v^{k}\le u^{n-1}$. Sommando gli $n$ addendi, la parentesi è compresa
tra $n v^{n-1}$ e $n u^{n-1}$; moltiplicando per $u-v>0$ si ottiene la tesi.
\end{dimostrazione}

\begin{teorema}[Esistenza e unicità della radice $n$-esima]\label{teo:radice-n}\index{radice $n$-esima!esistenza}
Per ogni $y\ge0$ e ogni $n\in\N$ esiste uno e un solo numero reale $x\ge0$ tale che
$x^{n}=y$.
\end{teorema}

\begin{schema}
\Ipotesi{$y\in\R$, $y\ge0$; \ $n\in\N$.}
\Tesi{$\esiste!\,x\ge0 : x^{n}=y$ \ (il simbolo $\esiste!$ si legge «esiste uno e un solo»).}
\end{schema}

\begin{strategia}
È la generalizzazione del ragionamento con cui nel \cref{sec:nzq} abbiamo «stretto»
$\sqrt2$ tra i razionali con quadrato minore e maggiore di $2$; ora però, grazie
all'estremo superiore, possiamo \emph{costruire} il numero cercato. Consideriamo
l'insieme $S=\{t\ge0 : t^{n}<y\}$ dei candidati «troppo piccoli» e poniamo $x=\sup S$.
Dimostriamo che $x^{n}=y$ escludendo le altre due possibilità.
\begin{itemize}
  \item Se fosse $x^{n}<y$, $x$ sarebbe ancora «troppo piccolo», e si potrebbe aumentarlo
  di poco, $x+h$, restando in $S$: ma $x$ è un maggiorante di $S$.
  \item Se fosse $x^{n}>y$, $x$ sarebbe «troppo grande», e si potrebbe diminuirlo di
  poco, $x-h$, ottenendo ancora un maggiorante di $S$: ma $x$ è il \emph{minimo}
  maggiorante.
\end{itemize}
In entrambi i casi il lavoro tecnico consiste nello scegliere $h$ abbastanza piccolo,
e il \cref{lem:stima-potenze} dice esattamente quanto. L'unicità segue dalla stretta
monotonia della potenza.
\end{strategia}

\begin{dimostrazione}
Se $y=0$ basta prendere $x=0$. Sia $y>0$ e consideriamo
\[
  S=\{t\in\R : t\ge0,\ t^{n}<y\}.
\]
$S$ è non vuoto: $0\in S$. $S$ è limitato superiormente da $\max\{1,y\}$: se infatti
$t>\max\{1,y\}$, allora $t>1$ e quindi $t^{n}\ge t>y$, cioè $t\notin S$. Per il
\cref{teo:esistenza-sup} esiste $x=\sup S$. Inoltre $x>0$: il numero
$t_{0}=\frac12\min\{1,y\}$ verifica $0<t_{0}<1$, dunque $t_{0}^{n}\le t_{0}<y$, cioè
$t_{0}\in S$ e $x\ge t_{0}>0$.

\emph{Non può essere $x^{n}<y$.} Supponiamolo per assurdo e scegliamo $h$ con
\[
  0<h<1 \qquad\text{e}\qquad h<\frac{y-x^{n}}{n(x+1)^{n-1}} .
\]
Per il \cref{lem:stima-potenze} con $u=x+h$, $v=x$, e poiché $x+h<x+1$,
\[
  (x+h)^{n}\le x^{n}+n(x+h)^{n-1}h\le x^{n}+n(x+1)^{n-1}h<x^{n}+(y-x^{n})=y .
\]
Dunque $x+h\in S$, contro il fatto che $x$ è un maggiorante di $S$ e $x+h>x$.

\emph{Non può essere $x^{n}>y$.} Supponiamolo per assurdo e scegliamo $h$ con
\[
  0<h<x \qquad\text{e}\qquad h<\frac{x^{n}-y}{n\,x^{n-1}} .
\]
Per il \cref{lem:stima-potenze} con $u=x$, $v=x-h\ge0$,
\[
  (x-h)^{n}\ge x^{n}-n\,x^{n-1}h>x^{n}-(x^{n}-y)=y .
\]
Allora $x-h$ è un maggiorante di $S$: se $t\in S$, si ha $t^{n}<y<(x-h)^{n}$, e quindi
$t<x-h$ (se fosse $t\ge x-h\ge0$, per la monotonia delle potenze sarebbe
$t^{n}\ge(x-h)^{n}$). Ma $x-h<x$, contro il fatto che $x$ è il \emph{minimo} dei
maggioranti.

Esclusi entrambi i casi, deve essere $x^{n}=y$: per la dicotomia
\ref{ass:dicotomia} vale $x^{n}\le y$ oppure $x^{n}\ge y$, e in nessuno dei due casi la
disuguaglianza può essere stretta. Infine, se anche $x'\ge0$ verifica
$x'^{n}=y$, non può essere né $x'<x$ né $x'>x$, perché la potenza $n$-esima è
strettamente crescente su $[0,+\infty)$ (\cref{prop:potenza-crescente}): dunque $x'=x$.
\end{dimostrazione}

\begin{marginfigure}
\centering
\begin{tikzpicture}
\begin{axis}[width=5.3cm, height=4.8cm, xmin=-0.15, xmax=1.95, ymin=-0.3, ymax=4.2,
  xtick={1.587}, xticklabels={$x$}, ytick={4}, yticklabels={$y$}]
  \addplot[tratteggio, domain=0:1.9] {4};
  \addplot[curva, domain=0:1.62, samples=80] {x^3};
  \addplot[blu!35, line width=3.5pt] coordinates {(0,0) (1.587,0)};
  \addplot[tratteggio] coordinates {(1.587,0) (1.587,4)};
  \node[font=\scriptsize, blu] at (axis cs:0.32,0.5) {$S$};
  \node[font=\scriptsize, blu] at (axis cs:1.15,2.6) {$t^{n}$};
\end{axis}
\end{tikzpicture}
\caption{La dimostrazione del \cref{teo:radice-n} (qui $n=3$, $y=4$): l'insieme $S$ dei
$t\ge0$ con $t^{n}<y$ è l'intervallo $[0,x)$, e il suo estremo superiore $x$ è la
radice cercata.}
\label{fig:radice-sup}
\end{marginfigure}

\begin{corollario}\label{cor:radice-2}
Esiste un unico numero reale positivo $\sqrt2$ tale che $(\sqrt2)^{2}=2$; esso è
irrazionale (\cref{teo:radice2}). Più in generale, la funzione radice $n$-esima
$x\mapsto\sqrt[n]{x}$ è ben definita su $[0,+\infty)$ ed è l'inversa della potenza
$x\mapsto x^{n}$ ristretta a $[0,+\infty)$, come anticipato nella \eqref{eq:radice}; per
il \cref{teo:monotona-invertibile} è strettamente crescente.
\end{corollario}

\needspace{10\baselineskip}
\begin{schema}
\Ipotesi{$n\in\N$; \ la potenza $x\mapsto x^{n}$ ristretta a $[0,+\infty)$.}
\Tesi{Esiste un unico $\sqrt2>0$ con $(\sqrt2)^{2}=2$, ed è irrazionale; la potenza
$x\mapsto x^{n}$ è una biiezione di $[0,+\infty)$ su sé stesso, la cui inversa
$\sqrt[n]{x}$ è strettamente crescente.}
\end{schema}

\begin{strategia}
Il \cref{teo:radice-n} dice esattamente che la potenza $n$-esima, ristretta a
$[0,+\infty)$, è suriettiva su $[0,+\infty)$ (esistenza) e iniettiva (unicità). Insieme
alla stretta monotonia della potenza, il criterio di invertibilità fornisce l'inversa e
ne garantisce la monotonia. L'affermazione su $\sqrt2$ è il caso $n=2$, $y=2$.
\end{strategia}

\begin{dimostrazione}
Con $n=2$ e $y=2$ il \cref{teo:radice-n} fornisce un unico $x\ge0$ con $x^{2}=2$; si ha
$x>0$ perché $0^{2}=0\neq2$, e $x\notin\Q$ per il \cref{teo:radice2}. In generale, la
funzione $f(x)=x^{n}$ manda $[0,+\infty)$ in $[0,+\infty)$, è suriettiva per il
\cref{teo:radice-n} ed è strettamente crescente per la \cref{prop:potenza-crescente}.
Per il \cref{teo:monotona-invertibile} è invertibile, e l'inversa
$\sqrt[n]{\cdot}\colon[0,+\infty)\to[0,+\infty)$ è strettamente crescente.
\end{dimostrazione}

\begin{osservazione}[Radici di ordine dispari dei numeri negativi]
Se $n$ è dispari, $(-x)^{n}=-x^{n}$ (regola dei segni applicata $n$ volte); quindi per
$y<0$ il numero $x=-\sqrt[n]{-y}$ verifica $x^{n}=y$, ed è l'unico. Per $n$ dispari la
potenza $x\mapsto x^{n}$ è dunque una biiezione di $\R$ su $\R$, e si pone ad esempio
$\sqrt[3]{-8}=-2$. Per $n$ pari, invece, $x^{n}\ge0$ per ogni $x$ e i numeri negativi non
hanno radici $n$-esime reali: è precisamente questo il «difetto» di $\R$ che i numeri
complessi del \cref{sec:complessi} rimedieranno.
\end{osservazione}

\begin{osservazione}[Dalla radice $n$-esima a ogni polinomio]\label{oss:radice-polinomi}
Nella dimostrazione del \cref{teo:radice-n} la forma di $t^{n}$ è servita solo a
stimare di quanto cambia la potenza quando la base cambia di poco
(\cref{lem:stima-potenze}). Sommando stime analoghe per i monomi, si vede che ogni
polinomio $P$ conserva il segno vicino ai punti in cui non si annulla; e allora lo stesso
ragionamento sull'estremo superiore di $\{t\in[s,u] : P(t)<0\}$ prova che \emph{un
polinomio negativo in $s$ e positivo in $u>s$ si annulla in $(s,u)$}. È la strada di
\cite{giusti}, che ottiene le radici come caso particolare ($P(t)=t^{n}-y$); il
\cref{cap:funzioni-continue} estenderà il risultato a tutte le funzioni continue
(\cref{teo:permanenza-segno-continue,teo:zeri}).
\end{osservazione}

\subsection{Potenze a esponente razionale: monotonia}

Grazie al \cref{teo:radice-n} le potenze a esponente razionale della
\eqref{eq:potenza-razionale}, $a^{m/n}=\sqrt[n]{a^{m}}$ per $a>0$, sono ben definite. Ci
servirà la loro monotonia rispetto all'esponente, che dimostriamo.

\begin{lemma}[Monotonia delle potenze razionali]\label{lem:monotonia-razionale}
Sia $a>1$. Se $y,z\in\Q$ e $y<z$, allora $a^{y}<a^{z}$. Inoltre, per ogni $K>0$ esiste
$m\in\N$ tale che $a^{m}>K$ e $a^{-m}<1/K$.
\end{lemma}

\begin{schema}
\Ipotesi{$a>1$; \ $y,z\in\Q$, $y<z$; \ $K>0$.}
\Tesi{$a^{y}<a^{z}$; \ $\esiste m\in\N : a^{m}>K$ \ e \ $a^{-m}<1/K$.}
\end{schema}

\begin{strategia}
Riduciamo i due esponenti allo stesso denominatore, $y=p/q$ e $z=r/q$ con $p<r$
interi: tutto si riduce a confrontare due potenze \emph{intere} della stessa base
$b=a^{1/q}$, e $b>1$. Per la seconda parte serve una stima dal basso di $a^{m}$ che cresca
illimitatamente con $m$: la disuguaglianza di Bernoulli, scrivendo $a=1+(a-1)$, dà
$a^{m}\ge1+m(a-1)$, e la proprietà di Archimede fa il resto.
\end{strategia}

\begin{dimostrazione}
Scriviamo $y=p/q$ e $z=r/q$ con $p,r\in\Z$, $q\in\N$, $p<r$, e poniamo $b=a^{1/q}$. Si ha
$b>1$: se fosse $b\le1$, sarebbe $a=b^{q}\le1$. Per ogni $k\in\N$ è $b^{k}>1$ (per
induzione: se $b^{k}>1$, allora $b^{k+1}=b^{k}\,b>b^{k}>1$), dunque
$a^{z}=b^{r}=b^{p}\,b^{r-p}>b^{p}=a^{y}$, essendo $r-p\in\N$ e $b^{p}>0$.

Per la seconda parte, per la disuguaglianza di Bernoulli (\cref{teo:bernoulli}) con
$x=a-1>0$ si ha $a^{m}=(1+(a-1))^{m}\ge1+m(a-1)$; per la proprietà di Archimede esiste
$m>\frac{K-1}{a-1}$, e per tale $m$ risulta $a^{m}\ge1+m(a-1)>K$, cioè anche
$a^{-m}=1/a^{m}<1/K$.
\end{dimostrazione}

\begin{lemma}[di densità]\label{lem:densita-esponenziale}\index{lemma di densità}
Sia $a>1$. Per ogni coppia $\alpha,\beta$ di numeri reali con $0<\alpha<\beta$ esiste
$y\in\Q$ tale che
\[
  \alpha<a^{y}<\beta .
\]
In altre parole, l'immagine $\{a^{y} : y\in\Q\}$ della funzione esponenziale definita
sui razionali è densa in $(0,+\infty)$.
\end{lemma}

\begin{schema}
\Ipotesi{$a>1$; \ $0<\alpha<\beta$.}
\Tesi{$\esiste y\in\Q : \alpha<a^{y}<\beta$.}
\end{schema}

\begin{strategia}
Cerchiamo $y$ della forma $(m+1)/n$. Elevando alla $n$, la condizione
$\alpha<a^{(m+1)/n}<\beta$ diventa $\alpha^{n}<a^{m+1}<\beta^{n}$: dobbiamo trovare una
potenza intera di $a$ dentro l'intervallo $(\alpha^{n},\beta^{n})$. Le potenze intere
$a^{0},a^{1},a^{2},\dots$ crescono «a salti» di un fattore $a$; se l'intervallo è
abbastanza «largo in senso moltiplicativo», cioè se $\beta^{n}/\alpha^{n}>a$, una potenza
ci cade dentro per forza. Basta quindi scegliere $n$ grande, in modo che
$(\beta/\alpha)^{n}>a$, e prendere come $a^{m}$ l'ultima potenza che non supera
$\alpha^{n}$: la successiva, $a^{m+1}$, supera $\alpha^{n}$ ma non può superare
$\beta^{n}$. È la stessa idea della densità di $\Q$ (\cref{fig:densita-Q}), in versione
moltiplicativa. Per semplicità trattiamo prima il caso $\alpha\ge1$, poi ci riconduciamo
a esso.
\end{strategia}

\begin{dimostrazione}
\emph{Primo caso: $\alpha\ge1$.} Poiché $\beta/\alpha>1$, per il
\cref{lem:monotonia-razionale} (applicato con base $\beta/\alpha$) esiste $n\in\N$ tale
che $(\beta/\alpha)^{n}>a$, cioè $\beta^{n}>a\,\alpha^{n}$. Consideriamo l'insieme
\[
  T=\{m\in\N\cup\{0\} : a^{m}\le\alpha^{n}\}.
\]
$T$ è non vuoto, perché $a^{0}=1\le\alpha^{n}$ (essendo $\alpha\ge1$), ed è limitato
superiormente: per il \cref{lem:monotonia-razionale} esiste $\bar m$ con
$a^{\bar m}>\alpha^{n}$, e allora ogni $m\in T$ verifica $m<\bar m$. Per il
\cref{teo:max-N} (o banalmente, se $T=\{0\}$) $T$ ha massimo $m$. Allora $m+1\notin T$,
cioè $a^{m+1}>\alpha^{n}$; e d'altra parte
\[
  a^{m+1}=a\cdot a^{m}\le a\,\alpha^{n}<\beta^{n}.
\]
Dunque $\alpha^{n}<a^{m+1}<\beta^{n}$, e per la stretta monotonia della radice $n$-esima
(\cref{cor:radice-2}) $\alpha<a^{(m+1)/n}<\beta$: basta porre $y=(m+1)/n$.

\emph{Caso generale: $0<\alpha<1$.} Per il \cref{lem:monotonia-razionale} esiste
$k\in\N$ con $a^{k}>1/\alpha$, cioè $a^{k}\alpha>1$. Applicando il primo caso alla
coppia $a^{k}\alpha<a^{k}\beta$ si trova $y'\in\Q$ con $a^{k}\alpha<a^{y'}<a^{k}\beta$;
dividendo per $a^{k}$ si ottiene $\alpha<a^{y'-k}<\beta$, e $y=y'-k$ è razionale.
\end{dimostrazione}

\subsection{La funzione esponenziale su \texorpdfstring{$\R$}{R}}

Siamo pronti per definire $a^{x}$ con $x$ reale. Il lemma di densità suggerisce
l'idea: $a^{x}$ deve essere l'estremo superiore delle potenze $a^{y}$ con esponente
razionale $y<x$ (\cref{fig:esponenziale-sup}).

\begin{definizione}[Potenza a esponente reale]\label{def:potenza-reale}\index{funzione!esponenziale}\index{potenza a esponente reale}
Sia $a>1$ e $x\in\R$. Si pone
\begin{equation}\label{eq:esponenziale-sup}
  a^{x}=\sup\{a^{y} : y\in\Q,\ y<x\}.
\end{equation}
Se $0<a<1$ si pone $a^{x}=(1/a)^{-x}$, e se $a=1$ si pone $1^{x}=1$.
\end{definizione}

La definizione è ben posta: l'insieme $E_{x}=\{a^{y} : y\in\Q,\ y<x\}$ è non vuoto
(contiene $a^{\lfloor x\rfloor-1}$) e limitato superiormente (per il
\cref{lem:monotonia-razionale}, $a^{y}<a^{r}$ per ogni $y\in E_{x}$, dove $r$ è un
qualsiasi razionale maggiore di $x$, ad esempio $r=\lfloor x\rfloor+1$). Inoltre
$a^{x}\ge a^{\lfloor x\rfloor-1}>0$.

\begin{teorema}[La funzione esponenziale]\label{teo:esponenziale}\index{funzione!esponenziale!biiettività}
Sia $a>1$. Allora:
\begin{enumerate}
  \item se $x\in\Q$, la \eqref{eq:esponenziale-sup} restituisce il valore $a^{x}$ già
  noto (\emph{coerenza});
  \item la funzione $x\mapsto a^{x}$ è strettamente crescente su $\R$;
  \item per ogni $c>0$ esiste uno e un solo $x\in\R$ tale che $a^{x}=c$.
\end{enumerate}
Quindi la funzione esponenziale è una biiezione strettamente crescente di $\R$ su
$(0,+\infty)$. Per $0<a<1$ essa è una biiezione strettamente decrescente di $\R$ su
$(0,+\infty)$.
\end{teorema}

\begin{schema}
\Ipotesi{$a>1$; potenze razionali con le proprietà dei \cref{lem:monotonia-razionale,lem:densita-esponenziale}.}
\Tesi{(i) coerenza su $\Q$; \ (ii) $x_{1}<x_{2}\implica a^{x_{1}}<a^{x_{2}}$; \
(iii) $\perogni c>0\ \esiste!\,x\in\R : a^{x}=c$.}
\end{schema}

\begin{strategia}
Le tre affermazioni si dimostrano con due soli strumenti, usati più volte: la
caratterizzazione dell'estremo superiore (\cref{prop:caratt-sup}) e la possibilità di
«infilare» un razionale ovunque, sia sull'asse degli esponenti (densità di $\Q$,
\cref{teo:densita-Q}) sia sull'asse dei valori (\cref{lem:densita-esponenziale}).
\begin{itemize}
  \item Per (i) mostriamo che il vecchio valore $a^{x}$ è il minimo maggiorante di
  $E_{x}$: è un maggiorante per la monotonia razionale, ed è il \emph{minimo} perché sotto
  di esso, per il lemma di densità, cade sempre qualche $a^{y}$ con $y<x$.
  \item Per (ii), dati $x_{1}<x_{2}$, infiliamo due razionali $r_{1}<r_{2}$ tra di essi:
  $a^{r_{1}}$ è un maggiorante di $E_{x_{1}}$ e $a^{r_{2}}$ è un elemento di $E_{x_{2}}$.
  \item Per (iii) imitiamo la dimostrazione del \cref{teo:radice-n}: definiamo $x$ come
  estremo superiore degli esponenti razionali $y$ con $a^{y}<c$, ed escludiamo per assurdo
  sia $a^{x}<c$ sia $a^{x}>c$, usando questa volta il lemma di densità al posto delle
  stime esplicite.
\end{itemize}
\end{strategia}

\begin{dimostrazione}
(i) Sia $x\in\Q$ e indichiamo provvisoriamente con $P$ il valore della potenza razionale
$a^{x}$. Per il \cref{lem:monotonia-razionale}, $a^{y}<P$ per ogni razionale $y<x$:
$P$ è un maggiorante di $E_{x}$. Sia $\epsilon>0$; se $\epsilon\ge P$ ogni elemento di
$E_{x}$, essendo positivo, supera $P-\epsilon\le0$; supponiamo quindi $0<\epsilon<P$. Per il
\cref{lem:densita-esponenziale} esiste $y\in\Q$ con $P-\epsilon<a^{y}<P$, e
necessariamente $y<x$ (se fosse $y\ge x$, sarebbe $a^{y}\ge P$). Dunque $a^{y}\in E_{x}$ e
$a^{y}>P-\epsilon$; per la \cref{prop:caratt-sup}, $\sup E_{x}=P$.

(ii) Siano $x_{1}<x_{2}$. Applicando due volte il \cref{teo:densita-Q} si trovano due
razionali $r_{1},r_{2}$ con $x_{1}<r_{1}<r_{2}<x_{2}$. Se $y\in\Q$ e $y<x_{1}$, allora
$y<r_{1}$ e $a^{y}<a^{r_{1}}$: dunque $a^{r_{1}}$ è un maggiorante di $E_{x_{1}}$, e
$a^{x_{1}}\le a^{r_{1}}$. D'altra parte $r_{2}<x_{2}$, quindi $a^{r_{2}}\in E_{x_{2}}$ e
$a^{r_{2}}\le a^{x_{2}}$. In conclusione
\[
  a^{x_{1}}\le a^{r_{1}}<a^{r_{2}}\le a^{x_{2}} .
\]

(iii) Sia $c>0$ e consideriamo $S=\{y\in\Q : a^{y}<c\}$. Per il
\cref{lem:monotonia-razionale} esiste $m\in\N$ con $a^{-m}<c$, quindi $-m\in S$ e $S$ è
non vuoto; esiste anche $m'\in\N$ con $a^{m'}>c$, e allora ogni $y\in S$ verifica $y<m'$
(se fosse $y\ge m'$ sarebbe $a^{y}\ge a^{m'}>c$): $S$ è limitato superiormente. Sia
$x=\sup S$.
\begin{itemize}
  \item Se fosse $a^{x}<c$, per il \cref{lem:densita-esponenziale} esisterebbe
  $r\in\Q$ con $a^{x}<a^{r}<c$. Allora $r\in S$, e quindi $r\le x$; ma da $a^{r}>a^{x}$ e
  dalla (ii) segue $r>x$. Assurdo.
  \item Se fosse $a^{x}>c$, esisterebbe $r\in\Q$ con $c<a^{r}<a^{x}$, e dalla (ii)
  $r<x$. Poiché $x=\sup S$ e $r<x$, per la \cref{prop:caratt-sup} (con $\epsilon=x-r$)
  esiste $y\in S$ con $y>r$; allora $a^{y}>a^{r}>c$, contro il fatto che $y\in S$.
  Assurdo.
\end{itemize}
Dunque $a^{x}=c$. L'unicità segue dalla stretta monotonia (ii). Infine, il caso
$0<a<1$ segue dalla definizione $a^{x}=(1/a)^{-x}$, essendo $1/a>1$: la composizione di
$x\mapsto-x$ (decrescente) con un esponenziale crescente è decrescente.
\end{dimostrazione}

\begin{figure}[htb]
\centering
\begin{tikzpicture}
\begin{axis}[width=0.92\textwidth, height=6cm, xmin=0.9, xmax=1.62, ymin=1.9, ymax=3.1,
  xtick={1,1.2,1.4,1.4142,1.6}, xticklabels={$1$,$1{,}2$,\makebox[0pt][r]{$1{,}4$\,},\makebox[0pt][l]{\,$\sqrt2$},$1{,}6$},
  ytick={2,2.5,2.6651}, yticklabels={$2$,$2{,}5$,$2^{\sqrt2}$}, axis lines=left]
  \addplot[curva, domain=0.9:1.6, samples=100] {2^x};
  \addplot[tratteggio] coordinates {(0.9,2.66514) (1.41421,2.66514) (1.41421,1.9)};
  \foreach \y in {1,1.2,1.3,1.4,1.41,1.414} {
    \addplot[only marks, mark=*, mark size=1.5pt, rame] coordinates {(\y,{2^\y})};}
  \addplot[only marks, mark=o, mark size=2.2pt, verde, thick] coordinates {(1.41421,2.66514)};
  \node[font=\scriptsize, rame, anchor=west, align=left] at (axis cs:1.0,2.82)
    {punti $(y,2^{y})$ con $y\in\Q$, $y<\sqrt2$:\\ $y=1;\ 1{,}2;\ 1{,}3;\ 1{,}4;\ 1{,}41;\ 1{,}414$};
\end{axis}
\end{tikzpicture}
\caption{La definizione \eqref{eq:esponenziale-sup}: $2^{\sqrt2}$ è l'estremo superiore
dei valori $2^{y}$ con $y$ razionale minore di $\sqrt2$. I punti color rame,
calcolati con esponenti razionali, salgono verso $2^{\sqrt2}\approx2{,}665$ senza mai
superarlo.}
\label{fig:esponenziale-sup}
\end{figure}

\begin{osservazione}[Che cosa abbiamo guadagnato]
Il \cref{teo:esponenziale} chiude un conto aperto nel \cref{sec:potenze}: la funzione
logaritmo, definita come inversa dell'esponenziale dalla \eqref{eq:log-def}, è ora definita su
\emph{tutto} $(0,+\infty)$, perché ogni $c>0$ è della forma $a^{x}$. Si noti che per
ottenere la suriettività non abbiamo usato il teorema dell'esistenza degli zeri, che
richiede la continuità, ma direttamente l'estremo superiore. Le proprietà algebriche
$a^{x}a^{z}=a^{x+z}$, $(a^{x})^{z}=a^{xz}$, $(ab)^{x}=a^{x}b^{x}$, note per esponenti
razionali, si estendono agli esponenti reali con ragionamenti dello stesso tipo
(approssimando $x$ e $z$ con razionali); le dimostreremo in modo più semplice quando
avremo a disposizione la continuità della funzione esponenziale. Si noti anche che,
per $a>1$, vale pure $a^{x}=\inf\{a^{q} : q\in\Q,\ q>x\}$: per la (ii) ogni $a^{q}$ con
$q>x$ è maggiore di $a^{x}$, e per ogni $\epsilon>0$ il \cref{lem:densita-esponenziale}
fornisce un razionale $r$ con $a^{x}<a^{r}<a^{x}+\epsilon$, e necessariamente $r>x$. Dunque
$a^{x}$ è l'\emph{unico} numero che separa le potenze con esponente razionale minore di
$x$ da quelle con esponente razionale maggiore di $x$, come anticipato nel
\cref{sec:potenze}. Per $0<a<1$ si può verificare (esercizio) che
$a^{x}=\sup\{a^{y} : y\in\Q,\ y>x\}$, che è la definizione adottata in
\cite{marcellini-sbordone}.
\end{osservazione}

\begin{notastorica}[Esponenti frazionari, negativi, irrazionali]
La notazione esponenziale moderna $a^{3}$, $a^{4}$ è di Cartesio (\emph{La Géométrie},
1637), che però per il quadrato scriveva ancora di solito $aa$. L'idea di una potenza con
esponente frazionario è molto più antica: nel Trecento Nicole Oresme, nel \emph{De
proportionibus proportionum}, ne faceva già uso, sia pure senza la notazione moderna. In
forma sistematica, gli esponenti negativi e frazionari furono introdotti da John Wallis
(\emph{Arithmetica infinitorum}, 1656) e resi d'uso comune da Newton, che in una lettera
del 1676 destinata a Leibniz (la \emph{Epistola prior}, inviata tramite Oldenburg)
scriveva tranquillamente $a^{1/2}$, $a^{-1}$, $a^{3/4}$. Gli esponenti irrazionali
comparivano già nel Settecento (Eulero, nell'\emph{Introductio} del 1748, osserva che
$a^{\sqrt7}$ è compreso tra $a^{2}$ e $a^{3}$), ma senza una definizione rigorosa, che
poteva venire solo con una teoria dei numeri reali. Fino ad allora anche l'esistenza della
radice $n$-esima (\cref{teo:radice-n}) era trattata come una proprietà evidente del
sistema numerico, da usare senza dimostrazione \cite{giusti}; e Dedekind, presentando nel
1872 le sue sezioni, notava con una punta di orgoglio che solo ora si potevano dare vere
dimostrazioni di uguaglianze come $\sqrt2\cdot\sqrt3=\sqrt6$, che a quanto gli risultava
nessuno aveva mai dimostrato. Una curiosità
che mostra quanto la questione sia profonda: nel 1900 Hilbert pose, come settimo dei suoi
celebri problemi, la domanda se $2^{\sqrt2}$ fosse trascendente (cioè non radice di
alcun polinomio a coefficienti interi). La risposta, affermativa, fu data da Rodion
Kuzmin nel 1930; il caso generale del problema fu risolto nel 1934--35 dal teorema di
Gelfond--Schneider.
\end{notastorica}

\begin{python}[Approssimare $2^{\sqrt2}$ con esponenti razionali]
# Approssimiamo √2 per difetto con i troncamenti decimali y_k
# (numeri razionali) e calcoliamo 2**y_k: per la definizione di
# potenza reale come sup, i valori crescono e si avvicinano
# a 2**√2 senza superarlo.
from fractions import Fraction as F
import math

r2 = math.sqrt(2)
for k in range(1, 8):
    y = F(math.floor(r2 * 10**k), 10**k)   # troncamento < √2
    print(f"y = {str(y):>15s} = {float(y):.7f}"
          f"   2^y = {2 ** float(y):.10f}")
print(f"{'2^√2':>37s} = {2 ** r2:.10f}")
\end{python}
\begin{risultato}
y =             7/5 = 1.4000000   2^y = 2.6390158215
y =         141/100 = 1.4100000   2^y = 2.6573716282
y =         707/500 = 1.4140000   2^y = 2.6647496502
y =       7071/5000 = 1.4142000   2^y = 2.6651190885
y =   141421/100000 = 1.4142100   2^y = 2.6651375618
y = 1414213/1000000 = 1.4142130   2^y = 2.6651431038
y = 2828427/2000000 = 1.4142135   2^y = 2.6651440275
                                 2^√2 = 2.6651441427
\end{risultato}

%% ----------------------------------------------------------------------------
%%  §2.4  Calcolo combinatorio
%% ----------------------------------------------------------------------------
\section{Calcolo combinatorio}
\label{sec:combinatoria}

Il calcolo combinatorio è l'arte di \emph{contare senza elencare}: stabilire quanti sono
gli oggetti di un certo tipo senza doverli scrivere uno per uno. Sia $A$ un insieme
costituito da $n$ elementi:
\begin{equation}\label{eq:insieme-n}
  A=\{a_{1},a_{2},\dots,a_{n}\}.
\end{equation}

\subsection{Disposizioni e permutazioni}

\begin{definizione}[Disposizione]\index{disposizione}
Sia $k$ un numero naturale minore o uguale a $n$. Una \emph{disposizione} (semplice) di
$k$ elementi tra gli $n$ dati è una $k$-upla ordinata $(b_{1},b_{2},\dots,b_{k})$ di
elementi \emph{distinti} di $A$. Due disposizioni sono considerate distinte se
differiscono per gli elementi, oppure solo per l'ordine in cui gli elementi compaiono.
\end{definizione}

Nel linguaggio del \cref{sec:invertibili}, una disposizione di $k$ elementi di $A$ non è
altro che una funzione \emph{iniettiva} $\{1,2,\dots,k\}\to A$: il primo posto va in
$b_{1}$, il secondo in $b_{2}$, e così via, e l'iniettività esprime il fatto che gli
elementi sono distinti.

\needspace{8\baselineskip}
\begin{teorema}[Numero delle disposizioni]\label{teo:disposizioni}\index{disposizione!numero delle}
Il numero $D_{n,k}$ delle disposizioni di $k$ elementi tra $n$ dati ($1\le k\le n$) è
\begin{equation}\label{eq:disposizioni}
  D_{n,k}=n(n-1)(n-2)\cdots(n-k+1),
\end{equation}
cioè il prodotto di $k$ numeri interi decrescenti a partire da $n$.
\end{teorema}

\begin{schema}
\Ipotesi{$A$ insieme di $n$ elementi; $k\in\N$, $1\le k\le n$.}
\Tesi{Le $k$-uple ordinate di elementi distinti di $A$ sono $n(n-1)\cdots(n-k+1)$.}
\end{schema}

\begin{strategia}
Esaminiamo un possibile modo di formare una disposizione: possiamo scegliere il primo
elemento tra gli $n$ dati; il secondo solo tra gli $n-1$ elementi rimasti dopo la prima
scelta; il terzo tra gli $n-2$ rimasti, e così via. Ogni scelta moltiplica il numero di
possibilità per il numero di elementi ancora disponibili (\cref{fig:albero-disposizioni}).
Per rendere rigoroso il «così via» facciamo induzione su $k$: ogni disposizione di $k+1$
elementi si ottiene, in uno e un solo modo, da una disposizione di $k$ elementi (i primi
$k$ posti) aggiungendo in coda uno degli $n-k$ elementi non ancora usati. Quindi
$D_{n,k+1}=D_{n,k}\cdot(n-k)$: è il cosiddetto \emph{principio di moltiplicazione}.
\end{strategia}

\begin{dimostrazione}
Procediamo per induzione su $k$, per $1\le k\le n$. Per $k=1$ le disposizioni sono le
$1$-uple $(b_{1})$ con $b_{1}\in A$: sono $n$, in accordo con la \eqref{eq:disposizioni}.

Supponiamo la formula vera per un certo $k<n$. Ogni disposizione
$(b_{1},\dots,b_{k},b_{k+1})$ di $k+1$ elementi individua la disposizione
$(b_{1},\dots,b_{k})$ dei suoi primi $k$ elementi e l'elemento $b_{k+1}$, che è uno degli
$n-k$ elementi di $A$ diversi da $b_{1},\dots,b_{k}$. Viceversa, da ogni disposizione di
$k$ elementi e da ogni scelta di uno degli $n-k$ elementi rimasti si ottiene una
disposizione di $k+1$ elementi, e disposizioni (o scelte) diverse danno disposizioni
diverse. Le disposizioni di $k+1$ elementi si ripartiscono quindi in $D_{n,k}$ gruppi
(uno per ogni possibile «inizio» di lunghezza $k$), ciascuno di $n-k$ elementi:
\[
  D_{n,k+1}=D_{n,k}\,(n-k)=n(n-1)\cdots(n-k+1)\,(n-k),
\]
che è la \eqref{eq:disposizioni} con $k+1$ al posto di $k$.
\end{dimostrazione}

Così ad esempio il numero delle disposizioni di $2$ elementi su $3$ dati è
$3\cdot2=6$. Infatti, se consideriamo l'insieme $\{a_{1},a_{2},a_{3}\}$, tutte le
disposizioni possibili con $2$ elementi sono:
\begin{equation}\label{eq:disposizioni-esempio}
  (a_{1},a_{2}),\quad (a_{1},a_{3}),\quad (a_{2},a_{1}),\quad (a_{2},a_{3}),\quad
  (a_{3},a_{1}),\quad (a_{3},a_{2}).
\end{equation}

\begin{figure}[htb]
\centering
\begin{tikzpicture}[
  grow=right, level distance=2.4cm,
  level 1/.style={sibling distance=1.9cm},
  level 2/.style={sibling distance=0.85cm},
  every node/.style={font=\footnotesize},
  nodo/.style={circle, draw=blu, fill=bluchiaro, inner sep=1.5pt, minimum size=6mm},
  edge from parent/.style={draw=grigio, -{Latex[length=1.6mm]}}]
  \node[nodo, fill=blu, text=white] {inizio}
    child {node[nodo] {$a_{3}$}
      child {node[nodo] {$a_{2}$} node[right=6mm, rame] {$(a_{3},a_{2})$}}
      child {node[nodo] {$a_{1}$} node[right=6mm, rame] {$(a_{3},a_{1})$}}}
    child {node[nodo] {$a_{2}$}
      child {node[nodo] {$a_{3}$} node[right=6mm, rame] {$(a_{2},a_{3})$}}
      child {node[nodo] {$a_{1}$} node[right=6mm, rame] {$(a_{2},a_{1})$}}}
    child {node[nodo] {$a_{1}$}
      child {node[nodo] {$a_{3}$} node[right=6mm, rame] {$(a_{1},a_{3})$}}
      child {node[nodo] {$a_{2}$} node[right=6mm, rame] {$(a_{1},a_{2})$}}};
  \node[grigio, align=center] at (2.4,-3.4) {1\textsuperscript{a} scelta:\\ $3$ possibilità};
  \node[grigio, align=center] at (4.8,-3.4) {2\textsuperscript{a} scelta:\\ $2$ possibilità};
  \node[rame, align=center] at (7.4,-3.4) {$3\cdot2=6$\\ disposizioni};
\end{tikzpicture}
\caption{L'albero delle scelte per le disposizioni di $2$ elementi tra $3$: ogni
livello moltiplica il numero dei rami per il numero di elementi ancora disponibili.}
\label{fig:albero-disposizioni}
\end{figure}

\begin{definizione}[Permutazioni e fattoriale]\index{permutazione}\index{fattoriale}
Una disposizione di $n$ elementi tra gli $n$ dati ($k=n$) si chiama \emph{permutazione}
degli $n$ elementi. Ponendo $k=n$ nel \cref{teo:disposizioni} si ottiene che il numero
delle permutazioni di $n$ elementi è
\begin{equation}\label{eq:fattoriale}
  n!=n\cdot(n-1)\cdot(n-2)\cdots2\cdot1;
\end{equation}
il simbolo $n!$, che indica il prodotto dei primi $n$ numeri naturali, si legge
«$n$ fattoriale». Si pone inoltre, per convenzione, $0!=1$.
\end{definizione}

Una permutazione è dunque una funzione biunivoca $\{1,\dots,n\}\to A$: ogni elemento va
in un posto, e ogni posto è occupato; componendola con la corrispondenza $a_{j}\mapsto j$
la si può vedere come una funzione biunivoca di $A$ in sé. Tenendo conto del simbolo di fattoriale, si
riconosce che il numero di disposizioni di $k$ elementi tra $n$ dati si può anche
scrivere
\begin{equation}\label{eq:disposizioni-fattoriale}
  D_{n,k}=n(n-1)\cdots(n-k+1)=\frac{n!}{(n-k)!}\,,
\end{equation}
perché nella frazione i fattori da $n-k$ fino a $1$ si semplificano. La convenzione
$0!=1$ rende la formula valida anche per $k=n$.

\begin{esempio}[Estensimetri su un'ala]\label{es:estensimetri}
Su un longherone alare sono predisposti $8$ punti di misura, e si dispone di $3$
estensimetri collegati a tre canali \emph{diversi} della centralina di acquisizione. In
quanti modi si può allestire la prova? Una configurazione è individuata dicendo quale
punto è collegato al canale 1, quale al canale 2, quale al canale 3: è una disposizione
di $3$ elementi tra $8$, e le configurazioni sono $D_{8,3}=8\cdot7\cdot6=336$. Se invece
si chiede in quanti \emph{ordini} diversi si possono eseguire $5$ ispezioni distinte su
un velivolo, la risposta è il numero di permutazioni di $5$ elementi: $5!=120$.
\end{esempio}

\begin{osservazione}[Disposizioni con ripetizione]\index{disposizione!con ripetizione}
Se si ammette che gli elementi della $k$-upla possano ripetersi, ogni posto può essere
occupato da uno qualsiasi degli $n$ elementi, indipendentemente dagli altri; lo stesso
ragionamento per induzione dà $n^{k}$ \emph{disposizioni con ripetizione}. Ad esempio le
stringhe di $k$ bit sono $2^{k}$: un byte ($k=8$) può assumere $256$ valori. Associando a
ogni sottoinsieme $X$ di $A$ la stringa di $n$ bit che ha $1$ al posto $j$ se $a_{j}\in X$
e $0$ altrimenti, si ritrova che un insieme di $n$ elementi ha $2^{n}$ sottoinsiemi
(esercizio del \cref{sec:induzione}).
\end{osservazione}

\begin{attenzione}[Il fattoriale cresce in fretta]
Il fattoriale cresce più rapidamente di qualsiasi potenza $a^{n}$: $10!=3\,628\,800$,
$20!\approx2{,}4\cdot10^{18}$, e $52!$ --- il numero dei modi di mescolare un mazzo di
carte francesi --- è circa $8\cdot10^{67}$. Questo spiega perché i problemi che
richiedono di esaminare tutte le permutazioni (come quello del commesso viaggiatore, che
cerca il percorso più breve tra $n$ città) diventano impraticabili anche per il più
potente dei calcolatori già per qualche decina di elementi.
\end{attenzione}

\subsection{Combinazioni e coefficienti binomiali}

\begin{definizione}[Combinazione]\index{combinazione}
Una \emph{combinazione} di $k$ elementi tra gli $n$ dati ($0\le k\le n$) è un
sottoinsieme di $A$ con $k$ elementi; consideriamo uguali due combinazioni che hanno gli
stessi elementi, indipendentemente dall'ordine.
\end{definizione}

\begin{teorema}[Numero delle combinazioni]\label{teo:combinazioni}\index{combinazione!numero delle}\index{coefficiente binomiale}
Il numero delle combinazioni di $k$ elementi tra $n$ dati è
\begin{equation}\label{eq:combinazioni}
  C_{n,k}=\binom{n}{k}=\frac{n!}{k!\,(n-k)!}\,.
\end{equation}
\end{teorema}

\begin{schema}
\Ipotesi{$A$ insieme di $n$ elementi; $0\le k\le n$.}
\Tesi{I sottoinsiemi di $A$ con $k$ elementi sono $\dfrac{n!}{k!\,(n-k)!}$.}
\end{schema}

\begin{strategia}
Non contiamo le combinazioni direttamente, ma le mettiamo in relazione con le
disposizioni, che sappiamo già contare. Una disposizione è una combinazione «con
l'ordine»: da ogni combinazione di $k$ elementi, ordinandone gli elementi in tutti i modi
possibili, si ottengono $k!$ disposizioni diverse (le permutazioni dei suoi $k$
elementi). E ogni disposizione proviene da una sola combinazione: quella formata dai suoi
elementi. Dunque le disposizioni sono esattamente $k!$ volte le combinazioni, e basta
dividere.
\end{strategia}

\begin{dimostrazione}
Per $k=0$ c'è un solo sottoinsieme con $0$ elementi, l'insieme vuoto, e infatti
$\frac{n!}{0!\,n!}=1$. Sia $1\le k\le n$. A ogni
disposizione $(b_{1},\dots,b_{k})$ associamo la combinazione $\{b_{1},\dots,b_{k}\}$ dei
suoi elementi. In questo modo l'insieme delle $D_{n,k}$ disposizioni si ripartisce in
gruppi, uno per ogni combinazione; il gruppo corrispondente alla combinazione
$\{c_{1},\dots,c_{k}\}$ è formato dalle $k$-uple ordinate dei suoi elementi, cioè dalle
loro permutazioni, che sono $k!$. Quindi
\[
  D_{n,k}=C_{n,k}\cdot k!, \qquad\text{cioè}\qquad
  C_{n,k}=\frac{D_{n,k}}{k!}=\frac{n!}{k!\,(n-k)!},
\]
dove nell'ultimo passaggio abbiamo usato la \eqref{eq:disposizioni-fattoriale}.
\end{dimostrazione}

Il simbolo $\binom nk$ si legge «$n$ su $k$» ed è chiamato \emph{coefficiente binomiale}
(il motivo del nome sarà chiaro nel prossimo paragrafo). Grazie alla convenzione $0!=1$
esso ha senso anche per $k=0$ e $k=n$:
\begin{equation}\label{eq:binomiale-estremi}
  \binom n0=\binom nn=\frac{n!}{n!\,0!}=1,
\end{equation}
in accordo con il fatto che c'è un solo sottoinsieme con $0$ elementi (l'insieme vuoto)
e un solo sottoinsieme con $n$ elementi ($A$ stesso).

Ritornando all'esempio di un insieme con $3$ elementi $\{a_{1},a_{2},a_{3}\}$, si hanno
$3!/(1!\,2!)=3$ combinazioni con due elementi, che sono
\begin{equation}\label{eq:combinazioni-esempio}
  \{a_{1},a_{2}\},\qquad \{a_{1},a_{3}\},\qquad \{a_{2},a_{3}\};
\end{equation}
ciascuna corrisponde a $2!=2$ delle disposizioni \eqref{eq:disposizioni-esempio}.

\begin{proposizione}[Simmetria dei coefficienti binomiali]\label{prop:simmetria-binomiale}
Per ogni $n\in\N$ e ogni $k\in\{0,1,\dots,n\}$,
\begin{equation}\label{eq:simmetria-binomiale}
  \binom nk=\binom{n}{n-k}.
\end{equation}
\end{proposizione}

\begin{schema}
\Ipotesi{$n\in\N$, \ $0\le k\le n$.}
\Tesi{$\binom nk=\binom{n}{n-k}$.}
\end{schema}

\begin{strategia}
Diamo due argomenti. Il primo è un calcolo: nella formula \eqref{eq:combinazioni} il
denominatore $k!\,(n-k)!$ è simmetrico rispetto allo scambio di $k$ con $n-k$. Il
secondo è un conteggio: si costruisce una corrispondenza biunivoca tra i due insiemi di
sottoinsiemi contati dai due membri, associando a ogni sottoinsieme il suo complementare.
\end{strategia}

\begin{dimostrazione}
\emph{Algebricamente}, la \eqref{eq:simmetria-binomiale} segue direttamente dalla
definizione \eqref{eq:combinazioni}: scambiando $k$ con $n-k$ il denominatore
$k!\,(n-k)!$ non cambia. \emph{Combinatoriamente}, scegliere i $k$ elementi da prendere
equivale a scegliere gli $n-k$ elementi da lasciare: la funzione $X\mapsto A\setminus X$
è una corrispondenza biunivoca tra i sottoinsiemi con $k$ elementi e quelli con $n-k$
elementi.
\end{dimostrazione}

\begin{idea}[Due dimostrazioni, due modi di pensare]
La proposizione precedente illustra un fenomeno frequente in combinatoria: un'identità
tra numeri si può dimostrare \emph{calcolando} (manipolando le formule) oppure
\emph{contando} (mostrando che i due membri contano gli stessi oggetti in due modi
diversi). La seconda strada, quando è disponibile, è di solito più illuminante: spiega
\emph{perché} l'identità è vera. Ritroveremo questa doppia possibilità nella formula del
binomio di Newton.
\end{idea}

\begin{esempio}
Riprendiamo l'\cref{es:estensimetri}, ma con $3$ estensimetri \emph{identici} collegati
in serie a un unico canale: conta solo \emph{quali} punti sono strumentati, non l'ordine,
e le configurazioni sono $\binom83=\frac{8\cdot7\cdot6}{3!}=56$, cioè $336/6$. Un esempio
meno ingegneristico: le sestine possibili al SuperEnalotto, in cui si estraggono $6$
numeri su $90$, sono $\binom{90}{6}=622\,614\,630$; giocando una sestina la probabilità
di fare «sei» è circa una su seicento milioni.
\end{esempio}

\begin{python}[Contare ed elencare con \texttt{itertools} e \texttt{math}]
import itertools, math

A = ["a1", "a2", "a3"]
disp = list(itertools.permutations(A, 2))   # disposizioni 2 su 3
comb = list(itertools.combinations(A, 2))   # combinazioni 2 su 3
print(len(disp), *(f"({x},{y})" for x, y in disp))
print(len(comb), *(f"{{{x},{y}}}" for x, y in comb))

# Le formule chiuse sono già pronte nella libreria standard
print(math.perm(8, 3), math.comb(8, 3), math.factorial(5))
print(math.comb(90, 6))                     # SuperEnalotto

# Verifica di D(n,k) = C(n,k) * k! per tutti i casi con n <= 10
print(all(math.perm(n, k) == math.comb(n, k) * math.factorial(k)
          for n in range(1, 11) for k in range(0, n + 1)))
\end{python}
\begin{risultato}
6 (a1,a2) (a1,a3) (a2,a1) (a2,a3) (a3,a1) (a3,a2)
3 {a1,a2} {a1,a3} {a2,a3}
336 56 120
622614630
True
\end{risultato}

%% ----------------------------------------------------------------------------
%%  §2.5  Il binomio di Newton
%% ----------------------------------------------------------------------------
\section{Il binomio di Newton}
\label{sec:binomio}

Un'importante applicazione dei risultati del paragrafo precedente (che peraltro è
anche conseguenza del principio di induzione, come mostreremo nella seconda parte del
paragrafo) è la formula che esprime la potenza $n$-esima di un binomio. Per scriverla in
modo compatto conviene introdurre un simbolo che useremo in tutto il resto del libro.

\begin{notazione}[Il simbolo di sommatoria]\index{sommatoria}
Dati i numeri reali $c_{p},c_{p+1},\dots,c_{q}$ (con $p\le q$ interi), la loro somma si
indica con
\[
  \sum_{k=p}^{q}c_{k}=c_{p}+c_{p+1}+\dots+c_{q},
\]
che si legge «somma per $k$ che va da $p$ a $q$ di $c_{k}$». La lettera $k$ è un
\emph{indice muto}: $\sum_{k=p}^{q}c_{k}$ e $\sum_{j=p}^{q}c_{j}$ indicano la stessa
somma. Valgono le proprietà (che seguono dalle proprietà associativa, commutativa e
distributiva)
\[
  \sum_{k=p}^{q}(c_{k}+d_{k})=\sum_{k=p}^{q}c_{k}+\sum_{k=p}^{q}d_{k}, \qquad
  \sum_{k=p}^{q}h\,c_{k}=h\sum_{k=p}^{q}c_{k}\quad(h\in\R),
\]
e la regola di \emph{traslazione dell'indice}: ponendo $j=k+1$,
\begin{equation}\label{eq:traslazione-indice}
  \sum_{k=p}^{q}c_{k}=\sum_{j=p+1}^{q+1}c_{j-1}.
\end{equation}
Con questa notazione, ad esempio, la \eqref{eq:gauss} si scrive
$\sum_{k=1}^{n}k=\frac{n(n+1)}{2}$.
\end{notazione}

\begin{teorema}[Formula del binomio di Newton]\label{teo:binomio}\index{binomio di Newton}
Per ogni coppia di numeri reali $a,b$ e per ogni $n\in\N$ vale l'identità
\begin{equation}\label{eq:binomio}
\begin{split}
  (a+b)^{n}&=\binom n0a^{n}+\binom n1a^{n-1}b+\binom n2a^{n-2}b^{2}+\dots\\
  &\qquad\dots+\binom nk a^{n-k}b^{k}+\dots+\binom{n}{n-1}ab^{n-1}+\binom nn b^{n}\\
  &=\sum_{k=0}^{n}\binom nk a^{n-k}b^{k}.
\end{split}
\end{equation}
\end{teorema}

\begin{schema}
\Ipotesi{$a,b\in\R$; $n\in\N$ (con la convenzione $c^{0}=1$ per ogni $c$, anche
$c=0$, nei termini estremi della somma).}
\Tesi{$(a+b)^{n}=\displaystyle\sum_{k=0}^{n}\binom nk a^{n-k}b^{k}$.}
\end{schema}

Diamo due dimostrazioni, che corrispondono ai due modi di pensare descritti nel
paragrafo precedente: \emph{contare} e \emph{calcolare}.

\begin{strategia}[Piano d'attacco della prima dimostrazione]
Scriviamo $(a+b)^{n}$ come prodotto di $n$ fattori uguali e immaginiamo di svolgerlo
con la proprietà distributiva senza raccogliere nulla: otteniamo una somma di $2^{n}$
prodotti, ciascuno dei quali si forma scegliendo da ogni parentesi o la $a$ o la $b$
(\cref{fig:binomio-scelte}). Un prodotto in cui la $b$ è stata scelta $k$ volte vale
$a^{n-k}b^{k}$. Quanti sono? Tanti quanti i modi di scegliere \emph{quali} $k$ parentesi
forniscono la $b$, cioè i sottoinsiemi di $k$ elementi di un insieme di $n$ parentesi:
$\binom nk$. Il coefficiente binomiale deve il suo nome proprio a questo.
\end{strategia}

\begin{dimostrazione}[Prima dimostrazione (combinatoria)]
Immaginiamo di eseguire il prodotto di $n$ fattori
\begin{equation}\label{eq:binomio-prodotto}
  (a+b)^{n}=(a+b)\cdot(a+b)\cdots(a+b) \qquad (n\ \text{fattori});
\end{equation}
applicando ripetutamente la proprietà distributiva, il risultato si scrive come somma
algebrica di addendi, ciascuno dei quali si ottiene scegliendo un termine ($a$ oppure
$b$) da ognuno degli $n$ fattori e moltiplicando i termini scelti. Fissiamo la nostra
attenzione sugli addendi del tipo $a^{n-k}b^{k}$: eseguendo la moltiplicazione indicata,
si ottiene tale addendo tante volte quante è possibile scegliere i $k$ fattori dai quali
prendere la $b$ tra gli $n$ dati (dagli altri $n-k$ si prende necessariamente la $a$).
Per il \cref{teo:combinazioni}, ciò avviene $\binom nk$ volte. Raccogliendo gli addendi
uguali si ottiene la \eqref{eq:binomio}.
\end{dimostrazione}

\begin{marginfigure}
\centering
\small
\renewcommand{\arraystretch}{1.15}
\begin{tabular}{@{}c@{\quad}l@{\quad}c@{}}
\toprule
n.\ di $b$ & prodotti & tot.\\
\midrule
$0$ & $aaa$ & $1$\\
$1$ & $aab,\ aba,\ baa$ & $3$\\
$2$ & $abb,\ bab,\ bba$ & $3$\\
$3$ & $bbb$ & $1$\\
\bottomrule
\end{tabular}
\caption{Gli $2^{3}=8$ prodotti che si ottengono svolgendo $(a+b)(a+b)(a+b)$, raggruppati
secondo il numero di $b$ scelte: $(a+b)^{3}=a^{3}+3a^{2}b+3ab^{2}+b^{3}$.}
\label{fig:binomio-scelte}
\end{marginfigure}

La formula del binomio \eqref{eq:binomio} fornisce in particolare risultati di facile
verifica per $n=2$ e $n=3$: infatti, per tali valori di $n$ si ottengono le ben note
relazioni
\begin{gather}
  \binom20=1,\quad \binom21=2,\quad \binom22=1, \notag\\
  \text{da cui}\quad (a+b)^{2}=a^{2}+2ab+b^{2};\\[3pt]
  \binom30=\binom33=1,\quad \binom31=\binom32=3, \notag\\
  \text{da cui}\quad (a+b)^{3}=a^{3}+3a^{2}b+3ab^{2}+b^{3}.
\end{gather}

\subsection{La formula di Stifel e il triangolo di Tartaglia}

Le considerazioni che seguono sono utili per sviluppare esplicitamente il secondo
membro della formula del binomio per valori di $n$ più grandi di $3$, e forniscono
l'ingrediente della seconda dimostrazione.

\begin{lemma}[Formula di Stifel]\label{lem:stifel}\index{formula di Stifel}
Per ogni $n\in\N$ e ogni $k\in\N$ con $1\le k\le n$ vale la formula
\begin{equation}\label{eq:stifel}
  \binom{n}{k-1}+\binom nk=\binom{n+1}{k}.
\end{equation}
\end{lemma}

\begin{schema}
\Ipotesi{$n\in\N$, \ $1\le k\le n$.}
\Tesi{$\binom{n}{k-1}+\binom nk=\binom{n+1}{k}$.}
\end{schema}

\begin{strategia}
È un calcolo con i fattoriali: si mettono le due frazioni a denominatore comune
raccogliendo i fattori comuni $n!$, $(k-1)!$ e $(n-k)!$, e si usano le relazioni
$k!=k\cdot(k-1)!$ e $(n-k+1)!=(n-k+1)\cdot(n-k)!$. Dopo il calcolo daremo anche
un'interpretazione combinatoria, che spiega l'identità senza conti.
\end{strategia}

\begin{dimostrazione}
Dalla definizione \eqref{eq:combinazioni} otteniamo
\begin{equation}\label{eq:stifel-calcolo}
\begin{split}
  \binom{n}{k-1}+\binom nk
  &=\frac{n!}{(k-1)!\,(n-k+1)!}+\frac{n!}{k!\,(n-k)!}\\
  &=\frac{n!}{(k-1)!\,(n-k)!}\Bigl(\frac{1}{n-k+1}+\frac1k\Bigr)\\
  &=\frac{n!}{(k-1)!\,(n-k)!}\cdot\frac{n+1}{k\,(n-k+1)}\\
  &=\frac{(n+1)!}{k!\,(n+1-k)!}=\binom{n+1}{k}.
\end{split}
\end{equation}
Il calcolo è lecito anche per $k=n$, grazie alla convenzione $0!=1$.
\end{dimostrazione}

\begin{osservazione}[Interpretazione combinatoria della formula di Stifel]
Consideriamo un insieme di $n+1$ elementi $\{a_{1},\dots,a_{n},a_{n+1}\}$ e contiamo i
suoi sottoinsiemi di $k$ elementi, che sono $\binom{n+1}{k}$, dividendoli in due gruppi:
quelli che \emph{non} contengono $a_{n+1}$, che sono i sottoinsiemi di $k$ elementi di
$\{a_{1},\dots,a_{n}\}$, cioè $\binom nk$; e quelli che \emph{contengono} $a_{n+1}$, che
si ottengono aggiungendo $a_{n+1}$ a un sottoinsieme di $k-1$ elementi di
$\{a_{1},\dots,a_{n}\}$, cioè $\binom{n}{k-1}$. La somma è la \eqref{eq:stifel}.
\end{osservazione}

La proprietà \eqref{eq:stifel} permette di scrivere il seguente \emph{triangolo di
Tartaglia} (noto anche come \emph{triangolo di Pascal}), nel quale ogni coefficiente è
uguale alla somma dei due coefficienti più vicini della riga precedente
(\cref{fig:tartaglia}).

\begin{figure}[htb]
\centering
\begin{tikzpicture}[x=0.82cm, y=0.62cm, every node/.style={font=\footnotesize\addfontfeatures{Numbers=Lining}},
  cella/.style={circle, inner sep=1.2pt, minimum size=5.4mm, fill=bluchiaro, draw=blu!40},
  celrame/.style={cella, fill=rame!30, draw=rame},
  celverde/.style={cella, fill=verde!30, draw=verde}]
  % (valori calcolati una volta per tutte: pgfmath non gestisce i fattoriali grandi)
  \node[grigio, anchor=east] at (-5.2,-0) {$n=0$};
  \node[cella] (t-0-0) at (0,-0) {1};
  \node[grigio, anchor=east] at (-5.2,-1) {$n=1$};
  \node[cella] (t-1-0) at (-0.5,-1) {1};
  \node[cella] (t-1-1) at (0.5,-1) {1};
  \node[grigio, anchor=east] at (-5.2,-2) {$n=2$};
  \node[cella] (t-2-0) at (-1,-2) {1};
  \node[cella] (t-2-1) at (0,-2) {2};
  \node[cella] (t-2-2) at (1,-2) {1};
  \node[grigio, anchor=east] at (-5.2,-3) {$n=3$};
  \node[cella] (t-3-0) at (-1.5,-3) {1};
  \node[cella] (t-3-1) at (-0.5,-3) {3};
  \node[cella] (t-3-2) at (0.5,-3) {3};
  \node[cella] (t-3-3) at (1.5,-3) {1};
  \node[grigio, anchor=east] at (-5.2,-4) {$n=4$};
  \node[cella] (t-4-0) at (-2,-4) {1};
  \node[cella] (t-4-1) at (-1,-4) {4};
  \node[cella] (t-4-2) at (0,-4) {6};
  \node[cella] (t-4-3) at (1,-4) {4};
  \node[cella] (t-4-4) at (2,-4) {1};
  \node[grigio, anchor=east] at (-5.2,-5) {$n=5$};
  \node[cella] (t-5-0) at (-2.5,-5) {1};
  \node[cella] (t-5-1) at (-1.5,-5) {5};
  \node[celrame] (t-5-2) at (-0.5,-5) {10};
  \node[celrame] (t-5-3) at (0.5,-5) {10};
  \node[cella] (t-5-4) at (1.5,-5) {5};
  \node[cella] (t-5-5) at (2.5,-5) {1};
  \node[grigio, anchor=east] at (-5.2,-6) {$n=6$};
  \node[cella] (t-6-0) at (-3,-6) {1};
  \node[cella] (t-6-1) at (-2,-6) {6};
  \node[cella] (t-6-2) at (-1,-6) {15};
  \node[celverde] (t-6-3) at (0,-6) {20};
  \node[cella] (t-6-4) at (1,-6) {15};
  \node[cella] (t-6-5) at (2,-6) {6};
  \node[cella] (t-6-6) at (3,-6) {1};
  \node[grigio, anchor=east] at (-5.2,-7) {$n=7$};
  \node[cella] (t-7-0) at (-3.5,-7) {1};
  \node[cella] (t-7-1) at (-2.5,-7) {7};
  \node[celrame] (t-7-2) at (-1.5,-7) {21};
  \node[celrame] (t-7-3) at (-0.5,-7) {35};
  \node[cella] (t-7-4) at (0.5,-7) {35};
  \node[cella] (t-7-5) at (1.5,-7) {21};
  \node[cella] (t-7-6) at (2.5,-7) {7};
  \node[cella] (t-7-7) at (3.5,-7) {1};
  \node[grigio, anchor=east] at (-5.2,-8) {$n=8$};
  \node[cella] (t-8-0) at (-4,-8) {1};
  \node[cella] (t-8-1) at (-3,-8) {8};
  \node[cella] (t-8-2) at (-2,-8) {28};
  \node[celverde] (t-8-3) at (-1,-8) {56};
  \node[cella] (t-8-4) at (0,-8) {70};
  \node[cella] (t-8-5) at (1,-8) {56};
  \node[cella] (t-8-6) at (2,-8) {28};
  \node[cella] (t-8-7) at (3,-8) {8};
  \node[cella] (t-8-8) at (4,-8) {1};
  \draw[-{Latex[length=1.6mm]}, rame, thick] (t-5-2) -- (t-6-3);
  \draw[-{Latex[length=1.6mm]}, rame, thick] (t-5-3) -- (t-6-3);
  \draw[-{Latex[length=1.6mm]}, rame, thick] (t-7-2) -- (t-8-3);
  \draw[-{Latex[length=1.6mm]}, rame, thick] (t-7-3) -- (t-8-3);
  \node[rame, anchor=west] at (5.0,-5.5) {$\binom52+\binom53=\binom63$};
  \node[rame, anchor=west] at (5.0,-7.5) {$\binom72+\binom73=\binom83$};
\end{tikzpicture}
\caption{Il triangolo di Tartaglia: nella riga $n$ compaiono i coefficienti
$\binom n0,\binom n1,\dots,\binom nn$. Ogni numero è la somma dei due che gli stanno
sopra (formula di Stifel): ad esempio $10+10=20$ e $21+35=56$. La simmetria di ogni riga
è la \eqref{eq:simmetria-binomiale}.}
\label{fig:tartaglia}
\end{figure}

Quindi, ad esempio, leggendo la riga $n=5$ risulta
\begin{equation}\label{eq:binomio-5}
  (a+b)^{5}=a^{5}+5a^{4}b+10a^{3}b^{2}+10a^{2}b^{3}+5ab^{4}+b^{5}.
\end{equation}

\subsection{La dimostrazione per induzione}

Di seguito esponiamo un'ulteriore dimostrazione della formula del binomio, con il metodo
di induzione.

\needspace{8\baselineskip}
\begin{strategia}[Piano d'attacco della seconda dimostrazione]
Supposta vera la formula per $n$, otteniamo $(a+b)^{n+1}$ moltiplicando lo sviluppo di
$(a+b)^{n}$ per $(a+b)$. Si ottengono due somme: in quella che proviene dalla
moltiplicazione per $a$ l'esponente di $a$ sale di uno, in quella che proviene dalla
moltiplicazione per $b$ sale l'esponente di $b$. Per poterle sommare termine a termine
occorre che i termini con la stessa potenza $a^{n+1-k}b^{k}$ abbiano lo stesso indice:
lo si ottiene traslando l'indice della seconda somma \eqref{eq:traslazione-indice}. A
quel punto ogni potenza $a^{n+1-k}b^{k}$ ha coefficiente $\binom nk+\binom n{k-1}$, che
per la formula di Stifel è proprio $\binom{n+1}k$. È, in formule, la costruzione del
triangolo di Tartaglia: la riga $n+1$ si ottiene dalla riga $n$.
\end{strategia}

\begin{dimostrazione}[Seconda dimostrazione (per induzione)]
Cominciamo con l'osservare che la formula del binomio è verificata per $n=1$; infatti
\begin{equation}\label{eq:binomio-base}
  (a+b)^{1}=\binom10a+\binom11b=a+b.
\end{equation}
Supponiamo vera la \eqref{eq:binomio} per un indice $n\in\N$ e dimostriamo la validità
della formula analoga con $n+1$ al posto di $n$. A tal fine moltiplichiamo entrambi i
membri della \eqref{eq:binomio} per $(a+b)$, ottenendo
\[
  (a+b)^{n+1}=\sum_{k=0}^{n}\binom nk a^{n+1-k}b^{k}+\sum_{k=0}^{n}\binom nk a^{n-k}b^{k+1}.
\]
Nella seconda somma poniamo $j=k+1$ (\eqref{eq:traslazione-indice}); poiché
$a^{n-k}b^{k+1}=a^{n+1-j}b^{j}$, essa diventa
$\sum_{j=1}^{n+1}\binom{n}{j-1}a^{n+1-j}b^{j}$. Rinominando $j$ in $k$ (l'indice è muto)
e isolando il primo termine della prima somma ($k=0$) e l'ultimo della seconda
($k=n+1$):
\begin{equation}\label{eq:binomio-passo}
\begin{split}
  (a+b)^{n+1}&=a^{n+1}+\sum_{k=1}^{n}\Bigl[\binom nk+\binom n{k-1}\Bigr]a^{n+1-k}b^{k}+b^{n+1}\\
  &=\binom{n+1}{0}a^{n+1}+\sum_{k=1}^{n}\binom{n+1}{k}a^{n+1-k}b^{k}+\binom{n+1}{n+1}b^{n+1}\\
  &=\sum_{k=0}^{n+1}\binom{n+1}{k}a^{n+1-k}b^{k},
\end{split}
\end{equation}
dove abbiamo usato la formula di Stifel \eqref{eq:stifel} per $1\le k\le n$ e la
\eqref{eq:binomiale-estremi}. È la \eqref{eq:binomio} con $n+1$ al posto di $n$; per il
principio di induzione la formula è dimostrata.
\end{dimostrazione}

\subsection{Conseguenze e applicazioni}

\begin{corollario}\label{cor:binomio}
Per ogni $n\in\N$:
\begin{enumerate}
  \item $\displaystyle\sum_{k=0}^{n}\binom nk=2^{n}$;
  \item $\displaystyle\sum_{k=0}^{n}(-1)^{k}\binom nk=0$;
  \item se $x\ge0$, \ $(1+x)^{n}\ge1+nx+\dfrac{n(n-1)}{2}\,x^{2}$.
\end{enumerate}
\end{corollario}

\begin{schema}
\Ipotesi{$n\in\N$; \ per la (iii), $x\ge0$.}
\Tesi{(i) $\sum_{k=0}^{n}\binom nk=2^{n}$; \ (ii) $\sum_{k=0}^{n}(-1)^{k}\binom nk=0$; \
(iii) $(1+x)^{n}\ge1+nx+\frac{n(n-1)}{2}x^{2}$.}
\end{schema}

\begin{strategia}
La formula del binomio è un'identità in $a$ e $b$: scegliendo valori particolari si
ottengono informazioni sui coefficienti. Con $a=b=1$ le potenze scompaiono e restano i
soli coefficienti; con $b=-1$ compaiono i segni alterni. Per la (iii) si sfrutta il
fatto che, se $b=x\ge0$, tutti gli addendi sono non negativi, e una somma di termini non
negativi non è minore di una parte di essi.
\end{strategia}

\begin{dimostrazione}
(i) e (ii) si ottengono dalla \eqref{eq:binomio} con $a=b=1$ e con $a=1$, $b=-1$. Per
(iii), con $a=1$ e $b=x\ge0$ tutti gli addendi della \eqref{eq:binomio} sono non
negativi; trascurando quelli con $k\ge3$ resta
$(1+x)^{n}\ge1+\binom n1x+\binom n2x^{2}=1+nx+\frac{n(n-1)}2x^{2}$ (per $n=1$ il
termine in $x^{2}$ è nullo e la disuguaglianza è un'uguaglianza).
\end{dimostrazione}

\begin{osservazione}
La (i) ha un significato combinatorio: $\binom nk$ conta i sottoinsiemi con $k$
elementi di un insieme di $n$ elementi, e sommando su tutti i $k$ si contano tutti i
sottoinsiemi, che sono $2^{n}$ (esercizio del \cref{sec:induzione}). La (iii) è un
rafforzamento della disuguaglianza di Bernoulli (\cref{teo:bernoulli}) per $x\ge0$, e
ci servirà nello studio delle successioni: per esempio, per dimostrare che
$\sqrt[n]{n}$ si avvicina a $1$ al crescere di $n$.
\end{osservazione}

\begin{esempio}[Sensibilità della portanza e della potenza alla velocità]\label{es:sensibilita}
La portanza di un'ala è $L=\frac12\rho V^{2}SC_{L}$. Se la velocità aumenta del $2\%$,
cioè passa da $V$ a $V(1+\epsilon)$ con $\epsilon=0{,}02$, a parità di $C_{L}$ la
portanza diventa
\[
  L(1+\epsilon)^{2}=L\,(1+2\epsilon+\epsilon^{2})=L\,(1+0{,}04+0{,}0004),
\]
cioè aumenta del $4{,}04\%$. La potenza necessaria a vincere la resistenza è
$P=DV=\frac12\rho V^{3}SC_{D}$; un aumento del $3\%$ della velocità ($\epsilon=0{,}03$)
la moltiplica per
\[
  (1+\epsilon)^{3}=1+3\epsilon+3\epsilon^{2}+\epsilon^{3}=1+0{,}09+0{,}0027+0{,}000027
  =1{,}092727,
\]
cioè più del $9\%$. La formula del binomio giustifica la regola pratica, usata
continuamente nelle stime ingegneristiche, $(1+\epsilon)^{n}\approx1+n\epsilon$ per
$\epsilon$ «piccolo»: l'errore commesso è costituito dai termini con $\epsilon^{2}$,
$\epsilon^{3}$, \dots, che per $\epsilon=0{,}03$ valgono complessivamente meno di
$0{,}003$. Una variazione percentuale piccola su una grandezza elevata alla $n$ produce,
in prima approssimazione, una variazione percentuale $n$ volte maggiore.
\end{esempio}

\begin{notastorica}[Tartaglia, Pascal, Newton --- e prima di loro]
Il triangolo porta in Italia il nome di Niccolò Fontana detto Tartaglia (Brescia,
1499 circa -- Venezia, 1557), che lo pubblicò nel \emph{General trattato di numeri et misure}
(1556), e in Francia quello di Blaise Pascal, che ne fece uno studio sistematico nel
\emph{Traité du triangle arithmétique} (scritto nel 1654, pubblicato postumo nel 1665),
usandolo per risolvere problemi di gioco
d'azzardo posti dal cavaliere de Méré: è l'atto di nascita del calcolo delle
probabilità. Ma il triangolo era noto da secoli: in Cina compare nel trattato di Yang Hui
(1261), che lo attribuisce a Jia Xian (XI secolo); in Persia era conosciuto da al-Karaji
e da Omar Khayyam; in India, in forma combinatoria, già nell'opera di Pingala sulla
metrica dei versi. La formula di Stifel prende il nome da Michael Stifel
(\emph{Arithmetica integra}, 1544). Il merito di Isaac Newton (1665) fu di estendere la
formula agli esponenti \emph{non interi}, ottenendo uno sviluppo in serie infinita di
$(1+x)^{\alpha}$ che ritroveremo studiando la formula di Taylor.
\end{notastorica}

\begin{python}[Il triangolo di Tartaglia costruito con la formula di Stifel]
import sympy as sp

def tartaglia(righe):
    riga = [1]
    for n in range(righe):
        yield riga
        # formula di Stifel: ogni elemento interno è la somma
        # dei due elementi che gli stanno sopra
        riga = ([1] + [riga[k - 1] + riga[k]
                       for k in range(1, len(riga))] + [1])

for riga in tartaglia(8):
    print(" ".join(f"{c:3d}" for c in riga).center(40))

a, b = sp.symbols("a b")
print(sp.expand((a + b)**5))
\end{python}
\begin{risultato}
                    1                   
                  1   1                 
                1   2   1               
              1   3   3   1             
            1   4   6   4   1           
          1   5  10  10   5   1         
        1   6  15  20  15   6   1       
      1   7  21  35  35  21   7   1     
a**5 + 5*a**4*b + 10*a**3*b**2 + 10*a**2*b**3 + 5*a*b**4 + b**5
\end{risultato}

\begin{figure}[htb]
\centering
\includegraphics[width=0.78\textwidth]{sierpinski}
\caption{Le prime $128$ righe del triangolo di Tartaglia, colorando i coefficienti
dispari e lasciando in bianco i pari (\texttt{python/fig\_sierpinski.py}). Compare il
\emph{triangolo di Sierpiński}, un frattale: la struttura si ripete identica a scale
diverse. Il motivo è proprio la formula di Stifel, letta «modulo $2$»: pari$+$pari e
dispari$+$dispari danno pari, pari$+$dispari dà dispari.}
\label{fig:sierpinski}
\end{figure}

%% ----------------------------------------------------------------------------
%%  §2.6  Appendice: i numeri complessi
%% ----------------------------------------------------------------------------
\section{Appendice: i numeri complessi}
\label{sec:complessi}

\subsection{Perché ampliare \texorpdfstring{$\R$}{R}}

Consideriamo una generica equazione di secondo grado nell'incognita $z$:
\begin{equation}\label{eq:secondo-grado}
  az^{2}+bz+c=0 \qquad (a\neq0),
\end{equation}
che ha come soluzioni (reali, se $b^{2}-4ac\ge0$)
\begin{equation}\label{eq:formula-risolutiva}
  z_{1,2}=-\frac{b}{2a}\pm\frac{\sqrt{b^{2}-4ac}}{2a}.
\end{equation}
I coefficienti $a$, $b$, $c$ sono legati alle soluzioni $z_{1}$, $z_{2}$ anche dalle
semplici relazioni, che si verificano immediatamente sommando e moltiplicando le
\eqref{eq:formula-risolutiva},
\begin{equation}\label{eq:viete}
  z_{1}+z_{2}=-\frac ba, \qquad z_{1}\cdot z_{2}=\frac ca .
\end{equation}
Proviamo a scrivere tali relazioni nel caso dell'equazione
\begin{equation}\label{eq:z2piu1}
  z^{2}+1=0;
\end{equation}
in tal caso le \eqref{eq:viete} diventano $z_{1}+z_{2}=0$ e $z_{1}\cdot z_{2}=1$.
Queste relazioni, anche se formalmente ben definite, non hanno alcun senso nell'ambito
dei numeri reali, perché non esistono $z_{1}$, $z_{2}$ soluzioni \emph{reali} della
\eqref{eq:z2piu1}: il quadrato di un numero reale non è mai negativo. Questo esempio
mostra come talvolta sia utile pensare che qualsiasi equazione di secondo grado ammetta
soluzione. Dato che ciò non vale nell'ambito dei numeri reali, si amplia $\R$
introducendo il campo $\C$ dei \emph{numeri complessi}.

\begin{notastorica}[Cardano, Bombelli e il «più di meno»]
Curiosamente, i numeri complessi non nacquero dalle equazioni di secondo grado (di fronte
a $z^{2}+1=0$ ci si limitava a dire che non ci sono soluzioni) ma da quelle di
\emph{terzo} grado. Nell'\emph{Ars Magna} (1545) Gerolamo Cardano pubblicò la formula
risolutiva dell'equazione $x^{3}=px+q$, scoperta da Scipione del Ferro e da Tartaglia.
Applicata a $x^{3}=15x+4$, che ha la soluzione reale $x=4$, la formula produce
\[
  x=\sqrt[3]{2+\sqrt{-121}}+\sqrt[3]{2-\sqrt{-121}},
\]
un'espressione con radici di numeri negativi. Fu il bolognese Rafael Bombelli, nella sua
\emph{L'Algebra} (1572), ad avere il coraggio di calcolare con questi oggetti
«sofistici», introducendo regole come «più di meno via più di meno fa meno» (cioè
$i\cdot i=-1$). Bombelli osservò che $(2+\sqrt{-1})^{3}=2+\sqrt{-121}$ e
$(2-\sqrt{-1})^{3}=2-\sqrt{-121}$, e quindi
$x=(2+\sqrt{-1})+(2-\sqrt{-1})=4$: i numeri «impossibili» servivano a trovare una
soluzione perfettamente reale. Il lettore potrà verificare il calcolo di Bombelli dopo
aver letto questo paragrafo.
\end{notastorica}

\subsection{Definizione e forma algebrica}

Il testo di riferimento \cite{marcellini-sbordone} introduce i numeri complessi
definendo «$i=\sqrt{-1}$» e calcolando «con le stesse regole dei reali». È un approccio
efficace, ma lascia aperta una domanda: \emph{esiste} davvero un sistema di numeri che
contiene $\R$ e un elemento con quadrato $-1$, nel quale valgano tutte le regole del
calcolo? La risposta più semplice, dovuta a William Rowan Hamilton (1837), consiste nel
costruirlo esplicitamente come insieme di coppie di numeri reali.

\needspace{8\baselineskip}
\begin{definizione}[Numeri complessi]\index{numeri complessi}
Si chiama insieme dei \emph{numeri complessi}, e si indica con $\C$, l'insieme
$\R^{2}=\R\times\R$ delle coppie ordinate di numeri reali, munito delle operazioni
\begin{align}
  (x,y)+(x',y')&=(x+x',\ y+y'), \label{eq:somma-complessi}\\
  (x,y)\cdot(x',y')&=(xx'-yy',\ xy'+x'y). \label{eq:prodotto-complessi}
\end{align}
\end{definizione}

La somma è quella «componente per componente»; il prodotto, a prima vista misterioso, è
scelto apposta perché funzioni la costruzione che segue.

\begin{teorema}[$\C$ è un campo]\label{teo:C-campo}\index{campo!dei numeri complessi}
Con le operazioni \eqref{eq:somma-complessi}, \eqref{eq:prodotto-complessi} l'insieme
$\C$ soddisfa gli assiomi (A1)--(A6) del \cref{sec:assiomi}, con elementi neutri
$(0,0)$ e $(1,0)$, opposto $-(x,y)=(-x,-y)$ e, per $(x,y)\neq(0,0)$, inverso
\begin{equation}\label{eq:inverso-complesso}
  (x,y)^{-1}=\Bigl(\frac{x}{x^{2}+y^{2}},\ \frac{-y}{x^{2}+y^{2}}\Bigr).
\end{equation}
\end{teorema}

\begin{schema}
\Ipotesi{$\C=\R^{2}$ con le operazioni \eqref{eq:somma-complessi},
\eqref{eq:prodotto-complessi}; \ gli assiomi di campo valgono in $\R$.}
\Tesi{$\C$ soddisfa (A1)--(A6), con zero $(0,0)$, unità $(1,0)$, opposto $(-x,-y)$ e,
per $(x,y)\neq(0,0)$, inverso dato dalla \eqref{eq:inverso-complesso}.}
\end{schema}

\begin{strategia}
Si tratta di verificare uno per uno gli assiomi, usando le proprietà dei numeri reali
sulle componenti. Le proprietà della somma sono immediate (si lavora componente per
componente). Per il prodotto, la commutatività si vede dalla simmetria della formula;
l'associatività e la distributiva richiedono un calcolo diretto, che svolgiamo; per
l'inverso basta verificare che la coppia proposta funziona, dopo aver osservato che il
denominatore $x^{2}+y^{2}$ non si annulla.
\end{strategia}

\begin{dimostrazione}
Le proprietà associativa e commutativa della somma, l'esistenza dello zero $(0,0)$ e
dell'opposto $(-x,-y)$ seguono immediatamente dalle corrispondenti proprietà dei numeri
reali, applicate alle due componenti. Il prodotto è commutativo, perché scambiando
$(x,y)$ con $(x',y')$ la \eqref{eq:prodotto-complessi} non cambia. È associativo: posto
$z=(a,b)$, $w=(c,d)$, $u=(e,f)$,
\begin{align*}
  (zw)u&=(ac-bd,\ ad+bc)\cdot(e,f)\\
       &=(ace-bde-adf-bcf,\ acf-bdf+ade+bce),\\
  z(wu)&=(a,b)\cdot(ce-df,\ cf+de)\\
       &=(ace-adf-bcf-bde,\ acf+ade+bce-bdf),
\end{align*}
e le due coppie coincidono. È distributivo:
\begin{align*}
  (a,b)\cdot\bigl[(c,d)+(e,f)\bigr]&=\bigl(a(c+e)-b(d+f),\ a(d+f)+b(c+e)\bigr)\\
  &=(ac-bd,\ ad+bc)+(ae-bf,\ af+be).
\end{align*}
L'elemento $(1,0)$ è neutro: $(x,y)(1,0)=(x\cdot1-y\cdot0,\ x\cdot0+1\cdot y)=(x,y)$; ed
è diverso da $(0,0)$. Infine, se $(x,y)\neq(0,0)$, allora $x^{2}+y^{2}>0$ (una somma di
quadrati di numeri reali è nulla solo se entrambi sono nulli, perché ogni quadrato è
$\ge0$), e
\[
  (x,y)\cdot\Bigl(\frac{x}{x^{2}+y^{2}},\frac{-y}{x^{2}+y^{2}}\Bigr)
  =\Bigl(\frac{x^{2}+y^{2}}{x^{2}+y^{2}},\ \frac{-xy+xy}{x^{2}+y^{2}}\Bigr)=(1,0). \qedhere
\]
\end{dimostrazione}

Poiché valgono gli assiomi (A1)--(A6), in $\C$ valgono anche \emph{tutte} le loro
conseguenze dimostrate nel \cref{sec:conseguenze}: leggi di semplificazione, legge di
annullamento del prodotto, regola dei segni, unicità di opposto e inverso. Ecco il
vantaggio dell'approccio assiomatico: non serve ridimostrare nulla.

\paragraph{I reali dentro i complessi.} Le coppie della forma $(x,0)$ si sommano e si
moltiplicano esattamente come i numeri reali:
\[
  (x,0)+(x',0)=(x+x',0), \qquad (x,0)\cdot(x',0)=(xx',0).
\]
Possiamo quindi \emph{identificare} il numero reale $x$ con la coppia $(x,0)$, e
considerare $\R$ come un sottoinsieme di $\C$. Poniamo poi
\[
  i=(0,1),
\]
e risulta
\begin{equation}\label{eq:unita-immaginaria}
  i^{2}=(0,1)\cdot(0,1)=(0\cdot0-1\cdot1,\ 0\cdot1+1\cdot0)=(-1,0)=-1.
\end{equation}
Il numero $i$ si chiama \emph{unità immaginaria}\index{unità immaginaria}. Infine, per
ogni coppia,
\[
  (x,y)=(x,0)+(0,y)=(x,0)+(y,0)\cdot(0,1)=x+iy .
\]
Abbiamo così ottenuto la \emph{forma algebrica} dei numeri complessi:
\begin{equation}\label{eq:forma-algebrica}
  \C=\{z=x+iy : x,y\in\R\}.
\end{equation}
Si dice che $x=\operatorname{Re}z$ è la \emph{parte reale} e $y=\operatorname{Im}z$ è
il \emph{coefficiente della parte immaginaria} di $z$. Due numeri complessi sono uguali
se e solo se hanno uguali parte reale e coefficiente dell'immaginario (è la
\eqref{eq:coppie}). Le operazioni sui numeri complessi si eseguono con le stesse regole
dei numeri reali, tenendo anche conto della \eqref{eq:unita-immaginaria}; quindi:
\begin{align}
  (x+iy)+(x'+iy')&=(x+x')+i(y+y'); \label{eq:somma-algebrica}\\
  (x+iy)\cdot(x'+iy')&=xx'+ixy'+ix'y+i^{2}yy' \notag\\
                      &=(xx'-yy')+i(xy'+x'y). \label{eq:prodotto-algebrico}
\end{align}
La \eqref{eq:prodotto-algebrico} spiega, a posteriori, la scelta del prodotto
\eqref{eq:prodotto-complessi}: è l'unica compatibile con la proprietà distributiva e con
$i^{2}=-1$.

\begin{notastorica}[Hamilton: dalle coppie ai quaternioni]
William Rowan Hamilton (1805--1865) pubblicò la sua teoria delle «coppie algebriche» nel
1837, negli atti dell'Accademia irlandese: voleva fondare i numeri complessi sulla sola
algebra, senza bisogno del piano di Wessel e di Argand di cui diremo più avanti. Il
successo lo spinse a cercare per lo \emph{spazio} ciò che le coppie fanno per il piano:
una moltiplicazione di terne $(x,y,z)$ che rappresentasse le rotazioni tridimensionali,
come quella dei complessi rappresenta le rotazioni piane. Non poteva riuscirci (oggi
sappiamo che in dimensione $3$ una moltiplicazione con le proprietà richieste non esiste),
ma il 16 ottobre 1843, camminando lungo il Royal Canal di Dublino, capì che bastava
passare alle \emph{quaterne}, con tre unità immaginarie legate dalle relazioni
$i^{2}=j^{2}=k^{2}=ijk=-1$, che incise sulla pietra del Brougham Bridge. Il prezzo era la
rinuncia alla commutatività: $ij=k$, ma $ji=-k$. I \emph{quaternioni}, a lungo considerati
una curiosità, sono oggi di uso quotidiano in ingegneria aerospaziale: i sistemi di
navigazione inerziale e di controllo d'assetto di velivoli e veicoli spaziali descrivono
l'orientamento con un quaternione di modulo $1$, che non soffre delle singolarità degli
angoli di Eulero.
\end{notastorica}

\begin{proposizione}[$\C$ non è ordinabile]\label{prop:C-non-ordinabile}
Non esiste in $\C$ alcuna relazione $\le$ che soddisfi gli assiomi dell'ordinamento
(O1)--(O4).
\end{proposizione}

\begin{schema}
\Ipotesi{$\C$ è un campo nel quale $i^{2}=-1$.}
\Tesi{Non esiste in $\C$ una relazione $\le$ che verifichi (O1)--(O4).}
\end{schema}

\begin{strategia}
Supponiamo per assurdo che un tale ordinamento esista. Allora valgono in $\C$ tutte le
conseguenze degli assiomi del \cref{sec:conseguenze}, e in particolare un fatto che si
deduce dai soli assiomi (lo verifichiamo nella dimostrazione): \emph{ogni quadrato è $\ge0$}.
Ma in $\C$ anche $-1$ è un quadrato, $-1=i^{2}$; e $1=1^{2}$ è un quadrato. Avremmo
dunque $1\ge0$ e $-1\ge0$, da cui $1=0$.
\end{strategia}

\begin{dimostrazione}
Se esistesse un tale ordinamento, per ogni $z\in\C$ si avrebbe $z^{2}\ge0$: se $z\ge0$
per l'assioma \ref{ass:ordine-prodotto}; se $z\le0$, allora $-z\ge0$
(\cref{prop:segno-opposto}) e $z^{2}=(-z)^{2}\ge0$ (\cref{prop:segni}). In particolare
$1=1^{2}\ge0$ e $-1=i^{2}\ge0$; dalla seconda, per la \cref{prop:segno-opposto},
$1\le0$, e per l'antisimmetria \ref{ass:antisimmetrica} $1=0$, contro l'assioma
\ref{ass:neutri}.
\end{dimostrazione}

\begin{attenzione}[Nei complessi non si scrivono disuguaglianze]
Una conseguenza pratica: scrivere $z<w$ con $z,w$ complessi non reali non ha senso. Le
disuguaglianze si scrivono solo tra numeri reali, ad esempio tra i \emph{moduli}
$\abs z$, $\abs w$, che definiamo tra poco.
\end{attenzione}

\subsection{Coniugato e modulo}

\begin{definizione}[Coniugato e modulo]\index{complesso coniugato}\index{modulo!di un numero complesso}
Il numero complesso $\bar z=x-iy$ si chiama \emph{complesso coniugato} del numero
$z=x+iy$. Il numero reale non negativo
\begin{equation}\label{eq:modulo-complesso}
  \abs z=\sqrt{x^{2}+y^{2}}
\end{equation}
si chiama \emph{modulo} di $z$. Se $z=x$ è reale, $\abs z=\sqrt{x^{2}}$ coincide con il
valore assoluto di $x$: la notazione è coerente.
\end{definizione}

\begin{proposizione}[Proprietà del coniugato e del modulo]\label{prop:coniugato-modulo}
Per ogni $z,w\in\C$:
\begin{enumerate}
  \item $\overline{z+w}=\bar z+\bar w$, \quad $\overline{zw}=\bar z\,\bar w$, \quad
  $\bar{\bar z}=z$;
  \item $z\,\bar z=\abs z^{2}$; \quad $\operatorname{Re}z=\dfrac{z+\bar z}{2}$, \quad
  $\operatorname{Im}z=\dfrac{z-\bar z}{2i}$;
  \item $\abs{zw}=\abs z\,\abs w$; \quad $\abs{\bar z}=\abs z$; \quad
  $\abs{\operatorname{Re}z}\le\abs z$, \quad $\abs{\operatorname{Im}z}\le\abs z$;
  \item $\abs{z+w}\le\abs z+\abs w$ \quad (disuguaglianza triangolare).
\end{enumerate}
\end{proposizione}

\begin{schema}
\Ipotesi{$z=x+iy$, \ $w=u+iv$, \ con $x,y,u,v\in\R$.}
\Tesi{Le identità (i)--(ii), le relazioni (iii) sui moduli e la disuguaglianza
triangolare $\abs{z+w}\le\abs z+\abs w$.}
\end{schema}

\begin{strategia}
(i) e (ii) sono calcoli diretti con la forma algebrica. Il punto (ii),
$z\bar z=\abs z^{2}$, è lo strumento principale per il resto: permette di trasformare
affermazioni sui moduli (che contengono radici) in affermazioni sui \emph{quadrati} dei
moduli, che sono prodotti e si calcolano con l'algebra. Così per (iii) calcoleremo
$\abs{zw}^{2}$ e per la disuguaglianza triangolare $\abs{z+w}^{2}$; alla fine si
torna ai moduli estraendo la radice, che è crescente (\cref{cor:radice-2}).
\end{strategia}

\begin{dimostrazione}
(i) Se $z=x+iy$, $w=u+iv$, si ha $zw=(xu-yv)+i(xv+yu)$ e
$\bar z\bar w=(x-iy)(u-iv)=(xu-yv)-i(xv+yu)=\overline{zw}$; le altre due sono immediate.
(ii) $z\bar z=(x+iy)(x-iy)=x^{2}-i^{2}y^{2}=x^{2}+y^{2}=\abs z^{2}$; inoltre
$z+\bar z=2x$ e $z-\bar z=2iy$.
(iii) Per (i) e (ii), $\abs{zw}^{2}=zw\,\overline{zw}=z\bar z\,w\bar w=\abs z^{2}\abs w^{2}$,
e poiché i moduli sono non negativi, $\abs{zw}=\abs z\abs w$. La seconda è ovvia; per la
terza, $\abs{\operatorname{Re}z}=\sqrt{x^{2}}\le\sqrt{x^{2}+y^{2}}$ per la monotonia della
radice, e analogamente per la quarta.
(iv) Usando (i), (ii) e il fatto che $\bar zw=\overline{z\bar w}$,
\[
  \abs{z+w}^{2}=(z+w)(\bar z+\bar w)=\abs z^{2}+\abs w^{2}+z\bar w+\overline{z\bar w}
  =\abs z^{2}+\abs w^{2}+2\operatorname{Re}(z\bar w).
\]
Per (iii), $\operatorname{Re}(z\bar w)\le\abs{z\bar w}=\abs z\abs w$, quindi
$\abs{z+w}^{2}\le\abs z^{2}+\abs w^{2}+2\abs z\abs w=(\abs z+\abs w)^{2}$; estraendo la
radice si ottiene la tesi.
\end{dimostrazione}

Il coniugato è particolarmente utile nel calcolo del quoziente di due numeri complessi:
moltiplicando numeratore e denominatore per il coniugato del denominatore, il
denominatore diventa reale. Se $z=x+iy\neq0$,
\begin{equation}\label{eq:quoziente-algebrico}
\begin{split}
  \frac{x'+iy'}{x+iy}&=\frac{(x'+iy')(x-iy)}{(x+iy)(x-iy)}=\frac{(x'+iy')(x-iy)}{x^{2}+y^{2}}\\
  &=\frac{xx'+yy'}{x^{2}+y^{2}}-i\,\frac{x'y-xy'}{x^{2}+y^{2}}.
\end{split}
\end{equation}
Come si vede dal conto precedente, la divisione tra due numeri complessi è possibile
purché il denominatore $z=x+iy$ sia diverso da zero, cioè purché non risulti
contemporaneamente $x=0$ e $y=0$ (è la \eqref{eq:inverso-complesso}).

\begin{esempio}
Calcoliamo un quoziente moltiplicando per il coniugato del denominatore:
\[
  \frac{3+2i}{1-i}=\frac{(3+2i)(1+i)}{(1-i)(1+i)}=\frac{3+3i+2i+2i^{2}}{2}=\frac{1+5i}{2}
  =\frac12+\frac52\,i .
\]
\end{esempio}

\subsection{Il piano di Gauss e la forma trigonometrica}

È utile la \emph{rappresentazione cartesiana} dei numeri complessi, che si ottiene
facendo corrispondere a ogni numero complesso $z=x+iy$ il punto $P$ di coordinate
$(x,y)$, come nella \cref{fig:piano-gauss}: l'asse delle ascisse rappresenta i numeri
reali (\emph{asse reale}), quello delle ordinate i multipli reali di $i$ (\emph{asse
immaginario}). Il piano così interpretato si chiama \emph{piano di Gauss} (o di
Argand--Gauss). Il modulo $\abs z$ è la distanza di $P$ dall'origine (per il teorema di
Pitagora), il coniugato $\bar z$ è il simmetrico di $z$ rispetto all'asse reale, e la
somma di due numeri complessi corrisponde alla regola del parallelogramma per la somma
di vettori.

\begin{figure}[htb]
\centering
\begin{tikzpicture}[scale=1.05, every node/.style={font=\footnotesize}]
  \draw[-Latex, grigio] (-0.5,0) -- (4.3,0) node[below] {$\operatorname{Re}$};
  \draw[-Latex, grigio] (0,-2.3) -- (0,2.6) node[left] {$\operatorname{Im}$};
  \coordinate (z) at (3,1.8);
  \draw[blu, thick, -Latex] (0,0) -- (z) node[midway, above left] {$\rho=\abs z$};
  \fill[blu] (z) circle (1.6pt) node[above right] {$z=x+iy$};
  \draw[tratteggio] (z) -- (3,0) node[below] {$x$};
  \draw[tratteggio] (z) -- (0,1.8) node[left] {$y$};
  \draw[rame, thick] (0.8,0) arc (0:31:0.8) node[midway, right] {$\vartheta$};
  \coordinate (zb) at (3,-1.8);
  \draw[verde, thick, -Latex] (0,0) -- (zb);
  \fill[verde] (zb) circle (1.6pt) node[below right] {$\bar z=x-iy$};
  \draw[tratteggio] (3,0) -- (zb);
  \node[below left] at (0,0) {$0$};
\end{tikzpicture}\hfill
\begin{tikzpicture}[scale=1.05, every node/.style={font=\footnotesize}]
  \draw[-Latex, grigio] (-0.5,0) -- (4.3,0) node[below] {$\operatorname{Re}$};
  \draw[-Latex, grigio] (0,-0.5) -- (0,3.3) node[left] {$\operatorname{Im}$};
  \coordinate (z) at (2.6,0.8); \coordinate (w) at (0.8,2.0); \coordinate (s) at (3.4,2.8);
  \fill[bluchiaro] (0,0) -- (z) -- (s) -- (w) -- cycle;
  \draw[blu, thick, -Latex] (0,0) -- (z) node[below right] {$z$};
  \draw[rame, thick, -Latex] (0,0) -- (w) node[above left] {$w$};
  \draw[verde, very thick, -Latex] (0,0) -- (s) node[above right] {$z+w$};
  \draw[tratteggio] (z) -- (s) -- (w);
\end{tikzpicture}
\caption{A sinistra: il piano di Gauss; il numero complesso $z$ è il punto di coordinate
$(x,y)$, individuato anche dal modulo $\rho$ e dall'argomento $\vartheta$; il coniugato
$\bar z$ è il simmetrico rispetto all'asse reale. A destra: la somma di numeri complessi
segue la regola del parallelogramma; la disuguaglianza triangolare dice che la diagonale
non supera la somma dei lati.}
\label{fig:piano-gauss}
\end{figure}

Un punto $P\neq0$ del piano può essere individuato anche dalla distanza $\rho$
dall'origine e dall'angolo $\vartheta$ che il segmento $OP$ forma con il semiasse reale
positivo. Dal teorema di Pitagora si deduce che
\begin{equation}\label{eq:rho}
  \rho=\sqrt{x^{2}+y^{2}}=\abs z;
\end{equation}
l'angolo $\vartheta$ è legato alla parte reale $x$ e al coefficiente dell'immaginario
$y$ dalle formule
\begin{equation}\label{eq:theta}
  \cos\vartheta=\frac x\rho, \qquad \sen\vartheta=\frac y\rho .
\end{equation}
Un tale $\vartheta$ esiste: poiché $(x/\rho)^{2}+(y/\rho)^{2}=1$, il punto
$(x/\rho,y/\rho)$ sta sulla circonferenza di raggio $1$, e per la definizione di seno e
coseno (\cref{sec:trigonometriche}) è il punto corrispondente a un opportuno angolo
$\vartheta$.

\begin{definizione}[Argomento e forma trigonometrica]\index{argomento di un numero complesso}\index{forma trigonometrica}
Se $z\neq0$, il numero $\vartheta$ della \eqref{eq:theta}, misurato in radianti, si
chiama \emph{argomento} di $z$; esso è definito a meno di multipli interi di $2\pi$.
Tenendo conto delle relazioni tra $\rho$, $\vartheta$ e $x$, $y$, il numero complesso
$z=x+iy$ si scrive in \emph{forma trigonometrica}
\begin{equation}\label{eq:forma-trigonometrica}
  z=\rho(\cos\vartheta+i\sen\vartheta).
\end{equation}
\end{definizione}

\begin{teorema}[Prodotto e quoziente in forma trigonometrica]\label{teo:prodotto-trigonometrico}
Siano
\begin{equation}\label{eq:due-complessi}
  z=\rho(\cos\vartheta+i\sen\vartheta), \qquad z'=\rho'(\cos\vartheta'+i\sen\vartheta').
\end{equation}
Allora
\begin{equation}\label{eq:prodotto-trigonometrico}
  z\cdot z'=\rho\rho'\bigl[\cos(\vartheta+\vartheta')+i\sen(\vartheta+\vartheta')\bigr],
\end{equation}
e, se $z'\neq0$,
\begin{equation}\label{eq:quoziente-trigonometrico}
  \frac{z}{z'}=\frac{\rho}{\rho'}\bigl[\cos(\vartheta-\vartheta')+i\sen(\vartheta-\vartheta')\bigr].
\end{equation}
Cioè: \emph{il prodotto di due numeri complessi ha per modulo il prodotto dei moduli e
per argomento la somma degli argomenti}; il quoziente ha per modulo il quoziente dei
moduli e per argomento la differenza degli argomenti.
\end{teorema}

\begin{schema}
\Ipotesi{$z$, $z'$ in forma trigonometrica \eqref{eq:due-complessi}, con
$\rho,\rho'\ge0$; \ per il quoziente, $z'\neq0$ (cioè $\rho'>0$).}
\Tesi{Valgono la \eqref{eq:prodotto-trigonometrico} e la
\eqref{eq:quoziente-trigonometrico}.}
\end{schema}

\begin{strategia}
Si moltiplica in forma algebrica e si riconoscono, nella parte reale e nel coefficiente
dell'immaginario, i secondi membri delle formule di addizione del coseno e del seno
(\eqref{eq:addizione-cos}, \eqref{eq:addizione-sen}). Per il quoziente conviene non
rifare il conto: si calcola l'inverso di $\cos\vartheta'+i\sen\vartheta'$, che è il suo
coniugato $\cos(-\vartheta')+i\sen(-\vartheta')$, e si applica la formula del prodotto.
\end{strategia}

\begin{dimostrazione}
Tenendo conto delle formule di addizione \eqref{eq:addizione-sen},
\eqref{eq:addizione-cos} si ottiene
\begin{equation}\label{eq:prodotto-passo}
\begin{split}
  z\cdot z'&=\rho\rho'(\cos\vartheta+i\sen\vartheta)(\cos\vartheta'+i\sen\vartheta')\\
  &=\rho\rho'\bigl[(\cos\vartheta\cos\vartheta'-\sen\vartheta\sen\vartheta')\\
  &\qquad\quad+i(\sen\vartheta\cos\vartheta'+\sen\vartheta'\cos\vartheta)\bigr]\\
  &=\rho\rho'\bigl[\cos(\vartheta+\vartheta')+i\sen(\vartheta+\vartheta')\bigr].
\end{split}
\end{equation}
Per il quoziente, osserviamo che
$(\cos\vartheta'+i\sen\vartheta')(\cos\vartheta'-i\sen\vartheta')=\cos^{2}\vartheta'+\sen^{2}\vartheta'=1$
per la relazione fondamentale \eqref{eq:relazione-fondamentale}; quindi, essendo il
coseno pari e il seno dispari,
\[
  \frac{1}{z'}=\frac{1}{\rho'}(\cos\vartheta'-i\sen\vartheta')
  =\frac{1}{\rho'}\bigl[\cos(-\vartheta')+i\sen(-\vartheta')\bigr],
\]
e la \eqref{eq:quoziente-trigonometrico} segue dalla \eqref{eq:prodotto-trigonometrico}
applicata a $z\cdot\frac{1}{z'}$.
\end{dimostrazione}

\begin{figure}[htb]
\centering
\begin{tikzpicture}[scale=1.25, every node/.style={font=\footnotesize}]
  \draw[-Latex, grigio] (-0.4,0) -- (3.6,0) node[below] {$\operatorname{Re}$};
  \draw[-Latex, grigio] (0,-0.3) -- (0,3.2) node[left] {$\operatorname{Im}$};
  \draw[grigio!50] (1,0) arc (0:100:1);
  % z: modulo 1.6, arg 20°; w: modulo 1.25, arg 45°; zw: modulo 2, arg 65°
  \draw[blu, thick, -Latex] (0,0) -- (20:1.6) node[right] {$z$};
  \draw[rame, thick, -Latex] (0,0) -- (45:1.25) node[above right] {$z'$};
  \draw[verde, very thick, -Latex] (0,0) -- (65:2.0) node[above] {$zz'$};
  \draw[blu] (0.55,0) arc (0:20:0.55); \node[blu] at (10:0.75) {$\vartheta$};
  \draw[verde] (20:0.95) arc (20:65:0.95); \node[verde] at (32:1.2) {$\vartheta'$};
  \node[grigio, align=left, anchor=west] at (1.9,2.6) {$\abs{zz'}=\abs z\,\abs{z'}$\\ $\arg(zz')=\vartheta+\vartheta'$};
\end{tikzpicture}
\caption{Il significato geometrico della \eqref{eq:prodotto-trigonometrico}: moltiplicare
per $z'$ significa \emph{ruotare} di un angolo $\vartheta'$ e \emph{dilatare} di un
fattore $\abs{z'}$. In particolare, moltiplicare per un numero di modulo $1$ è una
rotazione pura; moltiplicare per $i$ ruota di un angolo retto.}
\label{fig:prodotto-rotazione}
\end{figure}

\begin{notastorica}[Wessel, Argand, Gauss: i numeri diventano punti]
Per più di due secoli dopo Bombelli i numeri complessi furono usati con profitto, ma senza
che nessuno sapesse dire che cosa fossero: ancora Eulero, che li maneggiava con grande
disinvoltura, li chiamava «impossibili». La svolta venne dalla geometria. Il norvegese
Caspar Wessel (1745--1818), cartografo al servizio dell'Accademia danese delle scienze,
cercava un modo di calcolare con i segmenti orientati: in una memoria presentata
all'Accademia nel 1797, e pubblicata nel 1799, li sommava con la regola del
parallelogramma e li moltiplicava moltiplicando le lunghezze e sommando gli angoli,
proprio come nel \cref{teo:prodotto-trigonometrico}. Scritta in danese, la memoria restò
ignorata fino alla traduzione francese del 1897. Nel 1806 Jean-Robert Argand, un
dilettante che viveva a Parigi, fece circolare in poche copie, senza firmarlo, un
\emph{Essai} con la stessa idea; lo si conobbe solo nel 1813, quando Jacques Français ne
espose il contenuto negli \emph{Annales} di Gergonne invitando l'autore sconosciuto a
farsi avanti. A imporre la rappresentazione fu l'autorità di Gauss, che nel 1831 la difese
pubblicamente e coniò il nome di «numero complesso»: se $+1$, $-1$ e $\sqrt{-1}$, scrisse,
si fossero chiamate unità «diretta, inversa e laterale» anziché positiva, negativa e
immaginaria, nessuno vi avrebbe visto alcun mistero. È la lezione della
\cref{fig:prodotto-rotazione}: $i$ non è un numero misterioso, è una rotazione di un
angolo retto.
\end{notastorica}

\begin{esempio}[Dagli assi vento agli assi corda di un profilo]\label{es:assi-profilo}
Il significato «rotatorio» del prodotto è uno strumento di calcolo molto pratico. Un
profilo alare investito da una corrente di velocità $V_{\infty}$ con incidenza $\alpha$
è soggetto a una forza aerodinamica che nel riferimento \emph{vento} (asse $x$ nella
direzione della corrente, asse $y$ perpendicolare, verso l'alto) ha componenti $D$
(resistenza) e $L$ (portanza). Rappresentiamo la forza con il numero complesso
$F_{v}=D+iL$. Nel riferimento \emph{corda} (asse lungo la corda, dal bordo d'attacco al
bordo d'uscita, e asse normale), che con il profilo cabrato è ruotato di un angolo
$-\alpha$ rispetto al riferimento vento, le componenti si ottengono ruotando il vettore
in senso opposto, cioè moltiplicando per $\cos\alpha+i\sen\alpha$:
\[
  F_{c}=F_{v}(\cos\alpha+i\sen\alpha)=(D\cos\alpha-L\sen\alpha)+i\,(L\cos\alpha+D\sen\alpha).
\]
La parte reale è la \emph{forza assiale} $A=D\cos\alpha-L\sen\alpha$, il coefficiente
dell'immaginario è la \emph{forza normale} $N=L\cos\alpha+D\sen\alpha$: sono le formule
che si trovano in ogni testo di aerodinamica, ottenute qui con una sola moltiplicazione.
Con $L=1000$~N, $D=50$~N e $\alpha=5^{\circ}$ si trova $N\approx1000{,}6$~N e
$A\approx-37{,}3$~N: la forza assiale è
\emph{negativa}, cioè diretta verso il bordo d'attacco (è l'effetto della
«suzione» sul bordo d'attacco). Il passaggio inverso si ottiene moltiplicando per
$\cos\alpha-i\sen\alpha$, cioè dividendo: $F_{v}=F_{c}/(\cos\alpha+i\sen\alpha)$.
\end{esempio}

\begin{figure}[htb]
\centering
\begin{tikzpicture}[scale=1.0, every node/.style={font=\footnotesize},
  etich/.style={fill=white, inner sep=1pt, fill opacity=0.85, text opacity=1}]
  \def\alf{15}\def\Lv{2.6}\def\Dv{0.3}
  \pgfmathsetmacro{\Nv}{\Lv*cos(\alf)+\Dv*sin(\alf)}
  \pgfmathsetmacro{\Av}{\Dv*cos(\alf)-\Lv*sin(\alf)}
  % corrente
  \foreach \yy in {-0.7,-0.1,0.5} {\draw[-Latex, grigio!70] (-3.6,\yy) -- (-2.5,\yy);}
  \node[grigio] at (-3.05,1.0) {$V_{\infty}$};
  % profilo NACA 0012 cabrato di alfa, con il quarto di corda nell'origine
  \begin{scope}[rotate=-\alf, scale=6.5, shift={(-0.25,0)}]
    \filldraw[fill=bluchiaro, draw=blu, thick]
      plot[domain=0:1, samples=70, smooth] (\x, {0.6*(0.2969*sqrt(\x)-0.1260*\x-0.3516*\x^2+0.2843*\x^3-0.1015*\x^4)})
      -- plot[domain=1:0, samples=70, smooth] (\x, {-0.6*(0.2969*sqrt(\x)-0.1260*\x-0.3516*\x^2+0.2843*\x^3-0.1015*\x^4)}) -- cycle;
    \draw[grigio, dashed, thin] (-0.08,0) -- (1.25,0);
  \end{scope}
  % riferimento orizzontale e angolo di incidenza
  \draw[grigio, dotted] (0,0) -- (5.6,0);
  \draw[grigio] (4.6,0) arc (0:-\alf:4.6);
  \node[grigio] at ({4.85*cos(\alf/2)},{-4.85*sin(\alf/2)}) {$\alpha$};
  % risultante
  \draw[-Latex, grigio, thick] (0,0) -- (\Dv,\Lv) node[above=1pt, etich] {$R$};
  % assi vento: D, L
  \draw[-Latex, rame, very thick] (0,0) -- (0,\Lv) node[left=5pt, etich] {$L$};
  \draw[-Latex, rame, very thick] (0,0) -- (\Dv,0);
  \node[rame, etich] at (0.45,-0.28) {$D$};
  % assi corda: N, A
  \draw[-Latex, verde, very thick] (0,0) -- ({\Nv*sin(\alf)},{\Nv*cos(\alf)}) node[right=2pt, etich] {$N$};
  \draw[-Latex, verde, very thick] (0,0) -- ({\Av*cos(\alf)},{-\Av*sin(\alf)});
  \node[verde, etich] at (-1.05,0.75) {$A<0$};
  \draw[verde, thin] (-0.8,0.62) -- ({0.6*\Av*cos(\alf)},{-0.6*\Av*sin(\alf)+0.04});
  \fill (0,0) circle (1.3pt);
\end{tikzpicture}
\caption{Un profilo NACA 0012 a incidenza $\alpha$ (qui esagerata a $15^{\circ}$ per
leggibilità; forze non in scala). La risultante $R$ è disegnata nel quarto di corda, che
per un profilo simmetrico è, in prima approssimazione, il centro di pressione. La forza
aerodinamica ha componenti $(D,L)$ negli assi vento e $(A,N)$ negli assi corda: si
passa dalle une alle altre moltiplicando per
$\cos\alpha\pm i\sen\alpha$.}
\label{fig:profilo-assi}
\end{figure}

\subsection{Potenze e radici}

\begin{corollario}[Formula di De Moivre]\label{cor:de-moivre}\index{formula di De Moivre}
Se $z=\rho(\cos\vartheta+i\sen\vartheta)$, per ogni $n\in\N$
\begin{equation}\label{eq:de-moivre}
  z^{n}=\rho^{n}(\cos n\vartheta+i\sen n\vartheta).
\end{equation}
La formula vale anche per $n=0$ e, se $z\neq0$, per $n$ intero negativo.
\end{corollario}

\begin{schema}
\Ipotesi{$z=\rho(\cos\vartheta+i\sen\vartheta)$; \ $n\in\N$, oppure $n=0$, oppure
$n\in\Z$ negativo e $z\neq0$.}
\Tesi{$z^{n}=\rho^{n}(\cos n\vartheta+i\sen n\vartheta)$.}
\end{schema}

\begin{strategia}
Per $n\in\N$ si procede per induzione: per il \cref{teo:prodotto-trigonometrico}, ogni
moltiplicazione per $z$ moltiplica il modulo per $\rho$ e aggiunge $\vartheta$
all'argomento. Il caso $n=0$ si riduce alla convenzione $z^{0}=1$; per gli esponenti
negativi si scrive $z^{-m}=1/z^{m}$ e si usa la formula del quoziente.
\end{strategia}

\begin{dimostrazione}
Per $n=1$ è ovvia; se vale per $n$, per la \eqref{eq:prodotto-trigonometrico}
$z^{n+1}=z^{n}\cdot z=\rho^{n}\rho\bigl[\cos(n\vartheta+\vartheta)+i\sen(n\vartheta+\vartheta)\bigr]$,
che è la \eqref{eq:de-moivre} con $n+1$. Per $n=0$ entrambi i membri valgono $1$; per
$n=-m$ con $m\in\N$, si applica la \eqref{eq:quoziente-trigonometrico} a $1/z^{m}$.
\end{dimostrazione}

Un numero complesso $w$ si dice \emph{radice $n$-esima} di $z$ se risulta $w^{n}=z$.
Nel campo reale un numero positivo ha una sola radice $n$-esima positiva, e i numeri
negativi non hanno radici di ordine pari; nel campo complesso la situazione è molto più
regolare.

\begin{teorema}[Radici $n$-esime di un numero complesso]\label{teo:radici-complesse}\index{radice $n$-esima!di un numero complesso}
Sia $z=\rho(\cos\vartheta+i\sen\vartheta)\neq0$ e sia $n\in\N$. Allora $z$ ha esattamente
$n$ radici $n$-esime distinte, date da
\begin{equation}\label{eq:radici-complesse}
  w_{k}=\sqrt[n]{\rho}\,\Bigl(\cos\frac{\vartheta+2k\pi}{n}+i\sen\frac{\vartheta+2k\pi}{n}\Bigr),
  \qquad k=0,1,\dots,n-1.
\end{equation}
\end{teorema}

\begin{schema}
\Ipotesi{$z\in\C$, $z\neq0$, in forma trigonometrica con modulo $\rho>0$ e argomento
$\vartheta$; \ $n\in\N$.}
\Tesi{Le soluzioni dell'equazione $w^{n}=z$ sono esattamente gli $n$ numeri distinti
$w_{0},\dots,w_{n-1}$ della \eqref{eq:radici-complesse}.}
\end{schema}

\begin{strategia}
Cerchiamo $w$ in forma trigonometrica, $w=r(\cos\varphi+i\sen\varphi)$. Per la formula
di De Moivre l'equazione $w^{n}=z$ diventa un'uguaglianza tra due numeri in forma
trigonometrica, e due numeri complessi non nulli sono uguali se e solo se hanno lo
stesso modulo e argomenti che differiscono per un multiplo di $2\pi$. Si ottengono così
due condizioni separate: una sul modulo, $r^{n}=\rho$, che ha l'unica soluzione
$r=\sqrt[n]\rho$ per il \cref{teo:radice-n}; e una sull'argomento,
$n\varphi=\vartheta+2k\pi$, che ha infinite soluzioni $\varphi=(\vartheta+2k\pi)/n$, una
per ogni $k\in\Z$. Resta da contare quante di queste danno numeri \emph{diversi}: due
valori di $k$ danno lo stesso $w$ se e solo se gli argomenti differiscono per un multiplo
di $2\pi$, cioè se e solo se i due $k$ differiscono per un multiplo di $n$. I valori
$k=0,\dots,n-1$ sono quindi tutti e soli i rappresentanti distinti.
\end{strategia}

\begin{dimostrazione}
Poiché $z\neq0$, anche ogni radice $w$ è diversa da $0$ (se fosse $w=0$ sarebbe
$z=w^{n}=0$); possiamo quindi scrivere $w=r(\cos\varphi+i\sen\varphi)$, con $r>0$. Per la
\eqref{eq:de-moivre},
\begin{equation}\label{eq:radici-passo}
  w^{n}=z \quad\sse\quad r^{n}(\cos n\varphi+i\sen n\varphi)=\rho(\cos\vartheta+i\sen\vartheta).
\end{equation}
Prendendo i moduli, $r^{n}=\rho$, quindi $r=\sqrt[n]\rho$ (\cref{teo:radice-n}); allora
la \eqref{eq:radici-passo} equivale a $\cos n\varphi=\cos\vartheta$ e
$\sen n\varphi=\sen\vartheta$, cioè al fatto che gli angoli $n\varphi$ e $\vartheta$
individuino lo stesso punto della circonferenza di raggio $1$; ciò accade se e solo se
$n\varphi=\vartheta+2k\pi$ per qualche $k\in\Z$. Le radici sono dunque tutti e soli i
numeri della \eqref{eq:radici-complesse} con $k\in\Z$.

Siano $h,k\in\Z$. Gli argomenti $\frac{\vartheta+2h\pi}{n}$ e $\frac{\vartheta+2k\pi}{n}$
differiscono di $\frac{2(h-k)\pi}{n}$, che è un multiplo intero di $2\pi$ se e solo se
$h-k$ è un multiplo di $n$: quindi $w_{h}=w_{k}$ se e solo se $n$ divide $h-k$. In
particolare $w_{0},\dots,w_{n-1}$ sono distinti (due indici tra $0$ e $n-1$ differiscono
per meno di $n$, e non sono uguali). D'altra parte, ogni $k\in\Z$ si scrive come
$k=qn+s$ con $q=\lfloor k/n\rfloor\in\Z$ e $s\in\{0,\dots,n-1\}$
(\cref{prop:parte-intera}), e allora $w_{k}=w_{s}$. Le radici distinte sono quindi
esattamente $n$.
\end{dimostrazione}

In particolare il modulo di una radice $n$-esima è uguale alla radice $n$-esima
(aritmetica) del modulo, e gli argomenti delle $n$ radici sono equidistanziati di
$2\pi/n$: \emph{le radici $n$-esime di $z$ sono i vertici di un poligono regolare di $n$
lati inscritto nella circonferenza di centro $0$ e raggio $\sqrt[n]{\abs z}$}
(\cref{fig:radici-unita}).

\begin{esempio}[Le radici quadrate di $i$]\label{es:radici-i}
Calcoliamo le radici quadrate del numero complesso $z=i$. In forma trigonometrica $i$ ha
modulo uguale a $1$ e argomento uguale a $\pi/2$. Quindi, in base alla
\eqref{eq:radici-complesse}, le due radici quadrate di $i$ hanno modulo uguale a $1$ e
argomento uguale a
\begin{equation}\label{eq:argomenti-radici-i}
  \varphi_{0}=\frac{\pi/2}{2}=\frac\pi4, \qquad \varphi_{1}=\frac{\pi/2+2\pi}{2}=\frac54\pi .
\end{equation}
In corrispondenza si ottiene
\begin{align*}
  w_{0}&=\cos\frac\pi4+i\sen\frac\pi4=\frac{\sqrt2}{2}+i\,\frac{\sqrt2}{2},\\
  w_{1}&=\cos\frac54\pi+i\sen\frac54\pi=-\frac{\sqrt2}{2}-i\,\frac{\sqrt2}{2}.
\end{align*}
Quindi le due radici quadrate di $i$ sono
\begin{equation}\label{eq:radici-i}
  \pm\Bigl(\frac{\sqrt2}{2}+i\,\frac{\sqrt2}{2}\Bigr)=\pm\frac{\sqrt2}{2}(1+i).
\end{equation}
Verifica: $\bigl(\frac{\sqrt2}{2}(1+i)\bigr)^{2}=\frac12(1+2i+i^{2})=\frac12\cdot2i=i$.
\end{esempio}

\begin{figure}[htb]
\centering
\begin{tikzpicture}[scale=1.3, every node/.style={font=\footnotesize}]
  \draw[-Latex, grigio] (-1.45,0) -- (1.55,0) node[above left] {$\operatorname{Re}$};
  \draw[-Latex, grigio] (0,-1.35) -- (0,1.5) node[below left] {$\operatorname{Im}$};
  \draw[grigio!60] (0,0) circle (1);
  \draw[blu, thick, fill=bluchiaro, fill opacity=0.5] (0:1) \foreach \k in {1,...,5} { -- (\k*60:1)} -- cycle;
  \foreach \k in {0,...,5} {\fill[blu] (\k*60:1) circle (1.4pt);}
  \node[below right] at (0:1) {$w_{0}$};
  \node[above right] at (60:1) {$w_{1}$};
  \node[above left] at (120:1) {$w_{2}$};
  \node[below left] at (180:1) {$w_{3}$};
  \node[below left] at (240:1) {$w_{4}$};
  \node[below right] at (300:1) {$w_{5}$};
  \draw[rame] (0.35,0) arc (0:60:0.35); \node[rame] at (30:0.55) {$\frac{\pi}{3}$};
\end{tikzpicture}\hfill
\begin{tikzpicture}[scale=1.3, every node/.style={font=\footnotesize}]
  \draw[-Latex, grigio] (-1.45,0) -- (1.55,0) node[above left] {$\operatorname{Re}$};
  \draw[-Latex, grigio] (0,-1.35) -- (0,1.5) node[below left] {$\operatorname{Im}$};
  \draw[grigio!60] (0,0) circle (1);
  \fill[verde] (90:1) circle (1.5pt) node[above right] {$i$};
  \draw[rame, thick] (45:1) -- (225:1);
  \fill[rame] (45:1) circle (1.5pt) node[above right] {$w_{0}$};
  \fill[rame] (225:1) circle (1.5pt) node[below left] {$w_{1}$};
  \draw[verde, -Latex] (0.55,0) arc (0:90:0.55); \node[verde] at (68:0.75) {$\frac\pi2$};
  \draw[rame, -Latex] (0.3,0) arc (0:45:0.3); \node[rame] at (22:0.5) {$\frac\pi4$};
\end{tikzpicture}
\caption{A sinistra: le sei radici seste dell'unità $w_{k}=\cos\frac{k\pi}{3}+i\sen\frac{k\pi}{3}$
($w_{0}=1$, $w_{1}=\frac12+i\frac{\sqrt3}{2}$, $w_{3}=-1$, \dots), vertici di un esagono
regolare inscritto nella circonferenza unitaria. A destra: le due radici quadrate di $i$
(\cref{es:radici-i}), $w_{0}=\frac{\sqrt2}{2}(1+i)$ e $w_{1}=-w_{0}$, con argomenti $\pi/4$
e $5\pi/4$: il «poligono regolare di due lati» è un diametro.}
\label{fig:radici-unita}
\end{figure}

\subsection{Equazioni di secondo grado}

Possiamo ora mantenere la promessa iniziale: in $\C$ ogni equazione di secondo grado ha
soluzione.

\needspace{10\baselineskip}
\begin{proposizione}[Equazioni di secondo grado in $\C$]\label{prop:secondo-grado-C}
Siano $a,b,c\in\C$ con $a\neq0$, e sia $\Delta=b^{2}-4ac$. Detta $\delta$ una qualsiasi
radice quadrata complessa di $\Delta$ (che esiste per il \cref{teo:radici-complesse},
oppure è $\delta=0$ se $\Delta=0$), le soluzioni dell'equazione $az^{2}+bz+c=0$ sono
\begin{equation}\label{eq:soluzioni-C}
  z_{1,2}=\frac{-b\pm\delta}{2a},
\end{equation}
e valgono le relazioni \eqref{eq:viete}. In particolare, se $a,b,c$ sono reali e
$\Delta<0$, le soluzioni sono complesse coniugate:
$z_{1,2}=-\dfrac{b}{2a}\pm i\,\dfrac{\sqrt{-\Delta}}{2a}$.
\end{proposizione}

\begin{schema}
\Ipotesi{$a,b,c\in\C$, \ $a\neq0$; \ $\Delta=b^{2}-4ac$; \ $\delta\in\C$ con
$\delta^{2}=\Delta$.}
\Tesi{$az^{2}+bz+c=0\sse z=\frac{-b\pm\delta}{2a}$, \ e valgono le \eqref{eq:viete}; \
se $a,b,c\in\R$ e $\Delta<0$, \ $z_{1,2}=-\frac{b}{2a}\pm i\frac{\sqrt{-\Delta}}{2a}$.}
\end{schema}

\begin{strategia}
Si ripete il completamento del quadrato, che usa soltanto le proprietà di campo e
quindi vale in $\C$ senza modifiche. L'unico ingrediente nuovo è l'esistenza di una
radice quadrata $\delta$ di $\Delta$, garantita dal \cref{teo:radici-complesse} anche
quando $\Delta$ non è un reale non negativo. Scritto il trinomio come prodotto di due
fattori di primo grado, la legge di annullamento del prodotto conclude.
\end{strategia}

\begin{dimostrazione}
Il «completamento del quadrato» usa solo gli assiomi di campo, e vale quindi anche in
$\C$:
\[
  az^{2}+bz+c=a\Bigl[\Bigl(z+\frac{b}{2a}\Bigr)^{2}-\frac{b^{2}-4ac}{4a^{2}}\Bigr]
  =a\Bigl(z+\frac{b}{2a}-\frac{\delta}{2a}\Bigr)\Bigl(z+\frac{b}{2a}+\frac{\delta}{2a}\Bigr),
\]
dove abbiamo usato $\delta^{2}=\Delta$ e la scomposizione $u^{2}-v^{2}=(u-v)(u+v)$. Per
la legge di annullamento del prodotto (valida in $\C$, che è un campo) il primo membro è
nullo se e solo se uno dei due fattori lo è, cioè se e solo se $z$ è uno dei numeri
\eqref{eq:soluzioni-C}. Le relazioni \eqref{eq:viete} seguono sommando e moltiplicando:
$z_{1}+z_{2}=-b/a$ e $z_{1}z_{2}=(b^{2}-\delta^{2})/(4a^{2})=c/a$. Se infine $a,b,c$ sono
reali e $\Delta<0$, una radice quadrata di $\Delta$ è $\delta=i\sqrt{-\Delta}$, perché
$(i\sqrt{-\Delta})^{2}=i^{2}(-\Delta)=\Delta$.
\end{dimostrazione}

Ad esempio, l'equazione $z^{2}+1=0$ da cui eravamo partiti ha $\Delta=-4$, $\delta=2i$ e
soluzioni $z_{1,2}=\pm i$; e ora le relazioni $z_{1}+z_{2}=0$, $z_{1}z_{2}=1$ hanno
pienamente senso: $i+(-i)=0$ e $i\cdot(-i)=-i^{2}=1$.

\begin{notastorica}[Il teorema fondamentale dell'algebra]
La \cref{prop:secondo-grado-C} è il caso più semplice di un risultato molto più
generale, il \emph{teorema fondamentale dell'algebra}: \emph{ogni polinomio di grado
$n\ge1$ a coefficienti complessi ha almeno una radice in $\C$} (e quindi, scomponendo,
esattamente $n$ radici contate con la loro molteplicità). Enunciato già da Albert Girard
(1629) e tentato da d'Alembert (1746) e da Eulero, fu dimostrato in modo sostanzialmente
completo da Carl Friedrich Gauss nella sua tesi di dottorato (1799; la dimostrazione
conteneva una lacuna di natura topologica, colmata solo nel Novecento), e poi da
Jean-Robert Argand (1806, in una versione ripresa nel 1813), il dilettante incontrato
nella nota sul piano di Gauss, che fu anche il primo a enunciarlo esplicitamente per i
polinomi a coefficienti complessi. Il teorema dice che con $\C$ il processo di ampliamento iniziato passando da $\N$
a $\Z$ (per sottrarre), proseguito con $\Q$ (per dividere) e con $\R$ (per estrarre radici
e passare al limite) si \emph{chiude}: non occorrono altri numeri per risolvere le equazioni
algebriche.
\end{notastorica}

\subsection{La forma esponenziale}

Segnaliamo infine una notazione \emph{di tipo esponenziale} per i numeri complessi.
Posto
\begin{equation}\label{eq:eulero}
  e^{i\vartheta}=\cos\vartheta+i\sen\vartheta \qquad (\vartheta\in\R),
\end{equation}
un numero complesso $z=\rho(\cos\vartheta+i\sen\vartheta)$ si potrà rappresentare nella
forma
\begin{equation}\label{eq:forma-esponenziale}
  z=\rho\,e^{i\vartheta}.
\end{equation}
Per ora la \eqref{eq:eulero} è soltanto una \emph{definizione}, cioè un'abbreviazione:
nel capitolo sulle serie, dopo aver definito $e^{z}$ per ogni $z\in\C$ come somma della
serie esponenziale, vedremo che essa è in realtà un teorema, la \emph{formula di Eulero}
(\cref{teo:formula-eulero}), che lega la funzione esponenziale alle funzioni
trigonometriche. La notazione,
in ogni caso, è giustificata dal fatto che si comporta come una potenza. In tal modo,
infatti, la \eqref{eq:prodotto-trigonometrico} si riscrive come
\begin{equation}\label{eq:prodotto-esponenziale}
  z\cdot z'=\rho\,e^{i\vartheta}\cdot\rho'e^{i\vartheta'}=\rho\rho'\,e^{i(\vartheta+\vartheta')},
\end{equation}
che è la stessa regola $e^{u}e^{v}=e^{u+v}$ delle potenze reali; la formula di De Moivre
diventa $(\rho e^{i\vartheta})^{n}=\rho^{n}e^{in\vartheta}$, e le radici
\eqref{eq:radici-complesse} si scrivono
$w_{k}=\sqrt[n]\rho\,e^{i(\vartheta+2k\pi)/n}$. Se $w=r\,e^{i\varphi}$ è una radice
$n$-esima di $z=\rho\,e^{i\vartheta}$, sarà $r^{n}e^{in\varphi}=\rho\,e^{i\vartheta}$,
perciò $r=\sqrt[n]\rho$, $\cos n\varphi=\cos\vartheta$ e $\sen n\varphi=\sen\vartheta$:
ne seguono di nuovo le \eqref{eq:radici-complesse}.

\begin{notastorica}[La «formula più bella»]
Con $\vartheta=\pi$ la \eqref{eq:eulero} dà $e^{i\pi}=\cos\pi+i\sen\pi=-1$, cioè
\[
  e^{i\pi}+1=0,
\]
un'identità che lega le cinque costanti fondamentali della matematica: $0$, $1$, $e$,
$\pi$, $i$. Compare, nella forma $\cos x+i\sen x=e^{ix}$, nell'\emph{Introductio in
analysin infinitorum} di Leonhard Euler (1748); il simbolo $i$ per $\sqrt{-1}$ fu
introdotto dallo stesso Eulero nel 1777 e reso popolare da Gauss. In un sondaggio tra
matematici promosso dalla rivista \emph{The Mathematical Intelligencer} (1988) l'identità
di Eulero fu votata come «il più bel teorema della matematica».
\end{notastorica}

\begin{idea}[Oscillazioni e fasori]
Per un ingegnere i numeri complessi sono soprattutto il linguaggio delle
\emph{oscillazioni}. Uno spostamento armonico, come la flessione all'estremità di un'ala
che vibra, $x(t)=X\cos(\omega t+\phi)$, si scrive
$x(t)=\operatorname{Re}\bigl(X e^{i\phi}\,e^{i\omega t}\bigr)$: il numero complesso
$X e^{i\phi}$ (il \emph{fasore}) contiene ampiezza e fase, e il fattore $e^{i\omega t}$
lo fa ruotare nel piano di Gauss con velocità angolare $\omega$. Sommare due oscillazioni
della stessa frequenza diventa una somma di vettori (\cref{fig:piano-gauss}), derivare
rispetto al tempo diventa moltiplicare per $i\omega$, cioè ruotare di $90^{\circ}$ e
dilatare. È così che si studiano le vibrazioni strutturali, il flutter aeroelastico e i
circuiti in corrente alternata.
\end{idea}

\begin{python}[Numeri complessi in Python]
import cmath, math

z = 3 + 2j          # in Python l'unità immaginaria si scrive j
w = 1 - 1j
print(z / w)                       # quoziente: (0.5+2.5j)
print(abs(z), z.conjugate())       # modulo e coniugato
print(cmath.polar(1j))             # (modulo, argomento) di i

# Radici n-esime con la formula del teorema
def radici(z, n):
    rho, theta = cmath.polar(z)
    r = rho ** (1 / n)
    return [cmath.rect(r, (theta + 2 * k * math.pi) / n)
            for k in range(n)]

for w_k in radici(1j, 2):
    print(f"{w_k:.6f}   (w_k)^2 = {w_k**2:.6f}")

# Dagli assi vento agli assi corda (esempio): F_c = F_v * e^(i alfa)
L, D, alfa = 1000.0, 50.0, math.radians(5)
Fc = complex(D, L) * cmath.exp(1j * alfa)
print(f"A = {Fc.real:.1f} N,  N = {Fc.imag:.1f} N")

# Il calcolo di Bombelli: (2 + i)^3 = 2 + 11i, cioè 2 + sqrt(-121)
print((2 + 1j) ** 3)
\end{python}
\begin{risultato}
(0.5+2.5j)
3.605551275463989 (3-2j)
(1.0, 1.5707963267948966)
0.707107+0.707107j   (w_k)^2 = 0.000000+1.000000j
-0.707107-0.707107j   (w_k)^2 = 0.000000+1.000000j
A = -37.3 N,  N = 1000.6 N
(2+11j)
\end{risultato}

\begin{figure}[htb]
\centering
\includegraphics[width=\textwidth]{radici_unita}
\caption{Le radici $n$-esime dell'unità per $n=3,\dots,8$, calcolate con la
\eqref{eq:radici-complesse} (\texttt{python/fig\_radici\_unita.py}): al crescere di $n$
i poligoni regolari approssimano sempre meglio la circonferenza.}
\label{fig:radici-unita-python}
\end{figure}

%% ----------------------------------------------------------------------------
%%  §2.7  Esercizi di riepilogo
%% ----------------------------------------------------------------------------
\section{Esercizi di riepilogo}
\label{sec:esercizi-cap2}

Come nel capitolo precedente, ogni soluzione è accompagnata da un commento sulle
\emph{scelte logiche} che la guidano. Il lettore provi sempre prima da solo.

\subsection*{Estremo superiore e inferiore}

\begin{esercizio}\label{es:sup-differenze}
Determinare estremo superiore ed estremo inferiore dell'insieme
$A=\set{\dfrac1n-\dfrac1m : n,m\in\N}$, e stabilire se sono massimo e minimo.
\begin{soluzione}
Per ogni $n,m$ si ha $\frac1n\le1$ e $\frac1m>0$, quindi $\frac1n-\frac1m<1$: $1$ è un
maggiorante, e non appartiene ad $A$. Per ogni $\epsilon>0$, scelto $n=1$ e $m>1/\epsilon$
(Archimede), si ha $1-\frac1m>1-\epsilon$: per la \cref{prop:caratt-sup}, $\sup A=1$, e
non è un massimo. Poiché $A$ è simmetrico ($x\in A\sse-x\in A$, scambiando $n$ e $m$),
$\inf A=-1$, che non è un minimo. \emph{Scelta logica:} per un estremo superiore si
dimostrano sempre due cose separate, «è un maggiorante» e «nessun numero più piccolo lo
è»; la seconda richiede di \emph{costruire}, per ogni $\epsilon$, un elemento
dell'insieme che supera $\sup A-\epsilon$, ed è lì che entra in gioco Archimede. La
simmetria evita di rifare i conti per l'estremo inferiore.
\end{soluzione}
\end{esercizio}

\begin{esercizio}
Determinare estremo superiore e inferiore di $B=\set{(-1)^{n}\dfrac{n}{n+1} : n\in\N}$.
\begin{soluzione}
Per $n$ pari gli elementi sono $\frac{n}{n+1}=1-\frac{1}{n+1}\in[\frac23,1)$, per $n$
dispari sono $-\frac{n}{n+1}\in(-1,-\frac12]$. Si ha dunque $-1<b<1$ per ogni $b\in B$.
Fissato $\epsilon>0$, un $n$ pari con $\frac{1}{n+1}<\epsilon$ dà un elemento maggiore di
$1-\epsilon$, e un $n$ dispari con la stessa proprietà dà un elemento minore di
$-1+\epsilon$. Quindi $\sup B=1$, $\inf B=-1$, e né il massimo né il minimo esistono.
\emph{Scelta logica:} il segno alterno suggerisce di spezzare l'insieme in due
sottoinsiemi più semplici (indici pari e dispari), ciascuno «monotòno»; l'estremo
superiore dell'unione è il più grande dei due estremi superiori.
\end{soluzione}
\end{esercizio}

\begin{esercizio}[Estremo superiore della somma]\label{es:sup-somma}
Siano $A,B\subseteq\R$ non vuoti e limitati superiormente, e sia
$A+B=\{a+b : a\in A,\ b\in B\}$. Dimostrare che $\sup(A+B)=\sup A+\sup B$.
\begin{soluzione}
Siano $M=\sup A$, $M'=\sup B$. Per ogni $a\in A$, $b\in B$ si ha $a+b\le M+M'$: $M+M'$ è
un maggiorante di $A+B$. Dato $\epsilon>0$, per la \cref{prop:caratt-sup} applicata con
$\epsilon/2$ esistono $a\in A$ con $a>M-\frac\epsilon2$ e $b\in B$ con
$b>M'-\frac\epsilon2$; allora $a+b>(M+M')-\epsilon$. Quindi $\sup(A+B)=M+M'$.
\emph{Scelta logica:} è il primo esempio del celebre «trucco dell'$\epsilon/2$», che
useremo in moltissime dimostrazioni sui limiti: per ottenere un errore complessivo
minore di $\epsilon$ su una somma di due termini, si chiede a ciascun termine un errore
minore di $\epsilon/2$. Il «budget» di errore si divide tra i contributi.
\end{soluzione}
\end{esercizio}

\begin{esercizio}
Sia $A\neq\emptyset$ limitato inferiormente e sia $-A=\{-a : a\in A\}$. Dimostrare che
$-A$ è limitato superiormente e che $\inf A=-\sup(-A)$.
\begin{soluzione}
$\ell$ è un minorante di $A$ se e solo se $-\ell$ è un maggiorante di $-A$ (moltiplicare
per $-1$ rovescia le disuguaglianze, \cref{prop:ordine-prodotto}). Quindi l'insieme dei
maggioranti di $-A$ è l'opposto dell'insieme dei minoranti di $A$, e il minimo del primo
è l'opposto del massimo del secondo. \emph{Scelta logica:} questa simmetria permette di
dedurre tutti i risultati sull'estremo inferiore da quelli sull'estremo superiore,
«gratis»; è il motivo per cui nel testo ci siamo limitati a dimostrare i secondi.
\end{soluzione}
\end{esercizio}

\begin{esercizio}\label{es:sup-razionali}
Dimostrare che, in $\R$, $\sup\{x\in\Q : x^{2}<2\}=\sqrt2$.
\begin{soluzione}
Se $x\in\Q$ e $x^{2}<2$, allora $x<\sqrt2$ (se fosse $x\ge\sqrt2>0$, sarebbe
$x^{2}\ge2$): $\sqrt2$ è un maggiorante. Dato $\epsilon>0$, se $\epsilon\ge\sqrt2$ basta
l'elemento $x=1>0\ge\sqrt2-\epsilon$; se $\epsilon\in(0,\sqrt2)$, per la densità di
$\Q$ (\cref{teo:densita-Q}) esiste $x\in\Q$ con $\sqrt2-\epsilon<x<\sqrt2$; essendo
$0<x<\sqrt2$, si ha $x^{2}<2$, cioè $x$ appartiene all'insieme. \emph{Scelta logica:}
confrontare con l'\emph{Attenzione} del \cref{sec:sup}: lo stesso insieme non ha estremo
superiore in $\Q$, mentre in $\R$ lo ha. La densità di $\Q$ garantisce che l'insieme si
avvicini quanto si vuole a $\sqrt2$, la completezza di $\R$ garantisce che $\sqrt2$ ci sia.
\end{soluzione}
\end{esercizio}

\subsection*{Calcolo combinatorio e binomio}

\begin{esercizio}[Anagrammi]
Quanti sono gli anagrammi (anche privi di significato) della parola ELICA? E della
parola AEREO?
\begin{soluzione}
ELICA ha $5$ lettere distinte: gli anagrammi sono le permutazioni, $5!=120$. In AEREO la
lettera E compare due volte. Se distinguessimo le due E (E$_1$, E$_2$) avremmo $5!=120$
parole; ma ogni anagramma «vero» corrisponde a $2!=2$ di queste parole (quelle che
differiscono solo per lo scambio di E$_1$ ed E$_2$). Gli anagrammi sono quindi
$120/2=60$. \emph{Scelta logica:} è lo stesso ragionamento con cui abbiamo contato le
combinazioni (\cref{teo:combinazioni}): si conta un insieme più grande e più facile, e
si divide per il numero di oggetti che vengono identificati tra loro. In generale, una
parola di $n$ lettere in cui le lettere si ripetono $n_{1},n_{2},\dots,n_{r}$ volte ha
$\frac{n!}{n_{1}!\,n_{2}!\cdots n_{r}!}$ anagrammi.
\end{soluzione}
\end{esercizio}

\begin{esercizio}
Determinare il coefficiente di $x^{3}$ nello sviluppo di $(2x-1)^{7}$.
\begin{soluzione}
Per la \eqref{eq:binomio} con $a=2x$ e $b=-1$, il termine generico è
$\binom7k(2x)^{7-k}(-1)^{k}$; si ha $x^{3}$ per $7-k=3$, cioè $k=4$, e il coefficiente
è $\binom74\,2^{3}(-1)^{4}=35\cdot8=280$. \emph{Scelta logica:} non si sviluppa tutta la
potenza, si individua solo l'indice $k$ del termine richiesto. Attenzione ai segni: il
fattore $(-1)^{k}$ va sempre portato con sé.
\end{soluzione}
\end{esercizio}

\begin{esercizio}
Dimostrare che $\displaystyle\sum_{k=1}^{n}k\binom nk=n\,2^{n-1}$.
\begin{soluzione}
Per $n=1$ entrambi i membri valgono $1$; sia $n\ge2$.
\emph{Calcolando:} per $1\le k\le n$ si ha
$k\binom nk=\frac{n!}{(k-1)!\,(n-k)!}=n\binom{n-1}{k-1}$, quindi, traslando l'indice,
$\sum_{k=1}^{n}k\binom nk=n\sum_{j=0}^{n-1}\binom{n-1}{j}=n\,2^{n-1}$ per il
\cref{cor:binomio}(i). \emph{Contando:} entrambi i membri contano i modi di scegliere, in
un gruppo di $n$ persone, una commissione con un presidente. A sinistra si sceglie prima
la commissione ($k$ membri, $\binom nk$ modi) e poi il presidente tra i suoi membri ($k$
modi); a destra si sceglie prima il presidente ($n$ modi) e poi, per ciascuna delle
altre $n-1$ persone, se far parte o no della commissione ($2^{n-1}$ modi).
\emph{Scelta logica:} come nella \cref{prop:simmetria-binomiale}, la stessa identità
ammette una prova algebrica e una combinatoria; la seconda spiega il risultato.
\end{soluzione}
\end{esercizio}

\subsection*{Numeri complessi}

\begin{esercizio}
Calcolare $(1+i)^{10}$.
\begin{soluzione}
Sviluppare con il binomio sarebbe faticoso. In forma trigonometrica,
$1+i=\sqrt2\bigl(\cos\frac\pi4+i\sen\frac\pi4\bigr)$; per la formula di De Moivre,
$(1+i)^{10}=(\sqrt2)^{10}\bigl(\cos\frac{10\pi}{4}+i\sen\frac{10\pi}{4}\bigr)
=32\bigl(\cos\frac{\pi}{2}+i\sen\frac\pi2\bigr)=32\,i$, essendo
$\frac{10\pi}{4}=\frac\pi2+2\pi$. \emph{Scelta logica:} la forma algebrica è adatta
alle somme, quella trigonometrica (o esponenziale) ai prodotti e alle potenze; scegliere
la rappresentazione giusta trasforma un calcolo lungo in un calcolo di una riga.
\end{soluzione}
\end{esercizio}

\begin{esercizio}
Risolvere in $\C$ l'equazione $z^{3}=-8$ e rappresentare le soluzioni nel piano di Gauss.
\begin{soluzione}
$-8=8(\cos\pi+i\sen\pi)$. Per il \cref{teo:radici-complesse} le radici hanno modulo
$\sqrt[3]8=2$ e argomenti $\frac{\pi+2k\pi}{3}$, $k=0,1,2$, cioè $\frac\pi3$, $\pi$,
$\frac53\pi$:
\[
  w_{0}=2\Bigl(\frac12+i\frac{\sqrt3}{2}\Bigr)=1+i\sqrt3, \qquad w_{1}=-2, \qquad
  w_{2}=1-i\sqrt3 .
\]
Sono i vertici di un triangolo equilatero inscritto nella circonferenza di raggio $2$.
Nel campo reale avremmo trovato soltanto $w_{1}=-2$. \emph{Scelta logica:} si scrive il
secondo membro in forma trigonometrica \emph{prima} di applicare la formula; un errore
comune è dimenticare che l'argomento di un numero reale negativo è $\pi$, non $0$.
\end{soluzione}
\end{esercizio}

\begin{esercizio}
Dimostrare che, per $z\neq0$, si ha $\abs z=1$ se e solo se $\dfrac1z=\bar z$. Qual è
il significato geometrico?
\begin{soluzione}
Da $z\bar z=\abs z^{2}$ (\cref{prop:coniugato-modulo}) segue $\frac1z=\frac{\bar
z}{\abs z^{2}}$; quindi $\frac1z=\bar z$ se e solo se $\abs z^{2}=1$, cioè $\abs z=1$.
Geometricamente: sulla circonferenza unitaria l'inversione $z\mapsto1/z$ coincide con la
simmetria rispetto all'asse reale, perché l'inverso ha argomento opposto
(\eqref{eq:quoziente-trigonometrico}) e, per i numeri di modulo $1$, lo stesso modulo.
\emph{Scelta logica:} l'identità $z\bar z=\abs z^{2}$ è il ponte tra inverso, coniugato e
modulo; usarla evita di passare alle coordinate $x$, $y$ e riduce la dimostrazione a una
riga.
\end{soluzione}
\end{esercizio}

\begin{esercizio}[Con Python]
Verificare il calcolo di Bombelli: calcolare con Python le tre radici cubiche
complesse di $2+11i$ (con la funzione \texttt{radici} del \cref{sec:complessi}) e
controllare che una di esse sia $2+i$. Quali sono le altre due? Perché sommando la
radice $2+i$ di $2+11i$ con la radice $2-i$ di $2-11i$ si ottiene la soluzione reale $x=4$
dell'equazione $x^{3}=15x+4$, mentre altre scelte delle radici non funzionano?
\begin{soluzione}
\begin{python}
import cmath, math

def radici(z, n):
    rho, theta = cmath.polar(z)
    r = rho ** (1 / n)
    return [cmath.rect(r, (theta + 2 * k * math.pi) / n)
            for k in range(n)]

for w in radici(2 + 11j, 3):
    print(f"{w:.6f}")
\end{python}
\begin{risultato}
2.000000+1.000000j
-1.866025+1.232051j
-0.133975-2.232051j
\end{risultato}
Si ottengono $2+i$ e altri due numeri, circa $-1{,}866+1{,}232\,i$ e
$-0{,}134-2{,}232\,i$: sono $(2+i)\omega$ e $(2+i)\omega^{2}$, dove
$\omega=-\frac12+i\frac{\sqrt3}{2}$ è una radice cubica dell'unità (le radici cubiche di
un numero si ottengono da una di esse moltiplicando per le radici cubiche di $1$). Nella
formula di Cardano le due radici cubiche non sono indipendenti: il loro prodotto deve
valere $p/3=5$. Infatti $(2+i)(2-i)=5$; le altre coppie ammissibili,
$(2+i)\omega+(2-i)\bar\omega$ e $(2+i)\omega^{2}+(2-i)\bar\omega^{2}$, danno le altre due
soluzioni reali, $-2\pm\sqrt3$. \emph{Scelta logica:} il calcolatore conferma i numeri,
ma è la teoria (struttura delle radici, \cref{teo:radici-complesse}) a spiegare
perché le soluzioni sono tre e come si accoppiano.
\end{soluzione}
\end{esercizio}

%% ----------------------------------------------------------------------------
%%  §2.8  Punti fondamentali del capitolo
%% ----------------------------------------------------------------------------
\section{Punti fondamentali del capitolo}
\label{sec:sintesi-cap02}

L'idea centrale del capitolo è che l'assioma di completezza diventa uno strumento di
lavoro nella forma dell'\emph{estremo superiore}: con esso si dimostrano la proprietà di
Archimede, l'esistenza delle radici $n$-esime e la biiettività dell'esponenziale. La
seconda parte insegna a contare (calcolo combinatorio, binomio di Newton) e amplia $\R$
al campo $\C$, nel quale ogni equazione di secondo grado ha soluzione.

\begin{puntichiave}
\item \textbf{Estremo superiore.} Se $A\subseteq\R$ è non vuoto e limitato superiormente,
l'insieme dei suoi maggioranti ha minimo, che si chiama $\sup A$
(\cref{teo:esistenza-sup}); questa proprietà equivale all'assioma di completezza
(\cref{oss:completezza-sup}) ed è falsa in $\Q$.
\item \textbf{Caratterizzazione.} $M=\sup A$ se e solo se $M\ge a$ per ogni $a\in A$ e,
per ogni $\epsilon>0$, esiste $a\in A$ con $a>M-\epsilon$ (\cref{prop:caratt-sup});
analogamente per $\inf A$, con le disuguaglianze rovesciate.
\item \textbf{Sup e massimo.} Il massimo esiste se e solo se $\sup A\in A$
(\cref{oss:sup-max}): errore tipico è confondere i due concetti. Se $A$ non è limitato
superiormente si scrive $\sup A=+\infty$, ma $\pm\infty$ non sono numeri reali.
\item \textbf{Proprietà di $\N$.} Ogni sottoinsieme non vuoto di $\N$ ha minimo
(\cref{teo:buon-ordinamento}); se è anche limitato superiormente, ha massimo
(\cref{teo:max-N}).
\item \textbf{Proprietà di Archimede.} Per ogni $x\in\R$ esiste $n\in\N$ con $n>x$
(\cref{teo:archimede}); ne seguono $\inf\{1/n : n\in\N\}=0$ e, per $a>0$, l'esistenza di
$n$ con $na>b$ (\cref{cor:archimede}), e la parte intera (\cref{prop:parte-intera}).
\item \textbf{Densità.} Tra due reali distinti cadono sempre un razionale e un
irrazionale (\cref{teo:densita-Q,cor:densita-irrazionali}); ma essere denso non
significa essere completo.
\item \textbf{Radici $n$-esime.} Per ogni $y\ge0$ e $n\in\N$ esiste un unico $x\ge0$ con
$x^{n}=y$, ottenuto come $\sup\{t\ge0 : t^{n}<y\}$ (\cref{teo:radice-n}).
\item \textbf{Esponenziale reale.} Per $a>1$ si pone
$a^{x}=\sup\{a^{y} : y\in\Q,\ y<x\}$ (\cref{def:potenza-reale}): la funzione $a^{x}$ è
una biiezione strettamente crescente di $\R$ su $(0,+\infty)$, decrescente se $0<a<1$
(\cref{teo:esponenziale}); il logaritmo è quindi definito su tutto $(0,+\infty)$.
\item \textbf{Contare senza elencare.} Le disposizioni di $k$ oggetti tra $n$ sono
$n!/(n-k)!$, le permutazioni $n!$, le combinazioni $\binom nk=\frac{n!}{k!\,(n-k)!}$
(\cref{teo:disposizioni,teo:combinazioni}); prima di applicare una formula occorre
chiedersi se l'ordine conta.
\item \textbf{Binomio di Newton.} $(a+b)^{n}=\sum_{k=0}^{n}\binom nk a^{n-k}b^{k}$
(\cref{teo:binomio}); i coefficienti soddisfano la formula di Stifel
$\binom n{k-1}+\binom nk=\binom{n+1}k$ (\cref{lem:stifel}), e per $x\ge0$ vale
$(1+x)^{n}\ge1+nx+\frac{n(n-1)}2x^{2}$ (\cref{cor:binomio}).
\item \textbf{Applicazioni aeronautiche.} La velocità limite di un paracadutista è un
estremo superiore che non è un massimo (\cref{es:velocita-limite}); la stima
$(1+\epsilon)^{n}\approx1+n\epsilon$ quantifica la sensibilità della portanza ($\propto
V^{2}$) e della potenza ($\propto V^{3}$) alla velocità (\cref{es:sensibilita}).
\item \textbf{Il campo complesso.} $\C=\R^{2}$, con il prodotto
\eqref{eq:prodotto-complessi}, è un campo (\cref{teo:C-campo}) in cui $i^{2}=-1$, ma
non è ordinabile (\cref{prop:C-non-ordinabile}); vale $z\bar z=\abs z^{2}$ e la
disuguaglianza triangolare (\cref{prop:coniugato-modulo}).
\item \textbf{Forma trigonometrica.} Scritto $z=\rho(\cos\vartheta+i\sen\vartheta)$,
o in breve $z=\rho e^{i\vartheta}$, nel prodotto i moduli si moltiplicano e gli
argomenti si sommano (\cref{teo:prodotto-trigonometrico}), da cui la formula di De
Moivre (\cref{cor:de-moivre}) e la rotazione dagli assi vento agli assi corda
(\cref{es:assi-profilo}).
\item \textbf{Radici complesse.} Ogni $z\neq0$ ha esattamente $n$ radici $n$-esime, che
per $n\ge3$ sono i vertici di un poligono regolare (\cref{teo:radici-complesse}); in
$\C$ ogni equazione di
secondo grado ha soluzione (\cref{prop:secondo-grado-C}).
\end{puntichiave}

\paragraph{Il filo del discorso.} L'estremo superiore è il motore dei teoremi di
esistenza, ma finora l'abbiamo applicato a insiemi «statici». Nel \cref{cap:successioni}
lo useremo per studiare i limiti di successioni: una successione monotòna e limitata
converge proprio al suo estremo superiore o inferiore (\cref{teo:monotone}), e da qui
nasceranno il numero $e$ e il teorema di Bolzano--Weierstrass. La formula di Eulero,
qui introdotta come notazione, sarà dimostrata nel capitolo sulle serie
(\cref{teo:formula-eulero}).

